% !TeX program = xelatex
% Standalone English volume. Use XeLaTeX until the references stabilize.
\documentclass[a4paper,11pt]{article}
\usepackage[T1]{fontenc}
\usepackage{lmodern}
\usepackage[english]{babel}
\usepackage[left=26mm,right=26mm,top=25mm,bottom=25mm,headheight=15pt,headsep=7mm,footskip=11mm]{geometry}
\usepackage{amsmath,amssymb,mathtools,amsthm}
\newtheorem{proposition}{Proposition}[section]
\newtheorem{lemma}[proposition]{Lemma}
\newtheorem{theorem}[proposition]{Theorem}
\newtheorem{corollary}[proposition]{Corollary}
\newtheorem{maintheorem}[proposition]{Main Theorem}
\theoremstyle{definition}
\newtheorem{definition}[proposition]{Definition}
\theoremstyle{remark}
\newtheorem{remark}[proposition]{Remark}
\makeatletter
\renewcommand*\l@section{\@dottedtocline{1}{0em}{2.4em}}
\renewcommand*\l@subsection{\@dottedtocline{2}{1.45em}{3.45em}}
\renewcommand*\l@subsubsection{\@dottedtocline{3}{3.8em}{4.55em}}
\makeatother
\usepackage{xcolor,fancyhdr,titlesec,needspace,xurl}
\usepackage{longtable,multicol}
\usepackage{hyperref}
\definecolor{ink}{HTML}{1A3542}
\definecolor{muted}{HTML}{4F606A}
\hypersetup{unicode=true,colorlinks=true,linkcolor=ink,filecolor=ink,urlcolor=ink,citecolor=ink,bookmarksnumbered=true,
 pdfdisplaydoctitle=true,pdflang={en},
 pdftitle={The generalized Riemann hypothesis for Dirichlet L-functions. Edition 930, Volume XVI},
 pdfsubject={A natural-language reconstruction of the proposed proof},
 pdfauthor={A mathematical reconstruction}}
\titleformat{\section}{\Large\bfseries\color{ink}}{\thesection}{.65em}{}
\titlespacing*{\section}{0pt}{.8em}{.65em}
\setlength{\parindent}{1.2em}
\linespread{1.045}
\setlength{\parskip}{.22em plus .05em minus .03em}
\setlength{\emergencystretch}{2em}
\setcounter{topnumber}{2}
\setcounter{tocdepth}{3}\setcounter{secnumdepth}{3}
\widowpenalty=4500
\clubpenalty=4500
\pagestyle{fancy}\fancyhf{}
\fancyhead[L]{\small\color{muted}Dirichlet GRH \textperiodcentered\ Volume I}
\fancyhead[R]{\small\color{muted}Dirichlet zeros and power intervals}
\fancyfoot[C]{\small\thepage}
\renewcommand{\headrulewidth}{.3pt}
\newcommand{\lead}[1]{\par\smallskip\noindent{\bfseries\color{ink}#1.}\enspace}
\newcommand{\sour}[2]{\unskip\penalty150\hspace{.35em}{\small\hyperlink{ref#1}{[#2]}}}
\titleformat{\subsection}{\large\bfseries}{\thesubsection}{.6em}{}
\titlespacing*{\subsection}{0pt}{1.15em plus .15em}{.4em}
\setlength{\abovedisplayskip}{9pt plus 2pt minus 3pt}
\setlength{\belowdisplayskip}{9pt plus 2pt minus 3pt}
\newcommand{\RR}{\mathbb R}\newcommand{\CC}{\mathbb C}\newcommand{\QQ}{\mathbb Q}
\newcommand{\ZZ}{\mathbb Z}\newcommand{\ord}{\operatorname{ord}}
\newcommand{\Frob}{\operatorname{Frob}}\newcommand{\Gal}{\operatorname{Gal}}
\newcommand{\len}{\operatorname{length}}\newcommand{\Res}{\operatorname*{Res}}
\newcommand{\wmean}{\operatorname{av}}\newcommand{\Ker}{\operatorname{ker}}
\newcommand{\Lh}{\mathcal L}\newcommand{\Gg}{\mathcal G}
\newcommand{\real}{\operatorname{Re}}\newcommand{\imag}{\operatorname{Im}}
\newcommand{\abs}[1]{\left|#1\right|}\newcommand{\eps}{\varepsilon}
\newcommand{\anchoritem}[2]{\par\Needspace{2.8\baselineskip}\hypertarget{ref#1}{}\noindent{\bfseries\color{ink}#2.}\enspace}
\newcommand{\srcloc}[2]{{\small\nolinkurl{#1.lean}}, line~#2}
\numberwithin{equation}{section}
\raggedbottom

\newcommand{\Sel}{\mathcal S}
\newcommand{\Cl}{\operatorname{Cl}}
\newcommand{\Aut}{\operatorname{Aut}}
\newcommand{\GL}{\operatorname{GL}}
\newcommand{\OO}{\mathcal O}
\hypersetup{bookmarksdepth=3}


\addto\captionsenglish{\renewcommand{\partname}{Volume}}


\makeatletter
\expandafter\def\csname r@cor:actualdenom803\endcsname{{52.21}{80}{}{proposition.52.21}{Dirichlet-GRH-v930-Volume-X.pdf}}
\expandafter\def\csname r@cor:actualmarks714\endcsname{{12.10}{38}{}{proposition.12.10}{Dirichlet-GRH-v930-Volume-II.pdf}}
\expandafter\def\csname r@cor:complementreproduce812\endcsname{{77.4}{6}{}{proposition.77.4}{Dirichlet-GRH-v930-Volume-XV.pdf}}
\expandafter\def\csname r@cor:detheight810\endcsname{{73.23}{95}{}{proposition.73.23}{Dirichlet-GRH-v930-Volume-XIV.pdf}}
\expandafter\def\csname r@cor:gluedperiod811\endcsname{{73.25}{97}{}{proposition.73.25}{Dirichlet-GRH-v930-Volume-XIV.pdf}}
\expandafter\def\csname r@cor:normaldegree347\endcsname{{127.2}{27}{}{proposition.127.2}{Dirichlet-GRH-v930-Volume-XXIV.pdf}}
\expandafter\def\csname r@cor:rank111\endcsname{{9.14}{119}{}{proposition.9.14}{Dirichlet-GRH-v930-Volume-I.pdf}}
\expandafter\def\csname r@cor:trace165\endcsname{{93.6}{25}{}{proposition.93.6}{Dirichlet-GRH-v930-Volume-XVIII.pdf}}
\expandafter\def\csname r@eq:common177\endcsname{{72.8}{64}{}{equation.72.8}{Dirichlet-GRH-v930-Volume-XIV.pdf}}
\expandafter\def\csname r@eq:fixedmoment332\endcsname{{119.1}{6}{}{equation.119.1}{Dirichlet-GRH-v930-Volume-XXIII.pdf}}
\expandafter\def\csname r@eq:geom163\endcsname{{86.13}{18}{}{equation.86.13}{Dirichlet-GRH-v930-Volume-XVII.pdf}}
\expandafter\def\csname r@eq:letters111\endcsname{{9.1}{117}{}{equation.9.1}{Dirichlet-GRH-v930-Volume-I.pdf}}
\expandafter\def\csname r@eq:period94\endcsname{{121.7}{51}{}{equation.121.7}{Dirichlet-GRH-v930-Volume-XXIII.pdf}}
\expandafter\def\csname r@eq:region198\endcsname{{20.12}{33}{}{equation.20.12}{Dirichlet-GRH-v930-Volume-IV.pdf}}
\expandafter\def\csname r@eq:relchain810\endcsname{{73.1}{94}{}{equation.73.1}{Dirichlet-GRH-v930-Volume-XIV.pdf}}
\expandafter\def\csname r@eq:sieve\endcsname{{70.2}{4}{}{equation.70.2}{Dirichlet-GRH-v930-Volume-XIV.pdf}}
\expandafter\def\csname r@eq:tail114\endcsname{{19.15}{17}{}{equation.19.15}{Dirichlet-GRH-v930-Volume-IV.pdf}}
\expandafter\def\csname r@eq:weil161\endcsname{{93.1}{21}{}{equation.93.1}{Dirichlet-GRH-v930-Volume-XVIII.pdf}}
\expandafter\def\csname r@eq:window\endcsname{{{W}}{2}{}{AMS.2}{Dirichlet-GRH-v930-Volume-I.pdf}}
\expandafter\def\csname r@lem:auxiliaryroots339\endcsname{{99.16}{57}{}{proposition.99.16}{Dirichlet-GRH-v930-Volume-XIX.pdf}}
\expandafter\def\csname r@lem:auxmass805\endcsname{{73.4}{77}{}{proposition.73.4}{Dirichlet-GRH-v930-Volume-XIV.pdf}}
\expandafter\def\csname r@lem:basis59\endcsname{{65.8}{27}{}{proposition.65.8}{Dirichlet-GRH-v930-Volume-XIII.pdf}}
\expandafter\def\csname r@lem:centraltest736\endcsname{{50.2}{4}{}{proposition.50.2}{Dirichlet-GRH-v930-Volume-X.pdf}}
\expandafter\def\csname r@lem:complementfischer819\endcsname{{77.27}{30}{}{proposition.77.27}{Dirichlet-GRH-v930-Volume-XV.pdf}}
\expandafter\def\csname r@lem:dimension239\endcsname{{94.22}{55}{}{proposition.94.22}{Dirichlet-GRH-v930-Volume-XVIII.pdf}}
\expandafter\def\csname r@lem:dualcartan811\endcsname{{73.26}{97}{}{proposition.73.26}{Dirichlet-GRH-v930-Volume-XIV.pdf}}
\expandafter\def\csname r@lem:kappa62\endcsname{{66.12}{54}{}{proposition.66.12}{Dirichlet-GRH-v930-Volume-XIII.pdf}}
\expandafter\def\csname r@lem:mass32\endcsname{{121.17}{48}{}{proposition.121.17}{Dirichlet-GRH-v930-Volume-XXIII.pdf}}
\expandafter\def\csname r@lem:oscillatorcartan710\endcsname{{142.12}{45}{}{proposition.142.12}{Dirichlet-GRH-v930-Volume-XXVII.pdf}}
\expandafter\def\csname r@lem:ramifiedsum334\endcsname{{90.9}{92}{}{proposition.90.9}{Dirichlet-GRH-v930-Volume-XVII.pdf}}
\expandafter\def\csname r@lem:secondhecke829\endcsname{{78.7}{61}{}{proposition.78.7}{Dirichlet-GRH-v930-Volume-XV.pdf}}
\expandafter\def\csname r@lem:signkernel823\endcsname{{77.37}{43}{}{proposition.77.37}{Dirichlet-GRH-v930-Volume-XV.pdf}}
\expandafter\def\csname r@lem:size97\endcsname{{87.1}{33}{}{proposition.87.1}{Dirichlet-GRH-v930-Volume-XVII.pdf}}
\expandafter\def\csname r@lem:spacing781\endcsname{{51.19}{54}{}{proposition.51.19}{Dirichlet-GRH-v930-Volume-X.pdf}}
\expandafter\def\csname r@lem:sparsedistance814\endcsname{{77.9}{11}{}{proposition.77.9}{Dirichlet-GRH-v930-Volume-XV.pdf}}
\expandafter\def\csname r@lem:split61\endcsname{{65.12}{32}{}{proposition.65.12}{Dirichlet-GRH-v930-Volume-XIII.pdf}}
\expandafter\def\csname r@lem:tracebasis821\endcsname{{77.33}{37}{}{proposition.77.33}{Dirichlet-GRH-v930-Volume-XV.pdf}}
\expandafter\def\csname r@lem:twoelliptic806\endcsname{{73.7}{80}{}{proposition.73.7}{Dirichlet-GRH-v930-Volume-XIV.pdf}}
\expandafter\def\csname r@lem:unit41\endcsname{{64.5}{7}{}{proposition.64.5}{Dirichlet-GRH-v930-Volume-XIII.pdf}}
\expandafter\def\csname r@lem:zak242\endcsname{{67.7}{80}{}{proposition.67.7}{Dirichlet-GRH-v930-Volume-XIII.pdf}}
\expandafter\def\csname r@nav:VII-00\endcsname{{86.1}{3}{}{subsection.86.1}{Dirichlet-GRH-v930-Volume-XVII.pdf}}
\expandafter\def\csname r@prop:archcontraction804\endcsname{{52.22}{81}{}{proposition.52.22}{Dirichlet-GRH-v930-Volume-X.pdf}}
\expandafter\def\csname r@prop:branch275\endcsname{{94.29}{62}{}{proposition.94.29}{Dirichlet-GRH-v930-Volume-XVIII.pdf}}
\expandafter\def\csname r@prop:centeredflag341\endcsname{{99.18}{58}{}{proposition.99.18}{Dirichlet-GRH-v930-Volume-XIX.pdf}}
\expandafter\def\csname r@prop:chainpairing810\endcsname{{73.21}{93}{}{proposition.73.21}{Dirichlet-GRH-v930-Volume-XIV.pdf}}
\expandafter\def\csname r@prop:column110\endcsname{{86.29}{23}{}{proposition.86.29}{Dirichlet-GRH-v930-Volume-XVII.pdf}}
\expandafter\def\csname r@prop:doubledkernel711\endcsname{{142.14}{47}{}{proposition.142.14}{Dirichlet-GRH-v930-Volume-XXVII.pdf}}
\expandafter\def\csname r@prop:field12\endcsname{{6.4}{27}{}{proposition.6.4}{Dirichlet-GRH-v930-Volume-I.pdf}}
\expandafter\def\csname r@prop:markedfourier820\endcsname{{77.30}{33}{}{proposition.77.30}{Dirichlet-GRH-v930-Volume-XV.pdf}}
\expandafter\def\csname r@prop:sphericalframe334\endcsname{{90.11}{93}{}{proposition.90.11}{Dirichlet-GRH-v930-Volume-XVII.pdf}}
\expandafter\def\csname r@prop:split175\endcsname{{72.17}{60}{}{proposition.72.17}{Dirichlet-GRH-v930-Volume-XIV.pdf}}
\expandafter\def\csname r@prop:terminalcore255\endcsname{{95.4}{68}{}{proposition.95.4}{Dirichlet-GRH-v930-Volume-XVIII.pdf}}
\expandafter\def\csname r@prop:terminalmetric807\endcsname{{73.10}{83}{}{proposition.73.10}{Dirichlet-GRH-v930-Volume-XIV.pdf}}
\expandafter\def\csname r@prop:testedkernel785\endcsname{{52.4}{62}{}{proposition.52.4}{Dirichlet-GRH-v930-Volume-X.pdf}}
\expandafter\def\csname r@prop:torusorder736\endcsname{{50.1}{3}{}{proposition.50.1}{Dirichlet-GRH-v930-Volume-X.pdf}}
\expandafter\def\csname r@sec:arch59\endcsname{{65.3}{27}{}{subsection.65.3}{Dirichlet-GRH-v930-Volume-XIII.pdf}}
\expandafter\def\csname r@sec:auxdet805\endcsname{{73.4}{77}{}{subsection.73.4}{Dirichlet-GRH-v930-Volume-XIV.pdf}}
\expandafter\def\csname r@sec:commonsource806\endcsname{{73.5}{80}{}{subsection.73.5}{Dirichlet-GRH-v930-Volume-XIV.pdf}}
\expandafter\def\csname r@sec:complementperiod812\endcsname{{77.1}{3}{}{subsection.77.1}{Dirichlet-GRH-v930-Volume-XV.pdf}}
\expandafter\def\csname r@sec:exchange336\endcsname{{86.11}{26}{}{subsection.86.11}{Dirichlet-GRH-v930-Volume-XVII.pdf}}
\expandafter\def\csname r@sec:fixedcolumn332\endcsname{{119.3}{5}{}{subsection.119.3}{Dirichlet-GRH-v930-Volume-XXIII.pdf}}
\expandafter\def\csname r@sec:integralgluing811\endcsname{{73.10}{96}{}{subsection.73.10}{Dirichlet-GRH-v930-Volume-XIV.pdf}}
\expandafter\def\csname r@sec:matchingframes334\endcsname{{90.9}{92}{}{subsection.90.9}{Dirichlet-GRH-v930-Volume-XVII.pdf}}
\expandafter\def\csname r@sec:middle131\endcsname{{86.1}{3}{}{subsection.86.1}{Dirichlet-GRH-v930-Volume-XVII.pdf}}
\expandafter\def\csname r@sec:missing\endcsname{{6.6}{32}{}{subsection.6.6}{Dirichlet-GRH-v930-Volume-I.pdf}}
\expandafter\def\csname r@sec:rankonepair736\endcsname{{50.1}{3}{}{subsection.50.1}{Dirichlet-GRH-v930-Volume-X.pdf}}
\expandafter\def\csname r@sec:refinedclass824\endcsname{{77.13}{46}{}{subsection.77.13}{Dirichlet-GRH-v930-Volume-XV.pdf}}
\expandafter\def\csname r@sec:secondglobal833\endcsname{{78.7}{72}{}{subsection.78.7}{Dirichlet-GRH-v930-Volume-XV.pdf}}
\expandafter\def\csname r@sec:splitcost785\endcsname{{52.2}{62}{}{subsection.52.2}{Dirichlet-GRH-v930-Volume-X.pdf}}
\expandafter\def\csname r@sec:stackedcoordinates819\endcsname{{77.8}{30}{}{subsection.77.8}{Dirichlet-GRH-v930-Volume-XV.pdf}}
\expandafter\def\csname r@seriesguide\endcsname{{1.1}{4}{}{subsection.1.1}{Dirichlet-GRH-v930-Volume-I.pdf}}
\expandafter\def\csname r@thm:allrows251\endcsname{{71.11}{27}{}{proposition.71.11}{Dirichlet-GRH-v930-Volume-XIV.pdf}}
\expandafter\def\csname r@thm:arithmeticclass824\endcsname{{77.41}{48}{}{proposition.77.41}{Dirichlet-GRH-v930-Volume-XV.pdf}}
\expandafter\def\csname r@thm:auxdet805\endcsname{{73.5}{78}{}{proposition.73.5}{Dirichlet-GRH-v930-Volume-XIV.pdf}}
\expandafter\def\csname r@thm:auxiliaryflag339\endcsname{{99.17}{57}{}{proposition.99.17}{Dirichlet-GRH-v930-Volume-XIX.pdf}}
\expandafter\def\csname r@thm:auxoverlap341\endcsname{{99.19}{59}{}{proposition.99.19}{Dirichlet-GRH-v930-Volume-XIX.pdf}}
\expandafter\def\csname r@thm:class61\endcsname{{65.11}{31}{}{proposition.65.11}{Dirichlet-GRH-v930-Volume-XIII.pdf}}
\expandafter\def\csname r@thm:column242\endcsname{{121.3}{31}{}{proposition.121.3}{Dirichlet-GRH-v930-Volume-XXIII.pdf}}
\expandafter\def\csname r@thm:common111\endcsname{{9.13}{118}{}{proposition.9.13}{Dirichlet-GRH-v930-Volume-I.pdf}}
\expandafter\def\csname r@thm:commonfischer348\endcsname{{120.2}{15}{}{proposition.120.2}{Dirichlet-GRH-v930-Volume-XXIII.pdf}}
\expandafter\def\csname r@thm:commonsource806\endcsname{{73.8}{81}{}{proposition.73.8}{Dirichlet-GRH-v930-Volume-XIV.pdf}}
\expandafter\def\csname r@thm:complexseparation818\endcsname{{77.25}{27}{}{proposition.77.25}{Dirichlet-GRH-v930-Volume-XV.pdf}}
\expandafter\def\csname r@thm:coregram258\endcsname{{75.3}{115}{}{proposition.75.3}{Dirichlet-GRH-v930-Volume-XIV.pdf}}
\expandafter\def\csname r@thm:criterion16\endcsname{{1.2}{2}{}{proposition.1.2}{Dirichlet-GRH-v930-Volume-I.pdf}}
\expandafter\def\csname r@thm:criticalheights734\endcsname{{13.13}{72}{}{proposition.13.13}{Dirichlet-GRH-v930-Volume-II.pdf}}
\expandafter\def\csname r@thm:definiteclass822\endcsname{{77.36}{41}{}{proposition.77.36}{Dirichlet-GRH-v930-Volume-XV.pdf}}
\expandafter\def\csname r@thm:depth164\endcsname{{110.3}{39}{}{proposition.110.3}{Dirichlet-GRH-v930-Volume-XXI.pdf}}
\expandafter\def\csname r@thm:detquotient810\endcsname{{73.22}{94}{}{proposition.73.22}{Dirichlet-GRH-v930-Volume-XIV.pdf}}
\expandafter\def\csname r@thm:downcore254\endcsname{{71.13}{28}{}{proposition.71.13}{Dirichlet-GRH-v930-Volume-XIV.pdf}}
\expandafter\def\csname r@thm:elliptic312\endcsname{{67.12}{86}{}{proposition.67.12}{Dirichlet-GRH-v930-Volume-XIII.pdf}}
\expandafter\def\csname r@thm:extract275\endcsname{{88.8}{62}{}{proposition.88.8}{Dirichlet-GRH-v930-Volume-XVII.pdf}}
\expandafter\def\csname r@thm:finitecolumn313\endcsname{{122.10}{67}{}{proposition.122.10}{Dirichlet-GRH-v930-Volume-XXIII.pdf}}
\expandafter\def\csname r@thm:first250\endcsname{{71.10}{25}{}{proposition.71.10}{Dirichlet-GRH-v930-Volume-XIV.pdf}}
\expandafter\def\csname r@thm:firsthol817\endcsname{{77.22}{24}{}{proposition.77.22}{Dirichlet-GRH-v930-Volume-XV.pdf}}
\expandafter\def\csname r@thm:fixedconductor781\endcsname{{51.20}{55}{}{proposition.51.20}{Dirichlet-GRH-v930-Volume-X.pdf}}
\expandafter\def\csname r@thm:germ57\endcsname{{65.5}{24}{}{proposition.65.5}{Dirichlet-GRH-v930-Volume-XIII.pdf}}
\expandafter\def\csname r@thm:gluedcost811\endcsname{{73.27}{98}{}{proposition.73.27}{Dirichlet-GRH-v930-Volume-XIV.pdf}}
\expandafter\def\csname r@thm:graphtransport337\endcsname{{71.21}{38}{}{proposition.71.21}{Dirichlet-GRH-v930-Volume-XIV.pdf}}
\expandafter\def\csname r@thm:hurwitz62\endcsname{{66.13}{55}{}{proposition.66.13}{Dirichlet-GRH-v930-Volume-XIII.pdf}}
\expandafter\def\csname r@thm:index165\endcsname{{93.4}{24}{}{proposition.93.4}{Dirichlet-GRH-v930-Volume-XVIII.pdf}}
\expandafter\def\csname r@thm:markeddefinite820\endcsname{{77.32}{34}{}{proposition.77.32}{Dirichlet-GRH-v930-Volume-XV.pdf}}
\expandafter\def\csname r@thm:markedstrict337\endcsname{{22.11}{70}{}{proposition.22.11}{Dirichlet-GRH-v930-Volume-IV.pdf}}
\expandafter\def\csname r@thm:markedtype801\endcsname{{17.33}{119}{}{proposition.17.33}{Dirichlet-GRH-v930-Volume-III.pdf}}
\expandafter\def\csname r@thm:matchingframe334\endcsname{{90.10}{93}{}{proposition.90.10}{Dirichlet-GRH-v930-Volume-XVII.pdf}}
\expandafter\def\csname r@thm:mean56\endcsname{{65.22}{40}{}{proposition.65.22}{Dirichlet-GRH-v930-Volume-XIII.pdf}}
\expandafter\def\csname r@thm:midlevel131\endcsname{{86.1}{3}{}{proposition.86.1}{Dirichlet-GRH-v930-Volume-XVII.pdf}}
\expandafter\def\csname r@thm:midnorm131\endcsname{{86.2}{4}{}{proposition.86.2}{Dirichlet-GRH-v930-Volume-XVII.pdf}}
\expandafter\def\csname r@thm:negativeinterval740\endcsname{{15.3}{4}{}{proposition.15.3}{Dirichlet-GRH-v930-Volume-III.pdf}}
\expandafter\def\csname r@thm:norm177\endcsname{{72.21}{64}{}{proposition.72.21}{Dirichlet-GRH-v930-Volume-XIV.pdf}}
\expandafter\def\csname r@thm:one61\endcsname{{65.13}{33}{}{proposition.65.13}{Dirichlet-GRH-v930-Volume-XIII.pdf}}
\expandafter\def\csname r@thm:ordinary164\endcsname{{110.2}{39}{}{proposition.110.2}{Dirichlet-GRH-v930-Volume-XXI.pdf}}
\expandafter\def\csname r@thm:pol49\endcsname{{64.10}{12}{}{proposition.64.10}{Dirichlet-GRH-v930-Volume-XIII.pdf}}
\expandafter\def\csname r@thm:projectionindex784\endcsname{{52.2}{59}{}{proposition.52.2}{Dirichlet-GRH-v930-Volume-X.pdf}}
\expandafter\def\csname r@thm:range239\endcsname{{88.1}{52}{}{proposition.88.1}{Dirichlet-GRH-v930-Volume-XVII.pdf}}
\expandafter\def\csname r@thm:rangenew50\endcsname{{6.29}{62}{}{proposition.6.29}{Dirichlet-GRH-v930-Volume-I.pdf}}
\expandafter\def\csname r@thm:rankonepair736\endcsname{{50.3}{4}{}{proposition.50.3}{Dirichlet-GRH-v930-Volume-X.pdf}}
\expandafter\def\csname r@thm:rankoneselmer737\endcsname{{50.6}{6}{}{proposition.50.6}{Dirichlet-GRH-v930-Volume-X.pdf}}
\expandafter\def\csname r@thm:realcore253\endcsname{{71.16}{32}{}{proposition.71.16}{Dirichlet-GRH-v930-Volume-XIV.pdf}}
\expandafter\def\csname r@thm:realgram710\endcsname{{142.13}{46}{}{proposition.142.13}{Dirichlet-GRH-v930-Volume-XXVII.pdf}}
\expandafter\def\csname r@thm:scalarorbits332\endcsname{{119.4}{5}{}{proposition.119.4}{Dirichlet-GRH-v930-Volume-XXIII.pdf}}
\expandafter\def\csname r@thm:signedfibers331\endcsname{{99.12}{53}{}{proposition.99.12}{Dirichlet-GRH-v930-Volume-XIX.pdf}}
\expandafter\def\csname r@thm:signedseries823\endcsname{{77.38}{44}{}{proposition.77.38}{Dirichlet-GRH-v930-Volume-XV.pdf}}
\expandafter\def\csname r@thm:small274\endcsname{{88.7}{61}{}{proposition.88.7}{Dirichlet-GRH-v930-Volume-XVII.pdf}}
\expandafter\def\csname r@thm:splitlattice811\endcsname{{73.24}{96}{}{proposition.73.24}{Dirichlet-GRH-v930-Volume-XIV.pdf}}
\expandafter\def\csname r@thm:splitzeta827\endcsname{{78.1}{54}{}{proposition.78.1}{Dirichlet-GRH-v930-Volume-XV.pdf}}
\expandafter\def\csname r@thm:stackedheight819\endcsname{{77.29}{32}{}{proposition.77.29}{Dirichlet-GRH-v930-Volume-XV.pdf}}
\expandafter\def\csname r@thm:terminaldual325\endcsname{{122.6}{62}{}{proposition.122.6}{Dirichlet-GRH-v930-Volume-XXIII.pdf}}
\expandafter\def\csname r@thm:terminalformula347\endcsname{{127.1}{26}{}{proposition.127.1}{Dirichlet-GRH-v930-Volume-XXIV.pdf}}
\expandafter\def\csname r@thm:terminalwalk338\endcsname{{71.24}{40}{}{proposition.71.24}{Dirichlet-GRH-v930-Volume-XIV.pdf}}
\expandafter\def\csname r@thm:tests165\endcsname{{93.5}{25}{}{proposition.93.5}{Dirichlet-GRH-v930-Volume-XVIII.pdf}}
\expandafter\def\csname r@thm:thirdglobal843\endcsname{{79.20}{100}{}{proposition.79.20}{Dirichlet-GRH-v930-Volume-XV.pdf}}
\expandafter\def\csname r@thm:torsion312\endcsname{{66.35}{74}{}{proposition.66.35}{Dirichlet-GRH-v930-Volume-XIII.pdf}}
\expandafter\def\csname r@thm:trans57\endcsname{{65.6}{25}{}{proposition.65.6}{Dirichlet-GRH-v930-Volume-XIII.pdf}}
\expandafter\def\csname r@thm:transport59\endcsname{{65.9}{28}{}{proposition.65.9}{Dirichlet-GRH-v930-Volume-XIII.pdf}}
\expandafter\def\csname r@thm:types276\endcsname{{90.7}{89}{}{proposition.90.7}{Dirichlet-GRH-v930-Volume-XVII.pdf}}
\expandafter\def\csname r@thm:vac242\endcsname{{67.8}{81}{}{proposition.67.8}{Dirichlet-GRH-v930-Volume-XIII.pdf}}
\makeatother
\begin{document}
\setcounter{part}{15}
\fancyhead[L]{\small\color{muted}Dirichlet GRH \textperiodcentered\ Volume XVI}
\fancyhead[R]{\small\color{muted}Theta projections and determinants}
\thispagestyle{plain}\begin{center}
{\LARGE\bfseries The generalized Riemann hypothesis\\[2pt]for Dirichlet \(L\)-functions}\\[6pt]
{\small 27 November 2026 \quad\textperiodcentered\quad Edition 930}\end{center}
\part{Integral theta projections and terminal determinants}\label{part:XVI}
\begin{abstract}\noindent We descend the integral projection through the theta tower, evaluate the critical torus tests and study the terminal coefficient lattices. Finite valuations and archimedean norms give separate bounds on the same projected determinants.\end{abstract}
\setcounter{section}{80}
\noindent The \hyperref[seriesguide]{guide in Volume~I} lists all twenty-seven volumes.
\begingroup\small\setlength{\parskip}{0pt}\tableofcontents\endgroup
\clearpage

\section{Integral projection and descent in the theta tower}\label{nav:theta-descent}\label{nav:volume-XVI}

\subsection{An integral map on the definite target and its adjoint on good lifts}
\label{sec:gluedperiodmap844}
Return to the single lift from \(G=U(W_j)\), where \(s=\dim W_j=k+m-1\),
to the definite group of rank \(s+t\), where \(t=5s\). Decompose its
space as \(I_p\oplus V^\circ\) as in Section~\ref{sec:integralgluing811}.
Let \(\Theta_{j,\vartheta}^{\rm cl}\) denote the lift with its original
point test at the old marks, its determinant test at the new ones,
and the target class-sum metric of Section~\ref{sec:auxdet805}.
We specify the real normalization by testing the mixed Weyl vector
with its Fock polynomial in the flag-minor normalization of
Section~\ref{sec:stackedcoordinates819}. It is a common normalization
for all twists. Hermitian pairings in this subsection are linear in
the second argument.

\begin{lemma}[The row test used by the divided lattice]
\leavevmode\label{lem:gluedrow844}
For an index \(i\) with old matrix \(\Lambda\), stack the rows of
\(\Lambda\) over those of \(\varpi I_p\) to form \(w_i\) over
\(R=\mathcal O_E/\mathfrak a^2\) in the chosen components. Use
a simultaneous residue uniformizer \(\varpi\) at the old marks.
The first-block datum of Corollary~\ref{cor:gluedperiod811} equals
\(D_{w_i}\) of Theorem~\ref{thm:arithmeticclass824} at every finite
place. All stabilizers in the refined class sum of that theorem
are trivial for this choice of rows. The complementary class set
at level \(J_H\) of Section~\ref{sec:complementperiod812} also has
trivial stabilizers.
\end{lemma}
\begin{proof}
At an old mark the lower rows impose \(X=e+\pi U\), with \(U\)
integral. The upper rows then impose
\(\bar\Lambda\bar U=0\) over the residue field, exactly the condition
\(X\in e+(\pi K)^s\) in Theorem~\ref{thm:splitlattice811}.
On this set \(e^{\mathsf t}X\) reduces to the identity modulo \(\pi\),
so its tame determinant character is one. At a new mark the row
condition on the old residue ring is empty. The new datum on the
integral box is the character of the determinant of the \(e\)-rows
extended by zero, as required. Elsewhere the test is the lattice
vacuum. The upper rows already contain \(\Lambda\), so the
factorization condition on \(w\) in Theorem~\ref{thm:arithmeticclass824}
holds.

Each definite class has a finite rational stabilizer. The condition
\(w_i\bar A=w_i\) forces \(A\equiv I\bmod\pi\) at every old mark
by the lower rows. At a mark of rational residue characteristic
\(l\), the first principal congruence subgroup is pro-\(l\):
the successive congruence quotients embed in the additive groups
of residue matrices and their kernels intersect trivially.
Projection of a rational matrix into any one completion, followed
by conjugation by the class representative, is injective. Thus the
order of each stabilizer element is a power of both distinct old
rational primes \(p_0,p_1\). Such an element is the identity. The
complementary level is principal modulo \(\pi^2\) at both primes,
so the same argument applies there.
\end{proof}

Write \(T_j\) for the direct sum over pairs \((i,y)\) of the negative
Weyl coefficient spaces on \(V^\circ\), where \(y\) runs through
representatives of its classes at \(J_H\). Take their direct-sum
Hermitian metric. In a tuple \(\zeta\), integrality means membership
in the transported negative coefficient lattice at each \(y\), with
the contragredient action of Lemma~\ref{lem:dualcartan811}. Denote
the good source space by \(\mathcal H_j\). Use the integral
normalizing power \(M^{E_j}\) of Theorem~\ref{thm:gluedcost811}.
For each \(i\), choose its \(\gamma_i\) from
Corollary~\ref{cor:gluedperiod811} once for all twists. Let
\(\mathcal C_{i,\vartheta}\) denote the refined classes on \(I_p\)
in the lemma.

For a point \(z\) of the finite definite target and a mixed
coefficient vector \(v\), define \(\Delta_z(v)\) to be the form
supported on its level class with value
\(\sum_{\gamma\in\Gamma_z}\rho_z(\gamma)v\) at \(z\).
Here \(\Gamma_z\) is its finite rational stabilizer and \(\rho_z\)
is its action on the coefficient fiber, including the fixed level
character. Extend the value equivariantly to the class. It is
invariant under the stabilizer, so the extension is well defined.
With the reciprocal stabilizer weights one has
\[
          \langle\Delta_z(v),A\rangle
                        =\langle v,A(z)\rangle .
\]
Indeed \(A(z)\) is invariant and the action is unitary. Each term
in the sum contributes the same pairing and their number cancels
the stabilizer denominator.

Write \(u\odot v\) for the mixed Cartan coefficient obtained on
disjoint blocks as in Lemma~\ref{lem:dualcartan811}. Set
\[
 \begin{split}
 z_{i,h,y}&=\operatorname{diag}(h,y)\gamma_i,\\
 G_\vartheta\zeta
  &=M^{E_j}
    \sum_{i,y}\;\sum_{h\in\mathcal C_{i,\vartheta}}
       \Delta_{z_{i,h,y}}\bigl(
                \mathcal Y_{i,\vartheta}(h)\odot\zeta_{i,y}\bigr).
 \end{split}
\]
The first coefficient is the arithmetic vector of
Theorem~\ref{thm:signedseries823} for the indicated \(i\).
The formula is linear in the tuple. Repeated target classes are
summed in their common fiber.

\begin{theorem}[Integral tests and a common injective period map]
\leavevmode\label{thm:gluedperiodmap844}
At every coefficient place above a prime away from the new twist
conductor, \(G_\vartheta\) takes integral tuples to integral
definite forms in the original target coefficient lattice.
There is a linear injection \(P_j:\mathcal H_j\longrightarrow T_j\)
independent of \(\vartheta\) such that
\[
        b_{j,\vartheta}\,G_\vartheta^*\Theta_{j,\vartheta}^{\rm cl}
                  \big|_{\mathcal H_j}=P_j .
\]
In particular, if \(d_j=\dim\mathcal H_j\), the determinant of the
Gram operator of \(G_\vartheta^*\Theta_{j,\vartheta}^{\rm cl}\) on
this space equals
\[
             |b_{j,\vartheta}|^{-2d_j}\det(P_j^*P_j)>0 .
\]
The determinant on a zero-dimensional source is one.
\end{theorem}
\begin{proof}
The positive coefficients are integral in their transported lattices
by Theorem~\ref{thm:definiteclass822} and
Theorem~\ref{thm:signedseries823}. Theorem~\ref{thm:gluedcost811}
makes \(M^{E_j}(\mathcal Y_{i,\vartheta}(h)\odot\zeta_{i,y})\)
integral at \(z_{i,h,y}\). Stabilizer actions preserve this lattice:
in its finite frame they are integral level elements, and the
level character has values in roots of unity. Equivariant transport
takes the local lattice to the lattice at the transported point.
The definition of \(\Delta_z\) uses an unnormalized sum and the
lemma gives unit weights on the two block class sets. Thus the
whole map preserves integrality at the stated prime, including
when that prime divides \(M\).

We next compute the adjoint by testing against an arbitrary tuple.
The class evaluation formula and the definition of the single lift
express \(\langle G_\vartheta\zeta,\Theta_{j,\vartheta}^{\rm cl}F\rangle\)
as \(M^{E_j}\) times the sum of integrals against \(F\) of the scalar
theta kernels at \(z_{i,h,y}\), tested by the conjugate Cartan
coefficient. To check the real test, apply the mixed Cartan
inclusion of Lemma~\ref{lem:dualcartan811}. On the vector \(u\odot v\)
it gives the tensor of the two included vectors. Each flag
polynomial is read by \(T_c\) on the variables of its sign.
For evaluation at a matrix \(B\), the defining product of Gram
minors for \(T_c\) shows that including rows by a coordinate
matrix \(I\) replaces \(X\) by \(I^{\mathsf t}X\).
Evaluations span the dual coefficient space, so this substitution
holds on all its vectors. The matrix entries of these operations
are real. The mixed test consequently equals the product of
\(\mathcal L(\overline{\mathcal Y_{i,\vartheta}(h)})\) on the
positive block and the flag polynomial of
\(\overline{\zeta_{i,y}}\) on the negative block. The real Gaussian
also factors.

On the block target groups the source character is
\(\upsilon_1^{m-j}\) on each block determinant. On the shared
source the two target characters multiply to the target character
of the original pair. Indeed the first includes the splitting
power \(j+(t+k-1)/2\) in addition to the lift character, the
second has power \(m/2-j\), and their exponents sum to
\((s+t)/2\). The new tame character occurs only in the first
factor. Thus the splitting and residue choices in
Corollary~\ref{cor:gluedperiod811} apply. Its change of variables
in the good source integral replaces translation by \(\gamma_i\)
with the nonzero factor \(C_{\gamma_i}N_{\mathcal M}\) times the
two-block test. Both factors are independent of the new twist.
Every source integral is absolutely convergent since \(F\) is
cuspidal and the fixed theta kernels have moderate growth.
The class sets and the old congruence group in the change of
variables are finite.

Let \(\theta^\circ_{\bar\zeta}(g,y)\) denote the complementary kernel
with the polynomial of \(\bar\zeta\). By the lemma and
Theorem~\ref{thm:arithmeticclass824}, summing the first-block kernel
over \(h\) gives \(S_j L_{i,\vartheta}(g)\) exactly. Hence
\[
 \begin{split}
 b_{j,\vartheta}\langle\zeta,
           G_\vartheta^*\Theta_{j,\vartheta}^{\rm cl}F\rangle
  =M^{E_j}N_{\mathcal M}S_j
     \sum_{i,y} C_{\gamma_i}\!
       \int_{[G]} F(g)\,\mathcal B_i(g)\,
                  \theta^\circ_{\overline{\zeta_{i,y}}}(g,y)\,dg .
 \end{split}
\]
Here Theorem~\ref{thm:thirdglobal843} was used inside each integral.
The right side is conjugate-linear in \(\zeta\) and linear in \(F\).
It defines \(\langle\zeta,P_jF\rangle\) and contains no twist.
This gives the claimed equality of maps.

If \(P_jF=0\), take tuples supported at one \((i,y)\) and vary that
entry. Their polynomials include the compact translates used in
Corollary~\ref{cor:complementreproduce812}: evaluation at the
identity gives the product of leading minors, and equivariance of
the Gram-minor polynomial gives each translate. The constants
\(M^{E_j}N_{\mathcal M}S_j C_{\gamma_i}\) are nonzero.
Theorem~\ref{thm:complexseparation818} now implies \(F=0\).
Taking the adjoint followed by the map in the displayed operator
identity gives the Gram determinant formula.
\end{proof}

This result gives integral test forms for the projected periods.
To use the periods in a covolume estimate one must also control
the denominators of projection onto the good lift inside the
ambient integral space.
\clearpage

\subsection{A positive integral operator and its height at the trivial twist}
\label{sec:integralcomparison845}
Use the map \(G_\vartheta\) of Theorem~\ref{thm:gluedperiodmap844}.
Write \(f=[F:\mathbb Q]\), \(n_0=s+t\), \(N=t-9m-1=\dim V^\circ\)
and \(H=s^8\). The coefficient metrics below are the Gram-minor
metrics of Section~\ref{sec:stackedcoordinates819}, with the
determinant line of unit length in the standard frame. We take the
class-sum metric on the target and the sum metric on \(T_j\).
Set \(B_0=|\operatorname{GL}_{2fn_0}(\mathbb F_3)|\).

\begin{lemma}[Duals in the transported coefficient lattices]
\leavevmode\label{lem:glueddual845}
There is a positive integer \(h_j^-\) depending only on the
complementary diagram and \(N\) with the following properties.
At any rational prime, multiplication by \(h_j^-\) sends the
metric dual of the tuple lattice into the tuple lattice. For the
definite target, multiplying the class pairing of two integral
forms by \(B_0\) makes it integral. These assertions hold
simultaneously at the places of a coefficient field above the
rational prime. One may take
\[
 \log h_j^-\ \le\ 3L\log\max(1,L),\qquad L=2fHt^2 .
\]
\end{lemma}
\begin{proof}
First work in the identity frame. If \(D=C^*\) is a flag-minor
Weyl lattice, its Gram map is \(T_c:D\to C\), obtained by
contracting an integral kernel. Testing against the dual basis
of \(C\) shows that its metric dual amounts to
\(T_c^{-1}(C)\), with complex conjugation on the tested
argument. For the negative diagram at the \(f\) embeddings,
Lemma~\ref{lem:complementfischer819} provides an integer
\(h_j^-\) clearing this inverse. The number of boxes in its
complementary diagram is
\(\sum_{\tau,h}(N-h)c^-_{\tau,h}\).
There are at most \(2H\) columns at each of fewer than \(s\)
positive heights, and \(s,N\le t\), by
Theorem~\ref{thm:gluedcost811}. Hence this number is at most \(L\).
This gives the asserted logarithmic bound, including the empty
diagram. The ordinary mixed target Gram pairing is integral on its
standard Weyl lattice as well, directly by Cauchy--Binet on its
Gram-minor kernel.

For an algebraic formulation of the metric dual, enlarge the
coefficient field to a finite normal field \(\mathcal K\) containing
all entries under consideration. Use its ring of integers tensored
with \(\mathbb Z_\ell\). Complex conjugation permutes the factors
of this ring, and the pairing is conjugate-linear in one argument.
At a finite class one may bring its frame into the integral group
simultaneously above \(\ell\) by a rational unitary change of
coordinates. The Cayley chart and weak approximation prove this as
in Theorem~\ref{thm:definiteclass822}. The rational change is an
isometry of the algebraic Hermitian pairing and transports the
local lattices. The finite level character introduces only units.
Thus the assertions about the coefficient pairing and its dual
transfer to the moving frames, including the contragredient action
on the negative block. Tuple classes have unit mass by
Lemma~\ref{lem:gluedrow844}.

A target class stabilizer acts faithfully on a lattice of
\(\mathbb Z\)-rank \(2fn_0\). It is finite by positivity. Reduction
modulo three embeds each such group in
\(\operatorname{GL}_{2fn_0}(\mathbb F_3)\). Indeed, for an element
\(I+3^a U\) with \(a\ge1\), \(U\not\equiv0\bmod3\), a power of
prime order exponent different from three has first nonzero term
at valuation \(a\), whereas cubing has first nonzero term at
valuation \(a+1\). Thus the kernel has no element of prime order.
The stabilizer order divides \(B_0\), so \(B_0\) clears every
reciprocal mass in the target class pairing.
\end{proof}

\begin{proposition}[Uniform height of the untwisted integral map]
\leavevmode\label{prop:gluedheight845}
At every embedding \(\rho\) of the coefficient field one has
\[
  \log\max(1,\|\rho(G_1)\|_{\rm op})
       \ll_{q_0,d_{\rm K}} f m b_x\log b_x
\]
on the given sufficiently distant intervals. Here \(1\) denotes
the trivial new twist. The two fixed old marks are retained.
\end{proposition}
\begin{proof}
The Cartan product on disjoint blocks has norm
\(\|u\odot v\|=\|u\|\|v\|\). To see the normalization, the mixed
Cartan inclusion of Lemma~\ref{lem:dualcartan811} is equivariant
for the compact unitary group. The irreducible mixed highest-weight
space has an invariant Hermitian metric unique up to scale.
Evaluation at the identity in the mixed flag model maps to the
tensor of evaluation at the identity in the positive model and at
the reversal permutation in the negative model. Each evaluation has
norm one in its Gram-minor kernel, since the defining frame is
unitary. This fixes the scale of the inclusion. The two coordinate
block inclusions are themselves isometries: their columns are
orthonormal, so the matrices of pairwise Gram minors of evaluating
frames are preserved. Such evaluations span the dual flag space.
This reasoning applies to the tensor over embeddings. All maps
between the standard coefficient bases are rational, so the norm
identity applies to conjugate vectors as well.

In the class metric,
\(\|\Delta_z(v)\|\le|\Gamma_z|^{1/2}\|v\|\), by the triangle
inequality at \(z\) and the weight \(1/|\Gamma_z|\).
The same bound holds on conjugating the coefficients. Rational
unitary matrices remain unitary at every conjugate embedding of
the CM field, and finite character values remain roots of unity.
The matrices of \(G_\vartheta\) in class representatives are
algebraic: its coefficients are algebraic by
Theorem~\ref{thm:definiteclass822}, and the changes between
representatives act through rational isometries and finite
characters. For \(\vartheta=1\), that theorem bounds each squared
norm \(\|\rho(\mathcal Y_{i,1}(h))\|^2\) by
\(\exp(Cfmb_x\log b_x)\). Each individual tuple entry has norm
at most the total tuple norm.

We bound the number \(J\) of triples \((i,y,h)\) in
\(G_1\). All three class sets (the short-row class, the
complementary class, and the first-block refined class) have
ranks at most \(n_0\). At the trivial new twist their levels
contain the principal subgroups modulo the squared old marks.
By the full-level count in Lemma~\ref{lem:mass32} the logarithm
of the number of full-level classes in any of these ranks is
\[
       O\bigl((2fn_0)^4\log(2(2fn_0)q_1d_{\rm K})\bigr),
\]
since \(E\subseteq\mathbb Q(\zeta_{q_1d_{\rm K}})\). Over a
full-level class the number of classes at the finer level is at
most the index of the finer subgroup. At a split old mark in
rank \(u\) this index is at most \(q_v^{2u^2}\), by reduction
of all matrix entries modulo \(\pi_v^2\). As
\(\prod_{v\in\mathcal M}q_v\le M^f\), this adds at most
\(2f n_0^2\log M\) to each class count in the logarithm.

The number of old residue matrices is at most
\(|R|^{3mp}\le M^{6fmp}\).
Choose once the \(6m+1\) short rows for each short-row class
as in Theorem~\ref{thm:stackedheight819}. There are at most
\((6m+1)^{2mf\Delta'}\) choices of their ordered lists,
where \(\Delta'=s^8-(t+1)/2\le H\). Thus the sum of the
three logarithmic class bounds, \(6fmp\log M\), and
\(2mfH\log(6m+1)\) bounds \(\log J\).
On these intervals \(f\ll x^{12}\), \(s\ll x^5\),
\(m\ll\sqrt x\), \(q_1\ll_{q_0}x^{11}\log(2x)\),
\(\log M\le22\log b_x\), and \(b_x\ge x^{400}\).
In particular \((fn_0)^4\ll x^{68}\).
It follows that both \(\log J\) and
\(\log B_0\le(2fn_0)^2\log3\) are \(O_{q_0,d_{\rm K}}(fmb_x\log b_x)\).
Now sum the bounds for the \(J\) class vectors, including
\(M^{E_j}\), whose logarithm satisfies the same estimate by
Theorem~\ref{thm:gluedcost811}. This proves the proposition.
\end{proof}

\begin{theorem}[The positive integral comparison operator]
\leavevmode\label{thm:integralcomparison845}
On the definite target put
\[
        A_\vartheta=B_0 h_j^-\,G_\vartheta G_\vartheta^*.
\]
This is an algebraic positive semidefinite operator preserving the
target lattice at every prime away from the new conductor. Uniformly
at every complex embedding,
\[
   \log\max(1,\|\rho(A_1)\|_{\rm op})
                \ll_{q_0,d_{\rm K}} f m b_x\log b_x .
\]
With the good source and the common map of
Theorem~\ref{thm:gluedperiodmap844}, its pullback is
\[
 (\Theta_{j,\vartheta}^{\rm cl})^*A_\vartheta
       \Theta_{j,\vartheta}^{\rm cl}\big|_{\mathcal H_j}
       = \frac{B_0h_j^-}{|b_{j,\vartheta}|^2}\,P_j^*P_j ,
\]
where the adjoint on the left is taken for the lift restricted
to \(\mathcal H_j\). The right side is positive definite there.
\end{theorem}
\begin{proof}
Fix an allowed \(\ell\) and use the product of its coefficient
completions as in the lemma. For any integral \(\Phi\) in the
target and integral \(\zeta\) in the tuple lattice,
\[
    B_0\langle\zeta,G_\vartheta^*\Phi\rangle
             =B_0\langle G_\vartheta\zeta,\Phi\rangle
\]
is integral by Theorem~\ref{thm:gluedperiodmap844} and the lemma.
The dual-lattice assertion now gives \(B_0h_j^-G_\vartheta^*\Phi\)
integral in the tuple space. Applying \(G_\vartheta\) proves the
target assertion. All lattices and linear maps split by the
idempotents of the coefficient completions, so this also proves
preservation at each individual place.

For the archimedean assertion take full class tables in fixed
representatives. Their Gram matrices \(H_T,H_A\) have rational
entries from Gram minors and class masses. The matrix of the
adjoint on the target subspace is the restriction of
\(H_T^{-1}\overline{G_1}^{\,\mathsf t}H_A\) on full tables.
The latter matrix has algebraic entries. Write \(\rho^c(a)=
\overline{\rho(\bar a)}\). Under \(\rho\) it becomes the
adjoint of \(\rho^c(G_1)\) on the corresponding full tables.
It follows that
\[
       \|\rho(A_1)\|_{\rm op}
        \le B_0h_j^-\,\|\rho(G_1)\|_{\rm op}\,
                               \|\rho^c(G_1)\|_{\rm op}.
\]
The proposition applies to both embeddings.
Here \(L=2fHt^2\ll x^{62}\), so the lemma absorbs
\(\log h_j^-\) in the asserted bound. Finally the quadratic
form of \(G_\vartheta G_\vartheta^*\) on a vector \(u\) is
\(\|G_\vartheta^*u\|^2\). Insert the identity of
Theorem~\ref{thm:gluedperiodmap844} and polarize on the good
source. Its injectivity proves positive definiteness.
\end{proof}
\clearpage

\subsection{The split Hecke root and first occurrence in the descent tower}
\label{sec:descenttower846}
The integral comparison in Theorem~\ref{thm:integralcomparison845}
will be applied to the true lift range \(Z\). In
Proposition~\ref{prop:testedkernel785} a false lift at a good split
place \(w\) belongs to \(eV\cap\ker A_w\cap Z^\perp\).
The present argument uses the equation \(A_w=0\) on its
constituents. It applies even when only one split place with this
equation is specified.

\begin{lemma}[A bound for all entries of a unitary spherical parameter]
\leavevmode\label{lem:unitarysatake846}
Let \(k\) be a nonarchimedean local field with residue cardinality
\(q\). If an irreducible unitary spherical representation of
\(\mathrm{GL}_b(k)\) has unitary-normalized Satake entries
\(\alpha_1,\ldots,\alpha_b\), then
\[
 q^{-(b-1)/2}\le|\alpha_i|\le q^{(b-1)/2}\quad(1\le i\le b).
\]
\end{lemma}
\begin{proof}
Give the integral group mass one. The indicator of the integral
matrices of determinant valuation \(d\ge0\) is a compactly
supported spherical Hecke operator on the group. Its right cosets
correspond to sublattices of index \(q^d\). For diagonal valuations
\(a_1,\ldots,a_b\ge0\) of sum \(d\) in upper triangular Hermite
normal form, their number is \(q^{\sum_i(b-i)a_i}\): in row \(i\)
each of the \(b-i\) later entries is specified modulo
\(\varpi^{a_i}\). In normalized unramified induction the value
of the spherical function used in the Satake transform at such
a triangular representative is
\(\prod_i\alpha_i^{a_i}q^{-(b+1-2i)a_i/2}\).
Summing the cosets gives the eigenvalue
\[
 q^{(b-1)d/2}h_d(\alpha),\qquad
 h_d(\alpha)=\sum_{\sum a_i=d}\prod_i\alpha_i^{a_i}.
\]
This eigenvalue also holds on the spherical constituent by the
Satake isomorphism. In a unitary representation convolution by
an indicator has norm at most its mass. There are
\(\binom{d+b-1}{b-1}\) diagonal types, each with mass at most
\(q^{(b-1)d}\). Therefore
\(\sum_{d\ge0}h_d(\alpha)X^d\) converges for
\(|X|<q^{-(b-1)/2}\). As a rational function it is exactly
\(\prod_i(1-\alpha_i X)^{-1}\). Its numerator is one, so none
of its poles in the finite plane cancel. This proves the upper
bound, and the contragredient gives the lower bound.
\end{proof}

\begin{proposition}[A gap in the ranks of unitary local theta partners]
\leavevmode\label{prop:splitpartner846}
Let \(n=s+t\), and let \(\pi_w\) be a spherical representation at
the split place of the rank \(n\) definite group. Suppose \(A_w\)
kills its spherical line with the value \(X_w\) in
Section~\ref{sec:splitcost785}. At any rank \(b\) with \(s<b<t\)
there is no irreducible unitary local theta partner for \(\pi_w\)
with the unramified unitary splittings of the tower. The same
conclusion holds for its contragredient.
\end{proposition}
\begin{proof}
Write \(\beta_i\) for the unitary-normalized spherical entries of
\(\pi_w\). On its line the polynomial before evaluation at \(X_w\)
is \(\prod_i(1-q^{(n-1)/2}\beta_i X)\).
Indeed the \(i\)-th minuscule double coset has eigenvalue
\(q^{i(n-i)/2}\) times the \(i\)-th elementary symmetric
function, and its multiplier in \(A_w\) is \(q^{i(i-1)/2}\).
Since \(|X_w|=q^{-(t-1+s/2)}\), one of the \(\beta_i\)
has modulus \(q^{(t-1)/2}\).

At a split place the oscillator is the type II matrix
correspondence, up to unitary unramified determinant characters
and the choice of dual component in
Proposition~\ref{prop:split175}. We recall why a possibly
ramified partner cannot introduce other entries. By
\href{https://www.numdam.org/item/10.24033/asens.2080.pdf}
{M\'inguez, \emph{Correspondance de Howe explicite},
Theorem~1(2)}, the Langlands quotient data on the larger member
are the dual essentially square-integrable data on the smaller
one together with \(n-b\) characters of successive exponents
\((1-(n-b))/2,\ldots,((n-b)-1)/2\), reordered. For a spherical
larger member the unique Langlands data are all unramified
characters of rank one: its Satake multiset gives just this
Langlands quotient. Hence all the smaller data must themselves
be unramified characters. Its Langlands quotient is spherical.
The unitary determinant twists have absolute value one on the
uniformizer.

Lemma~\ref{lem:unitarysatake846} bounds the absolute values of the
smaller entries and their inverses by \(q^{(b-1)/2}\). The added
string and its inverse are bounded by \(q^{(n-b-1)/2}\). Both
bounds are strictly less than \(q^{(t-1)/2}\), since \(b<t\)
and \(n-b<t\). The same two bounds hold on the inverses of all
larger entries. This contradicts the entry already found, and
also excludes the contragredient for which its reciprocal occurs.
\end{proof}

\begin{theorem}[First occurrence for representations satisfying the split test]
\leavevmode\label{thm:descenttower846}
Let \(\pi\) be an irreducible automorphic representation on the
definite group \(U(I_n)\) having the local component in
Proposition~\ref{prop:splitpartner846}. In the rational Witt tower
through \(W_j\) use coherent splittings and let
\(W^{(a)}=W_j\oplus\mathbb H^a\), with \(\mathbb H\) a hyperbolic
plane. If the full global theta lift of \(\pi\) to \(U(W^{(a)})\)
is nonzero for some \(1\le a\le p=s+9m+1\), then its full theta
lift to \(U(W_j)\) is nonzero. This uses the fields and ranks of
the given intervals.
\end{theorem}
\begin{proof}
Here the full lift allows all Schwartz data and all input vectors
in the given automorphic copy. Coherence keeps the Weil splitting
characters fixed when two dimensions are added. In the external
norm-one convention of the manuscript, adding a plane multiplies
the smaller space's external character by \(\upsilon_1\):
base change multiplies by \(\mu^2\), cancelling the change of two
in the dimension exponent. This leaves the splitting on the
definite member compatible with passage down the tower.

Recall the tower property in this definite-source case. In the
mixed model at the parabolic of an isotropic \(h\)-plane, the
extra position variables lie in \(I_n^h\). In the constant term,
integration over the center of the unipotent radical kills every
rational tuple with nonzero Hermitian Gram matrix, by additive
character orthogonality. A zero Gram matrix forces the whole
tuple to vanish, since \(I_n\) is positive definite. On the zero
slice the remaining unipotents act trivially, and the kernel
is the lower theta kernel with Schwartz datum evaluated at zero
in the extra variables, up to the nonzero Levi and measure
factors. These evaluations are onto the lower Schwartz space:
tensor any lower datum with a Schwartz function in the extra
variables of value one at zero. The mixed-model change is an
invertible partial Fourier transform. All integrations for these
constant terms use a compact definite input quotient and a compact
unipotent quotient; at fixed adelic arguments the Schwartz sums
converge absolutely and uniformly on compact sets. Rational
parabolics are obtained from isotropic flags, and Witt's extension
theorem identifies planes of each fixed dimension. Thus a nonzero
lower lift persists at the next rank, and all proper constant
terms at the first nonzero rank vanish. This is Rallis's tower
property in the form used by
\href{https://arxiv.org/pdf/1207.4709}
{Gan, Qiu and Takeda, \S11.2}.

If the lift to \(W_j\) were zero, persistence would give vanishing
at the smaller members of its tower as well. Take the first
nonzero member through the specified \(W^{(a)}\). Its rank is
\(b=s+2a'\) with \(1\le a'\le p\). Its automorphic forms are
cuspidal by the constant-term calculation. The center is
anisotropic, so a nonzero such lift pairs nontrivially with a
vector in an irreducible cuspidal summand. This summand is
unitary. The pairing with the theta kernel and \(\pi\) gives a
nonzero local theta pairing at \(w\). More explicitly, test the
nonzero global pairing on decomposable data and vectors, fix
the factors away from \(w\), and vary the three local factors.
The resulting invariant trilinear functional is nonzero. A local
Schwartz vector has an open compact stabilizer, so its image on
the other two factors is in their smooth dual by admissibility.
This supplies the required local joint quotient, with either
dual convention covered by Proposition~\ref{prop:splitpartner846}.

Finally \(s+2p=3s+18m+2<5s=t\) for every sufficiently large
node, since \(s\ge x^5\) and \(m\le2\sqrt x\).
Thus \(s<b<t\). The proposition excludes the local pair at the
very place \(w\) where the original Hecke test was imposed.
This contradiction proves the theorem.
\end{proof}

To apply the tower result to the false-lift tests, one must pass
from vanishing against the prescribed datum at the original level
to vanishing against the Schwartz data needed for the block
period. This passage concerns the marked local tests and the
selected real weight.
\clearpage

\subsection{Finite tests at the bottom of the descent tower}
\label{sec:markedbottom847}
Fix \(W_j,I_n\) of ranks \(s,n=s+t\) on one of the intervals of
Theorem~\ref{thm:commonsource806}. Let \(S\) consist of the original
marked places and the primes of the actual new conductor, using
one component of each split pair. At a new place write \(\chi\)
for its nontrivial tame residue character. The untwisted case
has no new places. We use matrix coordinates in which the
\(s\) selected rows come first. Write \(O\) for the local
integer ring, \(\pi\) for a uniformizer, \(k=O/\pi O\), and
\(e=\binom{I_s}{0}\). On \(n\)-by-\(s\) matrices the block
function is the product of the following local functions:
\[
 \phi_{\rm old}(X)={\bf1}_{e+\pi\operatorname{Mat}_{n,s}(O)}(X),
 \qquad
 \phi_{\rm new}(X)
   ={\bf1}_{\operatorname{Mat}_{n,s}(O)}(X)\,
                         \chi(\det\overline{X'}).
\]
In the second expression take the value zero on singular top
blocks and outside the integral matrices. Here and below
\(X'\) denotes the square top block and \(X''\) the remaining
rows. The old level \(J\) in \(\mathrm{GL}_n(O)\) fixes
\(\bar e\). At a new place its reduction has zero upper-right
block, with invertible diagonal blocks. The character of \(J\)
on the form is \(d(h)=1\) at an old place and
\(d(h)=\chi(\det\bar h_{11})^{-1}\) at a new one. The operator
\(U_h\phi(X)=\phi(h^{-1}X)\) acts on the block function by this
same character. These are the physical translation conventions
of the proof of Theorem~\ref{thm:markedtype801}.

\begin{lemma}[The finite hyperplane average on the chosen types]
\leavevmode\label{lem:finitehyperplane847}
For every nonzero row \(\lambda\in k^{1\times s}\) put
\(H_\lambda=\{1+v\lambda:v\in k^{s\times1},\
1+\lambda v\ne0\}\). At each old place, let \(R(z)\) be
ordinary right translation on the level space by a lift of
\(\operatorname{diag}(z,I_t)\), for \(z\in\mathrm{GL}_s(k)\).
For every \(\Xi\in e_{\rm mark}V\), where \(e_{\rm mark}\)
is the idempotent of Proposition~\ref{prop:testedkernel785},
one has \(\sum_{z\in H_\lambda}R(z)\Xi=0\).
\end{lemma}
\begin{proof}
Translation is well defined because any two such lifts differ
by an element of the old level. In a basis with \(\lambda\) the
last coordinate row, \(H_\lambda\) is the group of matrices
\(\left(\begin{smallmatrix}I_{s-1}&u\\0&a\end{smallmatrix}\right)\)
with \(a\in k^\times\). It contains the whole unipotent radical
of the hyperplane parabolic and the last \(\mathrm{GL}_1\).
If an irreducible finite type has an \(H_\lambda\)-fixed
vector, its unipotent invariants therefore contain
\(\tau\boxtimes{\bf1}\) on the Levi for some irreducible
\(\tau\) of \(\mathrm{GL}_{s-1}(k)\). Indeed the last-factor
invariants are stable under the first factor and finite-group
representations over \(\mathbb C\) are semisimple.
Frobenius reciprocity and transitivity of parabolic induction
put the trivial degree-one character in its finite cuspidal
support. For \(s=1\) the assertion just says the character is
trivial. By the finite support calculation in
Theorem~\ref{thm:markedtype801}, the ordinary translation types
occurring in the marking have no such character. More explicitly,
on each orbit the marked label is tensored by the inverse of
the old determinant character, of order dividing \(Q\).
A degree-one label has order not dividing \(Q\), and a larger
orbit preserves its degree under this twist. Thus the average
over \(H_\lambda\), the projector onto its invariants times
its order, vanishes on every summand of \(e_{\rm mark}V\).
\end{proof}

\begin{theorem}[All finite marked data from the block test]
\leavevmode\label{thm:markedbottom847}
Let \(\Xi\in e_{\rm mark}V\). Fix a real Schwartz vector and
one matrix coefficient of \(\Xi\). Use the lattice vacuum
outside \(S\). If its descent from \(I_n\) to \(W_j\) by the
product block function is zero at every adelic argument, then
it vanishes for every Schwartz--Bruhat function on the product
of the matrix spaces at \(S\), at every argument. The real
vector and the coefficient in these two assertions are the same.
\end{theorem}
\begin{proof}
At argument \(g\) denote the integral with arbitrary marked
data by \(B_g(\phi,\Xi)\), linear in \(\phi\) and conjugate
linear in \(\Xi\). Use the split matrix model, with the external
character on the \(W_j\) argument removed if present; it is a
nonzero scalar at \(g\). The intrinsic splitting on integral
left translations in these coordinates is unramified.
Consequently simultaneous integral left translation of the test
and right translation of the form gives
\[
              B_g(U_h\phi,R(h)\Xi)=B_g(\phi,\Xi).
\]
This follows by the substitution \(u\mapsto uh\) in the compact
definite quotient integral, since the oscillator action on the
test is pullback by \(h^{-1}\). Right multiplication of its matrix
argument by \(A^{-1}\in\mathrm{GL}_s(F_v)\) is the oscillator
action of an element on the \(W_j\) side up to a nonzero
determinant scalar. Such elements at the different split places
are independent and can be taken to be the identity elsewhere.
Thus the hypothesis also kills all products of these right
translates of the block function, for all \(g\).

First consider a point \(X\) at a marked place where \(A=X'\)
is invertible over the local field and \(T=X''A^{-1}\) lies
in \(\pi\operatorname{Mat}_{t,s}(O)\) at an old place, or
in \(\operatorname{Mat}_{t,s}(O)\) at a new one. Then
\(X=\binom{I_s}{T}A\) belongs to \(J eA\). This orbit is exactly
the support of \(\phi_{\rm old}(\,\cdot\, A^{-1})\) or
\(\phi_{\rm new}(\,\cdot\, A^{-1})\), respectively. To verify
the reverse inclusion, write \(Y A^{-1}=\binom{C}{D}\) on this
support. Completing it by the remaining standard columns gives
\(\left(\begin{smallmatrix}C&0\\D&I_t\end{smallmatrix}\right)\in J\).
The orbit is compact and open. For a ball about \(X\) contained
in it, average its indicator with the operator
\(\int_J d(h)^{-1}U_h\,dh\) for probability Haar measure.
This leaves the pairing with \(\Xi\) unchanged by covariance.
The result is supported on one orbit and transforms by \(d\).
Such functions form a line, since their values are determined by
the value at \(eA\) and the translated block function is nonzero
there. Thus the average is a scalar multiple of this translate.

Suppose now the preceding conditions fail. There is an integral
row \(l=(l_1,l_2)\) with \(lX=0,\ \bar l_1\ne0\) and, at a new
place, \(\bar l_2=0\). If \(A\) is singular, scale a nonzero left
kernel vector to be primitive and take \(l_2=0\). Otherwise take
a row \(T_i\) offending the integral condition. Divide
\((-T_i,e_i^{\rm t})\) by an entry of \(T_i\) of largest absolute
value, where \(e_i^{\rm t}\) selects row \(i\) of \(X''\).
This value is at least one at an old place and strictly greater
than one at a new place. The row has all the required properties.

Put \(\lambda=\bar l_1\). For each \(z=1+v\lambda\in H_\lambda\)
lift \(v\) to \(\tilde v\in O^s\) and put
\[
           h_z=I_n+\binom{\tilde v}{0}\,l .
\]
Its determinant \(1+l_1\tilde v\) is a unit. It fixes \(X\)
exactly and preserves every ball about \(X\) taken in the
supremum norm, since both \(h_z\) and its inverse are integral.
At an old place \(h_z\) has the same first residue columns as
\(\operatorname{diag}(z,I_t)\). Its translation on \(\Xi\)
is \(R(z)\Xi\). At a new place \(h_z\in J\) and its translation
is \(\chi(\det z)^{-1}\Xi\). Now \(1+\lambda v\) runs over
\(k^\times\), so some \(h_z\) acts by a scalar different from
one in the new case. For a product ball with this component,
covariance gives the same pairing with \(R(h_z)\Xi\) as with
\(\Xi\). At an old place sum this equality and apply
Lemma~\ref{lem:finitehyperplane847}; at a new place use the
nontrivial scalar. Both give zero.

Finally give the product of the marked matrix spaces the nested
basis of balls obtained by taking the same integer depth at
each place. At every point where a compactly supported locally
constant test is nonzero, choose a ball on which it is constant.
If all components have the indicated integral ratios, shrink
the ball into the product of the open orbits. Its pairing
vanishes by averaging at every place and applying the block
hypothesis. If one component fails, its stabilizers already
give zero for the entire ball. Compactness gives a finite
cover of the support. Two balls from this nested basis which
meet are nested, so the maximal members form a disjoint cover.
Linearity proves the assertion. When \(S\) is empty it is the
given block vanishing itself.
\end{proof}
\clearpage

\subsection{The descent tower with the unramified tests fixed}
\label{sec:restrictedtower848}
Write \(S\) for the finite set from Section~\ref{sec:markedbottom847}.
At the real places and at \(S\) allow arbitrary Schwartz data.
At every other finite place keep the original lattice vacuum.
Call this space of data \(\mathcal S_S(B)\) for a member \(B\)
of the coherent Witt tower through \(W_j\). The variable \(g\)
in a descent on \(U(B)\) still ranges over its whole adelic
group. Fix a matrix coefficient \(y\) of one of the finite-level
algebraic forms on \(U(I_n)\), with its given real type. Write
\[
        D_B(\phi,y)(g)
        =\int_{[U(I_n)]}\theta_B(g,h;\phi)\overline{y(h)}\,dh .
\]
The compactness of the integration domain gives continuity in
the real Schwartz topology for each \(g\), by uniform convergence
of the theta sum on compact sets of group arguments.

\begin{lemma}[Transfer of the full Hecke polynomial on the lattice test]
\leavevmode\label{lem:rectangularpolynomial848}
At a split \(w\notin S\), write \(q=Nw\). In ranks \(b<n\)
put \(f={\bf1}_{\operatorname{Mat}_{n,b}(O_w)}\), in lattice
coordinates. Give the integral groups mass one. On rank \(d\)
let \(C_{d,i}^{\pm}\) be convolution by the unnormalized
elementary double coset with \(i\) elementary divisors
\(\varpi^{\pm1}\) and the rest one, and set
\(\mathsf E_{d,i}^{\pm}=q^{-i(d-i)/2}C_{d,i}^{\pm}\).
The two group actions on \(f\) satisfy
\[
 \sum_{h=0}^n(-Y)^h\mathsf E_{n,h}^{+} f
 =
 \prod_{i=0}^{n-b-1}(1-Yu q^{(n-b-1)/2-i})
       \sum_{l=0}^b(-Yu/v)^l\mathsf E_{b,l}^{-} f .
\]
Here \(u\) is the unitary phase of a positive elementary
divisor on the left and \(v\) that of a negative one on the
right, including the determinant twists of the chosen splitting.
\end{lemma}
\begin{proof}
The changes of variables are \(X\mapsto h^{-1}X\) on the left
and \(X\mapsto Xg\) on the right in the split model. Their
square-root Jacobian multipliers in the two cosets are
\(q^{bh/2}\) and \(q^{nl/2}\), respectively, where \(h\) in
the exponent counts elementary divisors. The coset phases are
\(u^h,v^l\). For a matrix of residue rank \(a\), the numbers of
contributing right cosets are
\({n-a\brack h}_q\) and \({b-a\brack l}_q\). The first count
chooses a subspace containing the column space, the second an
\(l\)-space in the kernel. Both sides vanish off the integral
box. With \(D=n-b\), the Grassmannian calculation of
Lemma~\ref{lem:secondhecke829} gives
\[
 {n-a\brack h}_q=\sum_l {b-a\brack l}_q {D\brack h-l}_q
                                       q^{l(D-h+l)} .
\]
Divide by the two Hecke normalizations and include the
Jacobians. For \(j'=h-l\) the coefficient multiplying
\(\mathsf E_{b,l}^{-} f\) is
\(u^h v^{-l}{D\brack j'}_q q^{-j'(D-j')/2}\).
Indeed the exponent is
\((bh-h(n-h)-nl+l(b-l))/2+l(D-h+l)\).
The normalized Gaussian coefficient is the elementary symmetric
function of degree \(j'\) in \(q^{(D-1)/2-i}\), \(0\le i<D\),
by expanding the product successively. Multiplying the two
generating polynomials proves the identity.
\end{proof}

\begin{lemma}[Constant terms and the prescribed vacuum]
\leavevmode\label{lem:restrictedconstant848}
Let \(B'=B\oplus\mathbb H^h\) be two members of the tower, using
the diagonal sign lattices. A constant term of
\(D_{B'}(\phi,y)\) along a parabolic stabilizing a rational
isotropic \(h\)-plane, at any adelic argument, is a nonzero Levi
and measure factor times a lower descent with datum in
\(\mathcal S_S(B)\), if \(\phi\in\mathcal S_S(B')\).
At the identity of the extra Levi and compact factors the
zero-slice map onto \(\mathcal S_S(B)\) is surjective.
\end{lemma}
\begin{proof}
We first choose an integral rational splitting of the standard
tower. The relative extension \(E/F\) is unramified at all
finite places, including the split ones, by
Lemma~\ref{lem:signkernel823}. The trace ideal of
\(\mathcal O_E\) is the unit ideal. Locally this follows at a
split prime by taking \((1,0)\), and in an unramified quadratic
field by the surjective finite-field trace; an element of
unit trace generates the local trace ideal. Hence there is
\(\lambda\in\mathcal O_E\) with \(\lambda+\bar\lambda=-1\).
For a positive and a negative coordinate with vectors \(u,v\)
the two columns \(u+v,\ u+\lambda(u+v)\) have norms zero,
pairing one and determinant \(-1\). They thus identify the
diagonal lattice of this pair with the integral hyperbolic
lattice at every finite place. The rational anisotropic kernel
of \(W_j=I_{k+j-1}\oplus(-I_{m-j})\) is the remaining definite
summand. Its definiteness bounds the Witt index at a real place,
so these pairs account for every rational rank in the tower.

For the standard parabolic write \(g=p k\) by local Iwasawa
decomposition. Away from \(S\) use the lattice stabilizer for
the compact factor, which is hyperspecial. In the split case
this is the integral basis theorem for flags. In the unramified
field case intersect an isotropic plane with the lattice to
obtain a primitive isotropic summand. Unimodularity gives
integral dual partners. Subtracting linear combinations of the
isotropic basis makes their Gram matrix zero: solve the
off-diagonal entries in pairs and use the surjectivity of the
integral trace on the diagonal. The orthogonal complements are
unimodular and integrally isometric by the diagonalization in
Section~\ref{sec:commonsource806}. Thus the compact group is
transitive on the planes, giving the asserted decomposition.
The other finite and the real compact factors can act freely
on the allowed data.

In the mixed model use \(h\) positions in \(I_n\) in addition
to the lower oscillator. As in Theorem~\ref{thm:descenttower846},
integration over the center of the radical selects zero
Hermitian Gram matrix, hence only the zero tuple, because
\(I_n\) is definite. The other radical coordinates act trivially
there. The remaining Levi action gives a scalar and translation
of the lower argument. At an unmarked finite place the compact
factor fixes the vacuum. The integral splitting identifies its
self-dual oscillator lattice with an orthogonal sum. Its vacuum
in the mixed model is the lower vacuum tensored with the lattice
indicator in the extra positions, up to a nonzero scalar.
Indeed partial Fourier transform of an integral lattice
indicator is its volume times the indicator of the dual lattice,
and the chosen measure is self-dual. This argument also works in
the scaled position coordinates of Lemma~\ref{lem:sparsedistance814}.
The extra indicator takes value one at zero. Almost everywhere
the changes of measure have scalar one. At the other places
extension from the zero slice is obtained by tensoring a lower
test with any Schwartz function of the extra positions having
value one at zero, then inverting the mixed-model change.
This proves both assertions by taking products and linear
combinations. Witt's extension theorem and automorphy reduce
any rational plane to the standard plane.
\end{proof}

\begin{theorem}[Propagation of vanishing with the same finite tests]
\leavevmode\label{thm:restrictedtower848}
Suppose \(y\) is spherical at one split \(w\notin S\) with
\(A_w y=0\) for the test and root of
Proposition~\ref{prop:testedkernel785}. Suppose
\(D_{W_j}(\phi,y)=0\) for all \(\phi\in\mathcal S_S(W_j)\).
Then for every \(1\le a\le p=s+9m+1\) and every
\(\phi\in\mathcal S_S(W_j\oplus\mathbb H^a)\) one has
\(D_{W_j\oplus\mathbb H^a}(\phi,y)=0\), at every argument.
\end{theorem}
\begin{proof}
Surjectivity of the zero slice first gives vanishing below
\(W_j\): take a constant term of the identically zero lift
at a Levi argument giving any prescribed lower argument.
If a larger member up to the specified rank has a nonzero
lift, choose the first such member \(B\) and put \(b=\dim B\).
All proper maximal constant terms of every allowed lift there
vanish by Lemma~\ref{lem:restrictedconstant848}. A constant
term along a flag is an iterated such integral. By continuity
and the density of Hermite polynomials times the Gaussian in
the Schwartz space, choose an allowed \(\phi\) with polynomial
Fock vector at infinity and \(F=D_B(\phi,y)\ne0\).

This function is cuspidal, of finite level and real compact
type. For clarity, it has moderate growth: on a fixed finite
level component its theta functional is a tempered lattice
sum applied to the moving real test. A fixed finite collection
of Schwartz seminorms bounds this sum with \(h\) ranging in
a compact set of representatives of the definite quotient.
These seminorms after oscillator translation grow polynomially
in the group norm: in the equivalent oscillator Sobolev norms
one commutes finitely many position and momentum operators
past the unitary action, which transforms each by the linear
symplectic matrix. Restriction to the real compact input
type of \(y\) gives a single infinitesimal character on each
factor of \(U(B)(F_\infty)\) by the compact dual-pair
decomposition (\href{https://msp.org/pjm/1987/129-1/pjm-v129-n1-p02-s.pdf}
{Davidson, \S2, Theorem~2.7}). Thus the classical rapid decay
theorem for cusp forms on Siegel sets makes \(F\) square
integrable; the center is anisotropic. One may apply it on the
finitely many arithmetic components at this level
(\href{https://arxiv.org/pdf/1106.2149}{Miller and Schmid,
Theorem~A and its adelic formulation}).

Normalize \(Y=q^{(n-1)/2}X_w\), so
\(|Y|=q^{-(t-1)/2}\). The Hecke polynomial of the hypothesis
is \(\sum_h(-Y)^h\mathsf E_{n,h}^+\).
It is normal on the definite spherical space since inverse
double cosets are its Hecke adjoints and all spherical operators
commute. Consequently its adjoint kills \(y\) as well.
Move this adjoint from the conjugated \(y\) through the definite
integral. Complex coefficients are conjugated and inverse
cosets are reversed by the change of integration variable. This
puts the polynomial with coefficients \((-Y)^h\) on the left
action of the local matrix vacuum. Lemma~\ref{lem:rectangularpolynomial848}
transfers it to the right action. The latter is precisely adelic
right translation of \(F\) on the smaller group, because acting
on the test before applying the theta functional replaces the
argument \(g\) by the corresponding right translate. We obtain
\[
 0=\prod_{i=0}^{n-b-1}(1-Yu q^{(n-b-1)/2-i})
          \sum_{l=0}^b(-Yu/v)^l\mathsf E_{b,l}^- F .
\]
Here \(b>s\), so \(n-b<t\), and the scalar product is nonzero.

The remaining polynomial is invertible on the spherical
\(L^2\)-space at these ranks. Its constant term is the identity.
Each \(\mathsf E_{b,l}^-\) has norm at most its mass,
\({b\brack l}_q q^{-l(b-l)/2}\), by unitarity of right
translation. The sum of the norms of the positive-degree terms
is thus at most
\[
 \prod_{i=0}^{b-1}\bigl(1+q^{-(t-b)/2-i}\bigr)-1
 \ \le\ \exp\bigl(q^{-(t-b)/2}/(1-q^{-1})\bigr)-1 .
\]
Since \(b\le s+2p=3s+18m+2\), \(t=5s\), \(m\le2\sqrt{x}\)
and \(s\ge x^5\), one has \(t-b\ge4\) throughout the sufficiently
large nodes. Already for \(q=2\) the upper bound is at most
\(\exp(1/2)-1<1\). Taking norms gives \(F=0\), a contradiction.
The argument uses the same fixed coefficient throughout, with
the lattice test at \(w\) on every member of the tower.
\end{proof}
\clearpage

\subsection{Vanishing of the descent of a false lift}
\label{sec:falsedescent849}
Fix one of the fields, intervals, lift indices and characters of
Theorem~\ref{thm:integralcomparison845}. Set
\(G=U(W_j)\), \(K'=U(I_n)(F_\infty)\), and let \(V_\rho\)
be the irreducible coefficient module of the definite target
at infinity, with its invariant Hermitian metric. For a
unit vector \(\zeta\in V_\rho\) write
\[
       y_\zeta(h)=\langle\zeta,\rho(h_\infty^{-1})\Xi(h_f)\rangle .
\]
Use an inner product linear in the second argument, and write
\(\zeta_0\) for the specified highest vector of length one.
Include in the standard block function the vacuum at every
unmarked place. Write \(\Phi\) for its tensor product with the
single lift's joint highest harmonic at infinity, chosen to
match the coefficient \(y_{\zeta_0}\). Put
\(d_\Xi(g)=\overline{D_{W_j}(\Phi,y_{\zeta_0})(g)}\), restoring
both external characters of the single lift in this formula.
The character on the \(G\) argument can subsequently be removed
by pointwise multiplication by its inverse.

\begin{proposition}[The block descent in the marked source space]
\leavevmode\label{prop:actualadjoint849}
If \(e_{\rm mark}\Xi=\Xi\), then \(d_\Xi\) belongs to the
original good source \(\mathcal H_j\). With the fixed measures
and joint harmonic normalization, it is the Hilbert adjoint of
\(\Theta^{\rm cl}_{j,\vartheta}\) on this space up to a common
nonzero scalar. In particular \(d_\Xi=0\) for
\(\Xi\in e_{\rm mark}V\cap Z^\perp\).
\end{proposition}
\begin{proof}
On the finite source use the level in
Theorem~\ref{thm:commonsource806}. Right multiplication of a
marked matrix by an old integral element of residue \(z\)
takes the point indicator to its left translate by
\(\operatorname{diag}(z^{-1},1)\), in the convention \(U_h\)
of Theorem~\ref{thm:markedbottom847}. The source splitting
supplies the inverse old determinant character in \(D\).
Conjugation in \(d_\Xi\) changes it to the direct character.
Changing variables on the definite group moves the left
translate to translation of \(\Xi\) by
\(\operatorname{diag}(z,1)\). Thus the right action on
\(d_\Xi\) is the marked action on \(\Xi\). In particular the
source principal subgroup fixes \(d_\Xi\). At a new prime the
factor \(\chi(\det\overline{X'})\) under right multiplication
by \(k\) acquires \(\chi(\det\bar k)\), cancelled by the
source splitting there. All other places have the fixed
hyperspecial test. This gives precisely the finite source
level, and applying the old central idempotent gives
\(e_{\rm mark}d_\Xi=d_\Xi\).

We spell out why the finite assertion already forces
cuspidality. Decompose the finite marked action into tuples
as in Lemma~\ref{lem:twoelliptic806}. Each tuple has an
irreducible finite cuspidal label at some original split
place \(v\). For a proper rational parabolic \(P\) its
localization is a proper parabolic of \(\mathrm{GL}_s(F_v)\).
The invariants of that finite cuspidal type under the residue
unipotent radical vanish by finite Frobenius reciprocity. To
evaluate the \(P\)-constant term at \(g\), write \(g_v=p_vk_v\)
with \(p_v\in P(F_v)\) and \(k_v\) integral. Averaging the
input by \(k_v^{-1}(N_P(F_v)\cap K_v)k_v\) leaves that constant
term unchanged: after multiplication by \(g_v\) this is left
translation by elements of the adelic unipotent radical, under
which the constant term is invariant. The same average is zero
on the finite type, since the integral radical maps onto the
residue radical. This proves vanishing for each tuple and all
proper \(P\). The descent is smooth of moderate growth by the
compact integral and seminorm argument of
Theorem~\ref{thm:restrictedtower848}.

We verify its real line at one embedding; tensor over the real
places. Write \(a=k+j-1,\ b=m-j,\ a+b=s,\ b'=(t-1)/2\),
and suppress the embedding in \(w_j,\sigma_i\) from
Theorem~\ref{thm:gluedcost811}. The positive and negative
column determinants in the polynomial use disjoint target rows.
Every source mixed lowering contraction kills the polynomial.
Triangular substitution on the ordered compact blocks makes
it a highest vector. The source action on a positive column is
dual and on a negative column is direct. The target character
has external weight \(-w_j+n/2\) after the rank twist and enters
inversely on the source. The vacuum further contributes one to
each negative column for each of the \(n\) target rows. Thus the
vacuum exponents are \(w_j-n/2,w_j+n/2\) on the two blocks.
For a positive column \(i=0,\ldots,a-1\), the number of letters
is \(\sum_{h>i}c_h^+=\sigma_i-(a-i)-(b'-w_j)\).
This telescopes the counts in Theorem~\ref{thm:gluedcost811}
to the truncated endpoint at \(h=a\). In a negative column
\(i=s-1-l\), \(0\le l<b\), the count is
\(\sum_{h>l}c_h^-=-w_j-b'-\sigma_i-b+l\), by the other endpoint
calculation of Theorem~\ref{thm:firsthol817}. The resulting two
compact weights, before conjugating the descent, are
\[
 -\sigma_i-i+(a-b-1)/2\quad(i<a),\qquad
 -\sigma_i+l+(a-b+1)/2\quad(i=s-1-l).
\]
These are the lowest compact type of the physical contragredient
discrete series, in its highest-line order, as computed in
Theorem~\ref{thm:firsthol817}. The case \(b=0\) uses only the
positive count.

For completeness, compact dual-pair decomposition identifies the
multiplicity space of the matching \(K'\)-type with an irreducible
real \(G\)-module, and the joint harmonic belongs to its lowest
compact type. This is the harmonic decomposition of
\href{https://msp.org/pjm/1987/129-1/pjm-v129-n1-p02-s.pdf}
{Davidson, \S2, Theorems~2.4, 2.5 and 2.7}, tensored by the
splitting characters. The coefficient of this vector is in
\(L^2\) by Lemma~\ref{lem:oscillatorcartan710}. Its pairing with
the indicated discrete-series coefficient is positive by
Theorem~\ref{thm:realgram710}, now applied to the compact weight
just calculated. Schur orthogonality therefore identifies the
irreducible module with that discrete series (with the usual
finite-dimensional orthogonality when \(b=0\)). Its highest
line in the lowest compact type is one dimensional. An
equivariant theta map on its algebraic vectors consequently
puts the output on that line in every nonzero copy. Applying
this observation to \(D\) and conjugating puts \(d_\Xi\) on
the prescribed source line. The cuspidal decay argument of
Theorem~\ref{thm:restrictedtower848} gives square integrability.
This proves \(d_\Xi\in\mathcal H_j\).

For \(f\in\mathcal H_j\), insert the bilinear kernel of the
single lift into \(\langle\Theta^{\rm cl}_{j,\vartheta} f,\Xi\rangle\).
On a definite class, integration of its highest coefficient and
the conjugate coefficient over the probability measure on \(K'\)
gives the vector inner product divided by \(\dim V_\rho\), by
compact Schur orthogonality. Disintegration at the finite level
gives the class weight \(1/|\Gamma|\) times the level volume.
Thus replacing the class pairing by the quotient integral
multiplies by a nonzero constant independent of \(f,\Xi\).
The two quotient integrations may be interchanged since the
definite one is compact, the kernel has moderate growth there
uniformly in that variable, and \(f\) is rapidly decreasing.
The result is this constant times
\(\int_{[G]}\overline{f(g)}\,d_\Xi(g)\,dg\), including any
fixed normalization of the lift. If \(\Xi\perp Z\), take
\(f=d_\Xi\). Its norm is zero and continuity gives pointwise
vanishing.
\end{proof}

\begin{theorem}[The restricted tower of every false lift]
\leavevmode\label{thm:falsedescent849}
Let \(w\) be a place tested in Proposition~\ref{prop:testedkernel785},
and take \(\Xi\in e_{\rm mark}V\cap\ker A_w\cap Z^\perp\).
Then for every \(\zeta\in V_\rho\), every \(0\le h\le s+9m+1\),
and all \(\phi\in\mathcal S_S(W_j\oplus\mathbb H^h)\), one has
\[
                    D_{W_j\oplus\mathbb H^h}(\phi,y_\zeta)=0 .
\]
\end{theorem}
\begin{proof}
The proposition gives zero for the block and highest harmonic
at every argument. Fix the finite block for the moment. Averaging
at \(K'\) in the definite integral kills all real oscillator
types except the one matching \(y_{\zeta_0}\). In that isotypic
space the integral at each \(g\) contracts the compact factor
against the coordinate of \(\zeta_0\). Indeed covariance under
simultaneous translation and compact Schur orthogonality give
the invariant pairing of the two compact modules, and that
pairing has multiplicity one. The chosen joint harmonic has
nonzero matched compact coordinate by construction: the mixed
flag-minor map sends the highest vector to the product of its
disjoint positive and negative flag minors, which has value
one at the corresponding identity frames.

In the remaining factor of the oscillator, the compact source
translates and real enveloping algebra applied to the harmonic
vector span all compact-finite vectors of that irreducible
multiplicity module. This uses algebraic cyclicity of the
Harish--Chandra module of an irreducible unitary admissible
representation. Every polynomial Fock vector is compact-finite;
projection to a fixed compact coordinate preserves this
property. The argument applies factor by factor at infinity.
The theta integral intertwines real source translations with
right translations on \(G\); real differentials may be passed
under the compact definite integral by continuity of the
Schwartz derivatives. Differentiating or translating the zero
function from the proposition gives zero on every one of these
polynomial vectors. Finite Hermite expansions are dense for the
Schwartz topology, so vanishing extends to all real Schwartz
vectors, still with the same finite block.

Apply Theorem~\ref{thm:markedbottom847} separately to each real
vector. This gives zero for every datum in \(\mathcal S_S(W_j)\)
for the coefficient \(y_{\zeta_0}\). For \(k'\in K'\) one has
\(y_{\rho(k')\zeta_0}(h)=y_{\zeta_0}(hk')\). Changing variables
in the integral moves this compact translation to the real
Schwartz datum. The translates of \(\zeta_0\) span \(V_\rho\)
by irreducibility. Thus the same bottom vanishing holds for
every coefficient \(y_\zeta\). Restoring or removing an
external character of the small argument multiplies each
function by a nowhere zero character.
Finally the map \(\Xi\mapsto y_\zeta\) is linear in \(\Xi\) and
commutes with the right finite action. Thus \(A_w y_\zeta=0\)
at precisely the same \(w\). Theorem~\ref{thm:restrictedtower848}
now proves the assertion at all the required ranks.
\end{proof}
\clearpage

\subsection{Fourier coefficients and orthogonality of the comparison}
\label{sec:testedcomparison850}
Use the rational orthogonal decomposition \(I_n=I_p\oplus I_N\)
of Section~\ref{sec:integralgluing811}, with \(p=s+9m+1\)
and \(N=t-9m-1\). Write \(H=U(I_n)\), \(H_1=U(I_p)\),
\(H_2=U(I_N)\). Normalize Haar measures consistently under
disintegration; their positive factors will not affect
vanishing. Put \(e_p=\binom{I_p}{0}\), a frame of Gram matrix
\(I_p\). The coherent splitting for \(\mathbb H^p\) is the
one in which \(H\) acts by substitution on the \(p\) positions
in \(I_n\). Indeed each hyperbolic plane contributes the ordinary
position representation with unit absolute Jacobian on \(H\);
the external power \(\upsilon_1^p\) is exactly the one in
Theorem~\ref{thm:descenttower846}.

\begin{lemma}[Detecting a complementary period with the fixed vacuum]
\leavevmode\label{lem:frameperiod850}
Let \(Y\) be a continuous function on \([H]\), right invariant
by the lattice stabilizers of \(H\) outside \(S\) at finite
places. If its integrals against the theta kernels for
\(\mathbb H^p\) vanish for all small-group arguments and data
in \(\mathcal S_S(\mathbb H^p)\), then
\[
          \int_{[H_2]}Y(\operatorname{diag}(h_1,h_2))\,dh_2=0
                   \qquad(h_1\in H_1(\mathbb A_F)).
\]
The integrals here pair the kernel with \(Y\) by multiplication.
\end{lemma}
\begin{proof}
Set \(g_0=\operatorname{diag}(h_1,1)\) and
\(J(h)=\int_{[H_2]}Y(\operatorname{diag}(1,h_2)h)\,dh_2\).
This continuous function descends to
\(H_2(\mathbb A_F)\backslash H(\mathbb A_F)\) by right
invariance of quotient measure. The map on this homogeneous
space given by \(h\mapsto h^{-1}e_p\) is the map to isometric
frames. At each local place its stabilizer is exactly \(H_2\),
and its image is the whole Gram-\(I_p\) orbit by Witt extension.

Take the Fourier coefficient for the character of \(I_p\) on
the Siegel radical of \(U(\mathbb H^p)\), at an element of
its \(\mathrm{GL}_p(\mathbb A_E)\) Levi. In position coordinates
the radical multiplies the term of \(X\in I_n^p(F)\) by the
additive character paired with \(X^*X\). This follows from
the Weil action of the upper unipotent in the polarization
by the two isotropic halves. Any fixed trace-pairing scalar
is used both in this character and in taking the coefficient.
Orthogonality on the additive quotient therefore retains
\(X^*X=I_p\). Witt extension over \(E/F\) makes these rational
frames one \(H(F)\)-orbit, with stabilizer \(H_2(F)\).
Unfolding the sum over that orbit and integrating first over
\([H_2]\) gives, up to a nonzero Levi scalar,
\[
       0=\int_{H_2(\mathbb A_F)\backslash H(\mathbb A_F)}
                         \phi(h^{-1}e_p c)\,J(h)\,d\dot h
             \qquad (c\in\mathrm{GL}_p(\mathbb A_E)).
\]
Here \(\phi\) is the position function in the hyperbolic model,
and the Levi action pulls its argument back by right multiplication
by \(c\), after choosing the corresponding Levi element. The
unitary \(H\)-action has no remaining factor in this splitting.
Absolute convergence before unfolding follows from uniform
Schwartz summation on the compact definite and radical quotients,
with \(c\) fixed. This is the unitary frame computation;
the corresponding orthogonal calculation is in
\href{https://arxiv.org/pdf/1501.01565}{Walls, Proposition~2.1.1}.

Away from \(S\) at finite places the position vacuum is a
nonzero scalar times the indicator of
\(\alpha_v\operatorname{Mat}_{n,p}(\mathcal O_{E_v})\), with
the finite scale of Lemma~\ref{lem:sparsedistance814}. Choose
\(c_v=\alpha_v h_{1,v}\) there. The support condition becomes
integrality of \(x_v=h_v^{-1}e_p h_{1,v}\), whose Gram matrix
is \(I_p\). These integral frames are a single orbit of
\(K_v=H(\mathcal O_v)\). Indeed the span of the columns splits
orthogonally in the integral lattice by the retraction \(x_v^*\).
The complement is unimodular. At an unramified field place
it has an orthonormal basis by the unit norm argument of
Section~\ref{sec:commonsource806}. At a split place choose
a basis on the first component and its dual on the other.
This proves that \(x_v=k_v^{-1}e_p\) for some \(k_v\in K_v\).
Equivalently the contributing class belongs to
\(H_2(F_v)g_{0,v}K_v\). On this single orbit \(J\) cannot
distinguish right translations by \(K_v\).

At infinity and at \(S\) take \(c_v=h_{1,v}\) and choose smooth
nonnegative bumps in the real position spaces and compact-open
indicators in the finite ones, positive at \(e_p\) and supported
in arbitrarily small neighborhoods of \(e_p\). Restricted to
the frame orbit these localize the coset at \(g_{0,v}\).
In detail, the local frame map has a continuous local section.
At a real place complete nearby orthonormal frames by projecting
the old complementary basis and orthonormalizing it. At a
marked split place the frame is a pair of matrices \(A,B\) with
\(B^{\mathsf t}A=I_p\). Complete \(A\) by applying
\(I-AB^{\mathsf t}\) to the old complementary basis, and use
the inverse transpose of the completed matrix on the other
component. Its determinant remains a unit near the standard
frame. These sections also show that the chosen bumps are
positive on a relatively open neighborhood of the desired
coset. The product of the unmarked compact orbits is compact
and open in the unmarked restricted product, and has positive
finite quotient measure. The function \(J\) is invariant on
the right by their full product of compact groups. Thus if
\(J(g_0)\ne0\), one constant complex rotation makes its real
part strictly positive throughout the support of sufficiently
small bumps on the quotient. The factors outside these places
are positive indicators times a common nonzero scalar; remove
that scalar in the integral. The integral must then have positive
real part, a contradiction. This proves the lemma.
\end{proof}

\begin{proposition}[The period with a theta kernel on the second block]
\leavevmode\label{prop:kernelperiod850}
For a false lift \(\Xi\) as in Theorem~\ref{thm:falsedescent849},
every coefficient \(y_\zeta\), and every finite
\(\gamma\in H(\mathbb A_{F,f})\) supported on \(S\), one has
\[
 \int_{[H_2]} \theta^\circ(g,h_2;\phi_-)\,
      \overline{y_\zeta(\operatorname{diag}(h_1,h_2)\gamma)}\,dh_2=0
\]
for all \(h_1,g\), with any allowed second-block test \(\phi_-\)
and the characters of Section~\ref{sec:complementperiod812}.
\end{proposition}
\begin{proof}
Replacing \(y_\zeta(h)\) by \(y_\zeta(h\gamma)\) in the tower
integrals moves \(\gamma^{-1}\) to the test by change of variables.
Its support only in \(S\) preserves the allowed space. Fix a
bottom test \(\phi_+\otimes\phi_-\) for the splitting of the
target into the two blocks; form the function
\[
        Y(h)=\theta_{W_j,I_n}(g,h;\phi_+\otimes\phi_-)\,
                                      \overline{y_\zeta(h\gamma)} .
\]
It is continuous, rationally invariant and fixed by the unmarked
finite compact group. The two blocks use the same source-space
splitting on their determinants. Their external characters on
the \(W_j\) argument may be taken inverse to one another,
with the prescribed character on the second block. On
\(\mathbb H^p\oplus W_j\) use the coherent splitting, and omit
any overall character depending only on the small-group argument.
For a tensor of the hyperbolic test and the bottom test its
theta kernel restricted to \(U(\mathbb H^p)\times U(W_j)\)
is the product of the two theta kernels. This follows because
the oscillator of an orthogonal sum is a tensor product, and the
rational positions in a product polarization sum independently.
All these tensors have the fixed vacuum away from the indicated
places. Theorem~\ref{thm:falsedescent849} at \(h=p\) therefore
gives the hyperbolic vanishing of Lemma~\ref{lem:frameperiod850}
for \(Y\).

At the block diagonal point the bottom theta kernel factors again,
now in the two target spaces. Its first factor can be made nonzero
at any fixed \((g,h_1)\) by choosing \(\phi_+\). To see this,
choose a rational polarization of the underlying symplectic
space and any nonzero allowed finite tensor. After action by
\((g_f,h_{1,f})\) it is a nonzero Schwartz--Bruhat function.
Some rational position lies where this function is nonzero,
since a vector space over \(F\) is dense in its finite adeles.
Rational positions in its compact support belong to a single
fractional lattice, discrete at infinity. Choose a real smooth
bump supported around the selected position and excluding all
the others, with value one at that point. Apply the inverse of
the real oscillator action to obtain the real test. The resulting
theta value has a single nonzero term. With this \(\phi_+\),
divide the period supplied by the lemma by its nonzero first
factor. This gives the required second-block integral.
\end{proof}

\begin{theorem}[The comparison of true lifts belongs to the tested space]
\leavevmode\label{thm:testedcomparison850}
With \(Z=\Theta_{j,\vartheta}^{\rm cl}\mathcal H_j\) and \(T\)
the simultaneous orthogonal complement of the false spaces of
Theorem~\ref{thm:projectionindex784}, the operator of
Theorem~\ref{thm:integralcomparison845} satisfies
\[
                        A_\vartheta Z\ \subseteq\ T .
\]
In fact \(G_\vartheta(P_j\mathcal H_j)\subseteq T\).
\end{theorem}
\begin{proof}
At each complementary class \(y\), let \(K(g,y)\) be the
Riesz vector in the negative coefficient space, with
\(\langle v,K(g,y)\rangle=\theta^\circ_{\bar v}(g,y)\) from
the formula in Theorem~\ref{thm:gluedperiodmap844}. Thus at
an infinite argument \(u\in H_2(F_\infty)\) the same coordinate
is \(\langle r(u)v,K(g,y)\rangle\), where \(r\) is the
negative coefficient representation. Indeed the compact
oscillator substitutes \(u\) in its negative row variables;
the flag-minor map on the conjugated coefficient gives
\(\overline{r(u)v}\). This is the contragredient block in
Lemma~\ref{lem:dualcartan811}; the vacuum and source determinant
on \(H_2\) have opposite weights, as for the original single
lift. For a fixed positive coefficient \(\xi\) the Cartan map
\(C_\xi v=\xi\odot v\) intertwines \(r(u)\) with the restriction
of \(\rho(\operatorname{diag}(1,u))\), by its tensor construction.

Fix \(i\) and a first-block finite representative \(h\).
Apply Proposition~\ref{prop:kernelperiod850} with \(h_1=(1,h)\),
\(\gamma=\gamma_i\), the second-block test of the gluing formula
with coefficient \(\bar v\), and the coefficient \(y_{C_\xi v}\).
Take the conjugate of its zero integral and sum over an
orthonormal basis \(v\). At \(h_2=(u,y)\) this gives
\[
  \sum_v \overline{\langle r(u)v,K(g,y)\rangle}\,
          \langle C_\xi r(u)v,\Xi(z_{i,h,y})\rangle
          =\langle C_\xi K(g,y),\Xi(z_{i,h,y})\rangle .
\]
The equality is expansion in the orthonormal basis \(r(u)v\).
The resulting integrand no longer depends on \(u\). It descends
to the complementary finite level \(J_H\): at the old marks
\(J_H\equiv1\bmod\pi^2\) and conjugation by \(\gamma_i\) costs
at most one power of \(\pi\), since
\(\pi L\subseteq\gamma_i L\subseteq L\) in
Theorem~\ref{thm:splitlattice811}. At a new mark the second
block acts trivially on the top determinant of the level, and
at other places the levels are the lattice stabilizers.
The kernel itself has at least this complementary level. By
Lemma~\ref{lem:gluedrow844} all complementary class stabilizers
are trivial. Disintegrating the period therefore gives the
unweighted sum of the displayed pairing over \(y\), times
one positive level and compact volume. The sum is zero.

Set \(\xi=\mathcal Y_{i,\vartheta}(h)\), sum over \(h\), and
use the definition of \(G_\vartheta\) and the class evaluation
identity preceding Theorem~\ref{thm:gluedperiodmap844}. It gives
\[
          \sum_y\langle K(g,y),(G_\vartheta^*\Xi)_{i,y}\rangle=0
                             \qquad(i,g).
\]
The integral formula for \(\langle\zeta,P_jF\rangle\) in that
theorem now gives zero for \(\zeta=G_\vartheta^*\Xi\): its inner
sum over \(y\) is the conjugate of the last display. Thus
\(\langle\Xi,G_\vartheta P_j F\rangle=0\) for every good \(F\)
and every false lift at each tested place. This proves the
second assertion. Since
\(b_{j,\vartheta}G_\vartheta^*\Theta_{j,\vartheta}^{\rm cl}
  =P_j\) on the good space and \(b_{j,\vartheta}\ne0\),
multiplication by \(B_0h_j^-G_\vartheta\) proves the first.
\end{proof}
\clearpage

\subsection{The projected determinants on the common period interval}
\label{sec:compressionbudget851}
Fix \(q_0,d_{\rm K}\), an interval starting at \(V_0\) in the
shorter family, and the simultaneous twists of
Corollaries~\ref{cor:actualdenom803} and~\ref{cor:detheight810}.
At a node \(x\), write \(b_x\) for its weight scale, \(s=k+m-1\),
\(f=[E^+:\mathbb Q]\) and
\[
        B_x^{\rm det}=\frac{fxb_x(\log b_x)^3}{mk}.
\]
For a single lift at index \(j\) write \(\beta\) for its prescribed
twist, \(Z_\beta\) for its true range and \(P_\beta\) for orthogonal
projection onto that range. Distinguish this projection from
the injective map \(P_j\) into the tuple space. Use the comparison
operator \(A_\beta\) of Theorem~\ref{thm:integralcomparison845}
and set
\[
 C_j(\beta)=P_\beta A_\beta|_{Z_\beta},\qquad
 \Delta_j(\beta)=\det C_j(\beta),\qquad
 D_x=\dim Z_\beta .
\]
The dimension is common to the indices and twists by
Theorems~\ref{thm:commonsource806} and~\ref{thm:gluedperiodmap844}.
In dimension zero take determinants equal to one.

\begin{theorem}[Finite and archimedean estimates for the same compression]
\leavevmode\label{thm:compressionbudget851}
Each \(\Delta_j(\beta)\) is positive and algebraic. With
\(v_\ell(\ell)=1\), the simultaneous local estimates are
\[
 \sum_\ell(\log\ell)\sup_{x,j,\iota}
 \frac{\max(0,-v_\ell(\iota\Delta_j(\beta_{\ell,\iota})))}
      {\max(1,D_x)\,sB_x^{\rm det}}
     \ll_{q_0,d_{\rm K}}(\log(2V_0))^{-2}.
\]
Here \(\beta_{\ell,\iota}\) is the twist at index \(j\) of
the family chosen for this coefficient embedding. At every
complex embedding,
\[
      \log\max(1,|\sigma\Delta_j(1)|)
          \ll_{q_0,d_{\rm K}} D_x f m b_x\log b_x .
\]
\end{theorem}
\begin{proof}
Theorem~\ref{thm:integralcomparison845} gives a positive
definite pullback by the single lift from \(\mathcal H_j\);
the lift is injective and onto \(Z_\beta\) by definition of
the range. Thus the compression is positive definite there.
At a coefficient embedding with prime \(\ell\) outside the
new conductor, the same operator preserves the ambient
integral lattice \(L\). Theorem~\ref{thm:testedcomparison850}
gives \(A_\beta Z_\beta\subseteq T\). Consequently
\[
           C_j(\beta)(L\cap Z_\beta)\ \subseteq\
                         P_\beta(L\cap T).
\]
Over a finite local field of definition take elementary
divisors \(\pi^{-a_i}\) for the larger lattice on the right
relative to \(L\cap Z_\beta\), with \(\pi\) a uniformizer.
The top exterior power of the inclusion shows
\(v_\ell(\iota\Delta_j(\beta))\ge-\sum_i a_i/v_\pi(\ell)\).
The quotient on the right is the normalized length
\(\lambda\) of Theorem~\ref{thm:projectionindex784}.
The twists attached to \(\ell\) have new conductor prime to
\(\ell\). Inserting the bound for the actual \(\lambda\)
in Corollary~\ref{cor:actualdenom803}, which takes the same
supremum over nodes, indices and coefficient embeddings, proves
the first estimate.

By Proposition~\ref{prop:archcontraction804}, on the transported
tested space the transported \(P_\beta\) is an orthogonal
projection in its positive class metric. Since \(A_\beta\)
takes the true range into the tested space before and after
transport, the norm of the transported compression is at most
\(\|\sigma(A_\beta)\|_{\rm op}\). The determinant in dimension
\(D_x\) has absolute value at most the \(D_x\)-th power of
this norm by the top exterior power. For \(\beta=1\) apply the
estimate of Theorem~\ref{thm:integralcomparison845}. For an
arbitrary fixed twist the ambient class tables, their metrics
and the entries of \(A_\beta\) are algebraic in fixed frames.
Their finitely many conjugates give a finite bound depending
on that twist. Thus the orbit of \(\Delta_j(\beta)\) under all
automorphisms of \(\mathbb C\) is bounded. A transcendental
number \(z\) has unbounded orbit. To see this, extend \(z\) to a
transcendence basis, replace it by \(az\) for positive integers
\(a\), and extend the resulting rational function field
automorphisms to \(\mathbb C\). This proves algebraicity also
when the original analytic source basis contains periods.
\end{proof}

Use the notation of Theorem~\ref{thm:thirdglobal843}. At index
\(j\) and \(c=\chi_M^j\beta\) put
\[
  \mathcal L_j(\beta)=
    \mathcal D_{m/2}(c\upsilon_1^{m/2})\,
    \mathcal D_m(c^{-1}\upsilon_1^m)\,
    \mathcal D_{e'}(\sigma_c),\qquad
       e'=\frac{k+7m}{2},\qquad
       \sigma_c=\alpha^{-1}\upsilon_1^{\,3m+e'} .
\]
Indeed the two ranks in the third exchange are \(3m\) and
\(p=s+9m+1\), giving \(e'=(p-3m)/2\). Each \(\mathcal D\)
includes its own conductor power and its squared absolute finite
\(L\)-values as in Section~\ref{sec:secondglobal833}.
Let \(\mathcal V_j(\beta)\) denote the Gram determinant of the
single lift with the class metric of Theorem~\ref{thm:auxdet805}.

\begin{proposition}[Exact compression identity for a single lift]
\leavevmode\label{prop:compressionidentity851}
On the fixed node one has
\[
 \frac{\mathcal V_j(1)}{\mathcal V_j(\beta)}
   =\frac{\Delta_j(\beta)}{\Delta_j(1)}
        \left(\frac{\mathcal L_j(1)}{\mathcal L_j(\beta)}\right)^{D_x}.
\]
\end{proposition}
\begin{proof}
Represent the lift by an invertible matrix \(U\) between
orthonormal bases of the good source and its true image.
Then the determinant of its pullback of \(A_\beta\) is
\(|\det U|^2\Delta_j(\beta)=\mathcal V_j(\beta)\Delta_j(\beta)\).
By Theorem~\ref{thm:integralcomparison845} that same determinant
is \((B_0h_j^-)^{D_x}|b_{j,\beta}|^{-2D_x}\det(P_j^*P_j)\),
which is nonzero. The map \(P_j\) and both integer factors are
common to the twists. Theorem~\ref{thm:thirdglobal843} gives
\(|b_{j,\beta}|^{-2}=\text{const.}\,\mathcal L_j(\beta)\)
with the same positive constant for \(\beta\) and \(1\). Division
gives the identity.
\end{proof}

\begin{corollary}[A common multiplier for the terminal determinant]
\leavevmode\label{cor:terminalcomparison851}
Write \(\mathfrak Z_\vartheta\) for the terminal class Gram
determinant of Theorem~\ref{thm:detquotient810}. At index \(j\)
use \(\beta=\vartheta^j\) and put
\[
 y_{\ell,\iota,x}
  =\frac{\mathfrak Z_1}{\mathfrak Z_\vartheta}
       \prod_{j=1}^m
        \left(\frac{\mathcal L_j(\vartheta^j)}{\mathcal L_j(1)}
        \right)^{D_x},\qquad
                    H_x=\prod_{j=1}^m\Delta_j(1)>0 .
\]
These \(y_{\ell,\iota,x}\) and \(H_x\) are algebraic. The multiplier
is independent of \(\ell,\iota\). Uniformly over the same nodes,
\[
 \begin{split}
 \log\max(1,|\sigma H_x|)
      &\ll_{q_0,d_{\rm K}}D_x f m^2 b_x\log b_x,\\
 \sum_\ell \sup_{\iota}
              \log\max(1,|\iota(H_x y_{\ell,\iota,x})|_\ell)
      &\le \max(1,D_x)\delta+
         O_{q_0,d_{\rm K}}\!\left(
             \frac{D_x m s B_x^{\rm det}}{(\log(2V_0))^2}\right).
 \end{split}
\]
Here \(\delta>0\) is the tolerance of
Corollary~\ref{cor:detheight810}. The two bounds can in particular
be added with the same order of error as in the second bound.
\end{corollary}
\begin{proof}
Let \(R_{\ell,\iota,x}\) be the left side of
\eqref{eq:relchain810}. That identity uses the product of the
\(m\) auxiliary Gram determinants. Proposition~\ref{prop:compressionidentity851}
at each index consequently gives the exact equality
\[
         H_x y_{\ell,\iota,x}
              =R_{\ell,\iota,x}^{-1}\prod_{j=1}^m\Delta_j(\vartheta^j).
\]
All factors on the right are nonzero and algebraic by
Corollary~\ref{cor:detheight810} and the theorem. The first
bound follows by adding the \(m\) archimedean bounds. For the
second, apply the positive part of the negative valuation to
the product. The inverse of \(R\) costs its positive valuation
times \(\log\ell\); the sum of those terms is at most
\(\max(1,D_x)\delta\) on the chosen interval. For each
compression the theorem bounds the corresponding sum by
\(O(D_xsB_x^{\rm det}/(\log(2V_0))^2)\) when \(D_x>0\).
In dimension zero all the factors are one. Finally,
\(f m b_x\log b_x/(sB_x^{\rm det})
=m^2k/(sx(\log b_x)^2)\le4/(\log b_x)^2\).
Since \(b_x\ge x^{400}\), this absorbs the archimedean cost.
\end{proof}

The prime-sum argument uses the ratio
\(\mathfrak Z_1/\mathfrak Z_\vartheta\) itself. Its remaining
factor in the last corollary is the inverse of the displayed
product of the \(\mathcal L_j\) ratios. In particular
\(7m/2< z\le e'\) in the third factor contributes \(k/2\)
critical integers at each \(j\); the direction-dependent
divisibility estimate for these factors enters before one
can deduce the required prime-sum inequality.
\clearpage

\section{Critical torus periods and their finite tests}\label{nav:critical-torus-tests}

\subsection{The norm of the torus lift at a critical integer}
\label{sec:torusnorm852}
Let \(E/F\) be a field of the shorter family, \(T=U(1)\), and
\(n=2z\) with \(z\ge2\). Let \(\nu\) be a unitary character of
\([T]\) with weight \(-c_\tau\) at each real place and put
\(d_\tau=c_\tau-z\ge0\). Its finite part is unramified off the
original marks \(S_0\). Let \(\beta\) be a prescribed weight-zero
tame twist with new split marks \(S_\beta\), disjoint from \(S_0\).
Take both sets as sets of places of \(F\), distinguishing a member
of the pair in \(E\) above each one. Write
\[
        Q_\beta=\prod_{v\in S_\beta}q_v,\qquad
        M_0=\prod_{v\in S_0}(q_v-1),\qquad
        \lambda_\eta(a)=\eta(a/\bar a).
\]
A superscript \(S\) on \(L_{E,f}\) omits the finite places above
\(S\). Absorb \(\nu\beta\) in the theta kernel from \(T\) to
\(H=U(I_n)\), with external source factor
\(((\nu\beta)\upsilon_1^z)^{-1}\) and trivial external target
factor. At infinity use the one-row highest polynomial of degree
\(d_\tau\), with a unit highest vector in the coefficient module
\(V_{\mathbf d}\). At the old marks use the indicator of a
specified nonzero residue column on the integral box. At a new
mark use the tame character on the specified entry of the
column, extended by zero, on that box. At other finite places
use the lattice vacuum. Let \(K_\beta(g)\) be the resulting
vector-valued kernel at a source class \(g\). Use the class
group \(C\) and the target class metric of
Section~\ref{sec:rankonepair736}. In particular
\[
        \phi_\beta=S_{\nu,z}\sum_{g\in C}K_\beta(g)
\]
with the same nonzero scaling \(S_{\nu,z}\) before and after
twisting. This convention allows the fixed vacuum and derivative
scales in an integral normalization. The measure on each compact
real group is probability Haar and the original maximal compacts
at finite places have mass one.

\begin{theorem}[Norm and conductor for the torus period]
\leavevmode\label{thm:torusnorm852}
In the Fock metric with monomial squared norm \(\alpha!\), there
is a positive constant \(a_{E,n}\) depending only on the adelic
spaces, the oscillator normalization and the measures such that
\[
\begin{split}
 \|\phi_\beta\|_{\rm cl}^2
   ={}&a_{E,n}\,|S_{\nu,z}|^2 M_0\,
       (\dim V_{\mathbf d})\prod_{\tau}d_\tau!\,
       Q_\beta^{\,n-1}\,
       \frac{L_{E,f}^{S_0\cup S_\beta}
                        (z,\lambda_{\nu\beta}^{-1})}{\zeta_F(n)} .
\end{split}
\]
In particular all these lifts are nonzero. If \(R_{\nu,z}(\beta)\)
is the ratio of the absolute values of the displayed numerator
and its value at \(\beta=1\), with their respective omissions,
then
\[
       R_{\nu,z}(\beta)
           =\frac{\|\phi_\beta\|_{\rm cl}^2}
                   {Q_\beta^{\,n-1}\|\phi_1\|_{\rm cl}^2}.
\]
\end{theorem}
\begin{proof}
First integrate the constant function on \([T]\) with Haar
measure. Apply the convergent unitary Siegel--Weil formula
to the double of the source line. The doubled group has
Siegel rank one, the definite target has Witt index zero and
\(n>2\). The section is evaluated at \((n-1)/2>1/2\), in the
region of absolute convergence. These are the hypotheses of
\href{https://msp.org/pjm/2013/264-1/pjm-v264-n1-p11-s.pdf}
{Yamana, Theorem~2}. For this double every rational isotropic
line is a graph of an isometry of the source line. Indeed a
nonzero isotropic vector in the double has both coordinates
nonzero and their norms agree. In the graph term the doubled
Schwartz function at zero, in difference and Fourier
coordinates, is the inner product of the original test and
its translate by that isometry. Unfolding the sum over the
rational isometries in the two torus integrals gives
\(\operatorname{vol}([T])\int_{T(\mathbb A_F)}
 \langle\Phi,\omega(u)\Phi\rangle\,du\), times the positive
Siegel--Weil and measure scalar. This is the graph calculation
of Proposition~\ref{prop:doubledkernel711}, now with the
constant source. The absolute local sum at a good split prime
is bounded by \(1+2q^{-z}/(1-q^{-z})\), as calculated below.
At a good inert prime it is one. Their product converges
absolutely for \(z>1\), which justifies unfolding and
multiplying the integrals.

At a real place the circle acts trivially on the polynomial
times the vacuum after inserting the specified character.
In the positive one-column model each of its \(d_\tau\)
letters has source weight \(-1\), the external factor has
weight \(c_\tau-z=d_\tau\), and the positive vacuum has source
weight zero in this convention. Its integral is therefore its
squared Fock norm. For the highest monomial this is \(d_\tau!\),
by applying the derivative of order \(d_\tau\) to the monomial
at zero. In passing from this one scalar coefficient to the
vector class metric, compact Schur orthogonality multiplies
by \(\dim V_{\mathbf d}\). In fact the compact average of the
projector onto the unit highest vector is
\(I/\dim V_{\mathbf d}\) by invariance and trace.

At a good split \(v\) put \(\xi=\lambda_{\nu\beta}
(\varpi_w)^{-1}\), where \(\varpi_w\) is on the chosen
component of \(E\). For a source scalar of valuation \(h\)
the change of position variable is \(X\mapsto X\varpi^h\)
on \(F_v^n\), with half-Jacobian \(q^{-nh/2}\).
The integral of the product of the two lattice indicators
has measure \(1\) for \(h\ge0\) and \(q^{nh}\) for \(h<0\).
The unitary multiplier is \(\xi^h\). Indeed for a single
target coordinate the intrinsic split multiplier at
\(\varpi\) is the splitting character on the distinguished
component, as in the calculation of
Theorem~\ref{thm:splitzeta827}. Its \(n=2z\) powers cancel
the external \(\upsilon_1^{-z}\), since on
\((a,a^{-1})\) the restriction of the splitting character
is the square of its value on \((a,1)\) by conjugate duality.
The remaining factor is \((\nu\beta)^{-1}\). Reversing the
chosen component exchanges \(\xi\) with \(\xi^{-1}\).
The support of each valuation shell has multiplicative mass
one. Summing the coefficients gives
\[
      \sum_{h\in\mathbb Z}\xi^h q^{-z|h|}
          =\frac{1-q^{-2z}}{(1-\xi q^{-z})(1-\xi^{-1}q^{-z})}.
\]
This is the factor of \(L_{E,f}(z,\lambda_{\nu\beta}^{-1})
/\zeta_F(2z)\). At an inert prime the entire local torus is
compact and fixes the vacuum. A uniformizer can be taken in
\(F_v\), so its base-change character is one and its norm
in \(E\) is \(q^2\). The same Euler quotient equals one there.
These cases exhaust the unmarked places by
Lemma~\ref{lem:signkernel823}.

At an old place the point column and its scalar translate
overlap only for a source unit congruent to one. The test
has squared norm \(q^{-n}\) and that subgroup has measure
\(1/(q-1)\). At a new place overlap requires a unit, and the
tame factor of its action on the entry is cancelled by the
source character. The squared norm is \(1-q^{-1}\).
Conversion to the target class measure multiplies these two
numbers of integration by \(q^n-1\) and
\((q^n-1)/(q-1)\), respectively. The first is the number of
nonzero columns, the second the number of residue hyperplanes,
as in Lemma~\ref{lem:auxmass805}. Thus the actual factors are
\((1-q^{-n})/(q-1)\) and \(q^{n-1}(1-q^{-n})\).
The factors \(1-q^{-n}\) at these places complete the
reciprocal of \(\zeta_F(n)\).

Finally the Haar integral on the source equals \(1/M_0\)
times the class sum. Indeed the source level has that
volume, its real group has mass one, its stabilizers are
trivial by Proposition~\ref{prop:torusorder736}, and the new
places require only the full unit group on the source. The
passage to \(\phi_\beta\) multiplies the squared norm by
\(|S_{\nu,z}|^2M_0^2\). The remaining measure constant is
independent of the character and the weight. At a split place
the two numerator Euler factors have positive product, the
inverse of a squared modulus. At an inert place the factor
is positive as well. Absolute convergence gives a positive
global value. Taking the quotient of the two norm identities
now proves the assertions.
\end{proof}

\begin{corollary}[Two tests for the inverse critical ratio]
\leavevmode\label{cor:torusvaluation852}
At a coefficient embedding \(\iota\) with \(\ell\nmid Q_\beta\),
use \(v_\iota(\ell)=1\) and the line projection \(P_\beta\) of
Section~\ref{sec:rankonepair736}. Put \(c(a)=\bar a\).
For two integral tested vectors \(U,V\) whose pairings with
\(\phi_\beta\) are nonzero, one has
\[
\begin{split}
 -v_\iota(R_{\nu,z}(\beta)^{-1})
   ={}&-v_\iota(\|\phi_1\|_{\rm cl}^2)
       -v_\iota(\langle P_\beta U,P_\beta V\rangle)\\
     &+v_\iota(\langle U,\phi_\beta\rangle)
       +v_{\iota\circ c}(\langle V,\phi_\beta\rangle).
\end{split}
\]
For a main pair, the positive part of the second term has the
uniform projected-period bound in Corollary~\ref{cor:actualdenom803}
on the same interval. The tests for the last two terms may
be chosen individually for \(\iota\) and \(\iota\circ c\) with
the same twist \(\beta\), provided they satisfy the respective
integrality conditions of Theorem~\ref{thm:rankonepair736}.
\end{corollary}
\begin{proof}
For a Hermitian metric linear in its second argument,
\(P_\beta v=\phi_\beta\langle\phi_\beta,v\rangle/
\|\phi_\beta\|_{\rm cl}^2\). Hence
\[
 \frac1{\|\phi_\beta\|_{\rm cl}^2}
    =\frac{\langle P_\beta U,P_\beta V\rangle}
            {\langle U,\phi_\beta\rangle\,
             \overline{\langle V,\phi_\beta\rangle}}.
\]
By the theorem \(R_{\nu,z}(\beta)^{-1}\) is this expression
times \(Q_\beta^{n-1}\|\phi_1\|_{\rm cl}^2\). The conductor
power is a unit at \(\ell\). Taking negative valuations proves
the equality, including the conjugated embedding in the last
summand. The period estimate is uniform over the allowed
tests of the given lift. A directional saving requires upper
bounds for the two last valuations, as well as the simultaneous
size of a multiplier for the untwisted norm.
\end{proof}
\clearpage

\subsection{Integral flat torus values at the original level}
\label{sec:flattorus853}
Keep the fields, old marks and measures of
Section~\ref{sec:torusnorm852}. For an even rank \(r=2a\ge4\)
put \(D=fr\), with \(f=[F:\mathbb Q]\), and \(M=p_0p_1\).
A flat pair of rank \(r\) has \(\nu=\psi\upsilon_1^{-a}\)
with \(\psi\) of weight zero and unramified off the old marks.
We take \(\psi\) tame of level at most one there. Thus its
external source character in the kernel, before inversion,
is \(\psi\beta\). The real polynomial is the constant one
and the target coefficient space is the scalar line. The
mark on the old residue column can be any nonzero column.
Use the standard unit box at all unmarked places.

In the notation of Lemma~\ref{lem:tracebasis821}, write \(t_E\)
for the integer denoted there by \(n\), chosen with its product
coordinates for the pair \(I_1,I_r\). Put
\[
 b_r=(M t_E)^D,\qquad
 B_r=|\mathrm{GL}_{2D}(\mathbb F_3)|.
\]
Write \(a_r\) for the scalar that sends the unit oscillator
vacuum to its marked elliptic normalization at degree zero.
Precisely, it is \(C_r^{-1}\) in the identification of a kernel
with the value of its marked section at the origin in
Theorem~\ref{thm:vac242}, using the product coordinates just
chosen. It is common to all finite data and characters in
this rank. Let \(\phi_{\psi,\beta}=a_r\sum_{g\in C}K_\beta(g)\).
This fixes the freedom in \(S_{\nu,a}\) of
Theorem~\ref{thm:torusnorm852} for the flat pairs. Define
\[
        J_r(\psi,\beta)=
         B_r b_r^2\,\|\phi_{\psi,\beta}\|_{\rm cl}^2/
                                      Q_\beta^{\,r-1}.
\]

\begin{theorem}[Flat coefficients with a common normalization]
\leavevmode\label{thm:flattorus853}
For every finite target argument \(h\) the value
\(b_r\phi_{\psi,\beta}(h)\) is algebraic and integral at all
places above every \(\ell\nmid Q_\beta\). The same holds for
\(J_r(\psi,\beta)\), which is positive. Moreover
\[
  J_r(\psi,\beta)/J_r(\psi,1)
       =R_{\psi\upsilon_1^{-a},a}(\beta).
\]
For \(4\le r\le 2s\) on the shorter interval one has
\[
 \sup_{\psi,\sigma}\log\max(1,|\sigma J_r(\psi,1)|)
          \ll_{q_0,d_{\rm K}} x^{70}(1+\log b_x)
          =o(B_x^{\rm det}),
\]
uniformly as the starting node tends to infinity. Here \(\sigma\)
runs through complex embeddings and \(s=k+m-1\).
\end{theorem}
\begin{proof}
All residue functions at an old place are indicators on
a group \(k^r\). Fourier inversion writes one point indicator
as a sum of additive characters with coefficients roots of
unity divided by \(q^r\). Their absolute values sum to one
also after any conjugation. At a new place the tame function
on the chosen entry, extended by zero, is an integral linear
combination of point indicators on the same group. The
Fourier coefficients thus have denominator dividing \(q^r\).
The product of these denominators at the old places is
\(M^{fr}\), and every new denominator is a unit above
\(\ell\nmid Q_\beta\). Multiplication by each additive
character on an integral position box is a translation in
the dual coordinates by a torsion point. This is finite
Fourier inversion in the split position model as in
Proposition~\ref{prop:markedfourier820}.

We explain the remaining scalar in these coordinates. Since
\(r\) is even, Lemma~\ref{lem:tracebasis821} supplies the
product period lattice \(B\) with
\(t_E B\subset L\subset t_E^{-1}B\), for the vacuum lattice
\(L\). In transferring the normalized product section to
\(L\), one sums unitary torsion translates with a factor
\(I^{-1/2}\), \(I=[L:L\cap B]\), by Lemma~\ref{lem:basis59}.
The group defining \(I\) embeds in \(t_E^{-1}B/B\). Thus
\(I\mid t_E^{2D}\) and \(t_E^D I^{-1/2}\) is an algebraic
integer. For the two finite arguments choose split
representatives off \(\ell\), the level and the discriminant
by Theorem~\ref{thm:class61}, at compact opens fixing the
given data. The change to the moved lattice \(L'\) has the
same normalized sum, now with
\([L':L\cap L']^{-1/2}\) by Lemma~\ref{lem:split61}. This is a
unit at \(\ell\) since \(L'_\ell=L_\ell\). All phases of the
torsion translations are roots of unity. The external
character \(\psi\beta\) has finite order since it is trivial
at infinity and factors through a finite torus ray class
group. The remaining intrinsic split phases are units above
\(\ell\) on these representatives by the determinant calculation
in Theorem~\ref{thm:one61}.

Each summand is now the value at the origin of a translate of
the product of marked elliptic sections. Its factors are
algebraic and integral at the prime under consideration by
the order-zero case of Theorem~\ref{thm:torsion312}. This
case applies even when the translation is of \(\ell\)-power
order, since its derivative order is zero. The passage to
these marked sections accounts exactly for \(a_r\) by
Theorem~\ref{thm:vac242}. Source real rotations act trivially
after the flat character cancellation, and the target real
representation is trivial. Hence rational changes of representatives
move no positive-degree polynomial or coefficient lattice. These
calculations give the asserted integrality and algebraicity of
\(b_r a_r K_\beta(g)(h)\) for all finite arguments. The source
sum has integer multiplicities and trivial stabilizers by
Proposition~\ref{prop:torusorder736}.

For a target class representative \(h\) its stabilizer is
a finite group acting faithfully on its underlying integer
lattice of rank \(2D\). Reduction modulo three injects this
group into \(\mathrm{GL}_{2D}(\mathbb F_3)\) by the binomial
argument in Theorem~\ref{thm:definiteclass822}. Thus its
order divides \(B_r\). The class sum for the norm of
\(b_r\phi_{\psi,\beta}\) after multiplication by \(B_r\)
is integral at every such prime, also for the conjugated
coefficients, since the coefficient integrality holds at
all primes over \(\ell\). The power of \(Q_\beta\) is a unit
there. Positivity and the ratio identity follow from
Theorem~\ref{thm:torusnorm852} with this single choice of \(a_r\).

For the size bound we only need \(\beta=1\). The old Fourier
coefficients have absolute sum one, so the normalized finite
vector is a linear combination of unitary translates of the
vacuum with that absolute sum in every embedding. After
moving the finite arguments, Theorems~\ref{thm:one61}
and~\ref{thm:transport59} bound its conjugated marked section
at the origin by
\[
  \bigl[e\max(1,2r\,\operatorname{rd}(E)/\pi)\,
                 \max(1,\nu_\sigma)\bigr]^{D/2}.
\]
Here \(\nu_\sigma\) is the squared metric normalization of
the marked elliptic section of that theorem. The marked curve
and differential for \(K=\mathbb Q(\sqrt{-d_{\rm K}})\) are
fixed on the interval and have only finitely many conjugates.
This is the degree-zero estimate in the proof of
Theorem~\ref{thm:markeddefinite820}. Its proof uses the
orthogonality of the translated vacuum sections, so the
indices of moved lattices introduce no archimedean factor.
We have \(\operatorname{rd}(E)\le q_1d_{\rm K}\) by the
cyclotomic discriminant bound, \(q_1\ll_{q_0}x^{11}\log(2x)\),
\(\log t_E=O_{q_0,d_{\rm K}}(\log(2x))\) and
\(\log M\le22\log b_x\). It follows that every conjugate of
each normalized kernel times \(b_r\) has positive logarithmic
size at most \(O_{q_0,d_{\rm K}}(D(\log(2x)+\log b_x))\).

It remains to count terms in the norm. At full level
the logarithm of the number of target classes is at most
\(O((2D)^4\log(2Dq_1d_{\rm K}))\) by the definite lattice
count used in Theorem~\ref{thm:definiteclass822}. The same
bound applies to the rank-one source since \(f\le D\).
On passing to marked levels each full class splits into at
most the index of the smaller compact subgroup. The target
index at \(\beta=1\) is \(\prod_{v\in S_0}(q_v^r-1)\le M^D\),
and the source index is at most \(M^f\). The stabilizer
weights are at most one, while \(\log B_r\le(2D)^2\log3\).
For any conjugate of a squared absolute value \(u\bar u\)
bound the two factors under \(\sigma\) and \(\sigma\circ c\)
separately. The triangle inequality, first over source
classes and then over target classes, now bounds the
logarithm in the statement by
\[
 O_{q_0,d_{\rm K}}\bigl((2D)^4\log(2Dq_1d_{\rm K})
                     +D(\log(2x)+\log b_x)\bigr).
\]
Since \(D=fr\ll_{q_0}x^{17}\), this gives the stated estimate.
Finally \(m\le2\sqrt x,\ k\ll x^5,\ b_x\ge x^{400}\) and
\(B_x^{\rm det}=fxb_x(\log b_x)^3/(mk)\). Dividing the bound
by this quantity gives a quantity tending to zero uniformly.
\end{proof}

\begin{corollary}[A multiplier for every factorization of a torus character]
\leavevmode\label{cor:flatmean853}
Let \(G=\widehat C\) be the finite group of characters of the
original torus classes. For two even ranks \(r_1,r_3\) in the
range of the theorem, choose a root \(L\) of
\[
        L^{|G|}=\prod_{\eta\in G}J_{r_1}(\eta,1)J_{r_3}(\eta,1).
\]
It is a nonzero algebraic integer, and its conjugates satisfy
the same order of logarithmic upper bound as in the theorem.
For every coefficient valuation \(v\) and every \(\psi\in G\)
one can factor \(\psi=\psi_1\psi_3\) so that
\[
           v\bigl(J_{r_1}(\psi_1,1)J_{r_3}(\psi_3,1)\bigr)\le v(L).
\]
\end{corollary}
\begin{proof}
The class characters inflate to the allowed tame weight-zero
characters by the definition of \(C\). All factors in this
equation are nonzero algebraic integers by the theorem at
\(Q_\beta=1\). A root of the monic equation is integral over
the algebraic integers and therefore is an algebraic integer.
Taking absolute values at any complex embedding and the
\(|G|\)-th root proves the size estimate. For fixed \(\psi\),
as \(\eta\) runs over \(G\), so does \(\psi\eta^{-1}\).
Consequently the valuations of
\(J_{r_1}(\eta,1)J_{r_3}(\psi\eta^{-1},1)\) have mean \(v(L)\).
Any pair attaining their minimum supplies the factorization.
\end{proof}
\clearpage

\subsection{Four-block tests for the critical torus lift}
\label{sec:fourblock854}
Take a main critical pair \((\nu,z)\) with
\(7m/2<z\le (k+7m)/2\) from Section~\ref{sec:compressionbudget851}.
We use that \(\nu\) has nontrivial tame character at each old mark,
as in Theorem~\ref{thm:rankoneselmer737}; the prescribed \(\beta\)
is ramified at each of its new marks. Set \(n=2z\),
\(t=n-1\), \(d_\tau=c_\tau-z\), and use \(C,K_\beta\) of
Section~\ref{sec:torusnorm852}. At these nodes \(n\ge46\) and
\(z\le s\) since \(k\ge5m+2\) eventually. Set
\[
   I_n=I_{12}\oplus I_4\oplus I_4\oplus I_{n-20}.
\]
Call the first summand \(I_0\). Index the remaining summands by
\(2,3,1\) in that order, with ranks \(r_2=r_3=4,\ r_1=n-20\).
Write \(H_i=U(I_{r_i})\) for these three groups. The character
\(\psi_i\in\widehat C\) specifies flat data on \(H_i\) as in
Theorem~\ref{thm:flattorus853}. Take
\(\beta_2=\beta,\ \beta_3=\beta_1=1\). Use one fixed standard
vector in each flat block for the old mark, and use its
coordinate in block \(2\) for the new marks.

\begin{lemma}[The complementary period of a false torus lift]
\leavevmode\label{lem:falseline854}
Let \(Y\in\operatorname{im}e\cap\ker A_w\cap\ker D\) be a false
form as defined in Section~\ref{sec:rankonepair736}, at a split
\(w\) off the marked places. For a coefficient \(v\) of \(Y\)
fixed by \(H_1(F_\infty)\), a finite integral translation \(A\)
supported on the marks, any argument \(h_+\) on the first
twenty rows, and any flat kernel \(K_1(g,h)\) on the last block,
\[
  \int_{[H_1]}K_1(g,h)\,
      \overline{y_v(\operatorname{diag}(h_+,h)A)}\,dh=0 .
\]
Here \(y_v\) is the scalar coefficient of the algebraic form as
in Section~\ref{sec:falsedescent849}. The kernel on the last
block has the lattice data away from the marks.
\end{lemma}
\begin{proof}
The equation \(DY=0\) gives \(\langle K_\beta(g),Y\rangle=0\)
on the full set of source classes \(C\). This includes every
finite source argument by its level invariance. Integrating
a highest real coefficient of \(Y\) against the block theta
kernel gives the conjugate of this class pairing up to the
positive compact Schur and level factors of
Proposition~\ref{prop:actualadjoint849}. In this rank both real
groups are compact. The homogeneous polynomials of degree \(d\)
in the holomorphic Fock model on one column constitute one
copy of the degree-\(d\) symmetric power, and different degrees
have different central weights. Indeed they are the homogeneous
summands of the polynomial algebra in \(n\) variables, with the
compact group acting by substitution, up to the fixed character.
Schur orthogonality kills all other summands. On the matching
summand simultaneous compact translation carries the highest
polynomial and highest coordinate to their translates, with
the same invariant contraction. Thus zero for that pair gives
zero for every real polynomial and every coefficient, by
irreducibility at each real place. Density of the Hermite
expansions in Schwartz space extends this to all real data,
by continuity of the lattice integral over the compact quotient.

At an old mark the normalized residue action on the form is
ordinary right translation by \(\operatorname{diag}(u,I_t)\)
times the tame character of \(\nu\) on \(u\). It fixes \(Y\)
since the idempotent is the group mean. Consequently the
ordinary translations over \(u\in k^\times\) sum to zero on
\(Y\) by nontriviality of that character. This is exactly the
hyperplane average needed for one column in the proof of
Theorem~\ref{thm:markedbottom847}. More explicitly, on the
old open orbit the top entry of the column is nonzero and the
ratios of its other entries to that entry belong to the maximal
ideal. At a new place the ratios need only be integral.
Averaging over the corresponding target level carries each ball
inside an open orbit into the line of the right scalar translate
of the block test. The descent of every such translate vanishes
because all source arguments are allowed. Outside these orbits
the annihilating row and shears constructed in that proof
fix every sufficiently small ball there. At an old place their
mean is the zero translation sum above. At a new place one shear
acts by a nontrivial tame scalar. Compact disjoint subdivision
therefore proves vanishing for all marked data at the bottom
with the other finite vacua fixed.

We need only the tower through the positive line with at most
twenty hyperbolic planes adjoined. If the first nonzero
restricted lift occurs at rank \(b=1+2h,\ 1\le h\le20\), the
zero-slice calculation of Lemma~\ref{lem:restrictedconstant848}
kills every proper constant term. Its integral hyperbolic
splitting applies to the same unramified extension \(E/F\).
One may take a polynomial real test giving a nonzero cusp
form in \(L^2\), by the continuity, moderate-growth and
cuspidal decay argument of Theorem~\ref{thm:restrictedtower848}.
The normalized variable in \(A_w\) has modulus
\(q^{-(n-2)/2}=q^{-(t-1)/2}\), since its specified Satake
entry is the top of the length-\(t\) string. Applying
Lemma~\ref{lem:rectangularpolynomial848} at the same \(w\)
as in that theorem, the scalar product of degree \(n-b\) is
nonzero because \(n-b<t\). On the spherical \(L^2\) space of
the smaller group the sum of the norms of the positive-degree
terms of the remaining polynomial is at most
\(\exp(q^{-(t-b)/2}/(1-q^{-1}))-1<1\). Here \(b\le41\) and
\(t-b\ge4\). The constant term is the identity, giving a
contradiction. Thus all the restricted descents through
rank \(41\) vanish for each coefficient.

Translate the form by \(A\), which can be moved onto the marked
tests in these integrals. Restrict the rank-\(41\) kernel to
\(U(\mathbb H^{20})\times T\) and use a tensor test. It is
the product of the hyperbolic and the bottom kernels by the
orthogonal-sum oscillator identity. The Fourier coefficient
of the former at Gram matrix \(I_{20}\) unfolds to the
complementary period by Lemma~\ref{lem:frameperiod850} with
frame length twenty. Its proof uses only integral Gram-one
frames, the unmarked lattice vacuum and definiteness of \(I_n\).
The bottom kernel factors over the first twenty and the last
\(n-20\) target rows. Assign the prescribed flat source character
to the latter and the quotient character to the former. Both
have trivial external target character. For fixed \(g,h_+\)
one can make the first factor nonzero by choosing a real bump
isolating one rational position in the finite support, exactly
as in Proposition~\ref{prop:kernelperiod850}. Division by this
nonzero value proves the stated period.
\end{proof}

Choose a residue column \(x_0\) of length twelve at each old
mark and a multi-index \(\alpha=(\alpha_{\tau i})\), \(1\le i\le12\),
with \(\sum_i\alpha_{\tau i}=d_\tau\). Set
\[
    m_\alpha=\prod_\tau\frac{d_\tau!}{\prod_i\alpha_{\tau i}!},
   \qquad
    \chi=\frac{\nu\beta\upsilon_1^z}
                    {(\psi_2\beta)\psi_3\psi_1}.
\]
Let \(K_{0,\chi,\alpha}(g;x_0)\) be the unscaled twelve-row
kernel with external source factor \(\chi^{-1}\), old residue
column \(x_0\), no new residue condition, polynomial
\(\prod_{\tau,i}X_{\tau i}^{\alpha_{\tau i}}\), and target
identity. In particular the first kernel has the indicated
old point test even if \(x_0\) is zero.

\begin{theorem}[A tested integral form with a prescribed derivative monomial]
\leavevmode\label{thm:fourblock854}
For the above choices there is a form \(\Xi\) on \(U(I_n)\),
integral at all coefficient places away from \(Q_\beta\)
and orthogonal to every false torus form in the sense of
Section~\ref{sec:rankonepair736}, such that for each \(g\)
\[
 |C|^3\left(\prod_{i=1}^3 a_{r_i}\right)
           \langle\Xi,K_\beta(g)\rangle
    =m_\alpha K_{0,\chi,\alpha}(g;x_0)
                          \prod_{i=1}^3 J_{r_i}(\psi_i,\beta_i).
\]
The same \(\Xi\) works for every source argument at this twist.
\end{theorem}
\begin{proof}
In the one-row coefficient representation take at each
embedding the tensor \(v_{\alpha,\tau}\) which is the sum of
the distinct words with \(\alpha_{\tau i}\) copies of \(e_i\).
It belongs to the integral lattice of symmetric tensors.
Pairing it with \(X^{\otimes d_\tau}\) gives the monomial times
the multinomial coefficient, because each word contributes
the same monomial with coefficient one. This includes degree
zero. The tensor product \(v_\alpha\) over the embeddings is
fixed by all matrices acting identically on the first twelve
columns. The coefficient representation has no determinant
twist since the source line is positive and the external
target character is trivial.

At an old mark let \(A\) carry the originally specified
column to \((x_0,e_2,e_3,e_1)^{\mathsf t}\), in block notation,
where \(e_i\) are the standard columns in the flat blocks.
This is always a nonzero residue vector and may be completed
to a basis and lifted integrally. At a new mark choose a
permutation for which the original coordinate of \(A^{-1}X\)
reads the prescribed entry of \(X\) in block \(2\). Take
\(A=1\) elsewhere. For \(h=(h_2,h_3,h_1)\) put
\(y=\operatorname{diag}(1,h_2,h_3,h_1)A\) and define the
point form \(F(h)\) by
\[
       \langle F(h),Y\rangle=\langle v_\alpha,Y(y)\rangle
                                  \quad\text{for all target forms }Y .
\]
In the class metric this has support on the class of \(y\);
its value at \(y\) is the sum of the stabilizer translates
of \(v_\alpha\), including the level character. Indeed division
by the stabilizer order in the metric cancels the size of this
sum on taking the pairing with an invariant value. The
stabilizer action preserves the moving coefficient lattice
at \(y\). The inverse of the block diagonal point fixes
\(v_\alpha\), and \(A\) is integral with integral inverse
at every place. Thus \(v_\alpha\) is integral in this lattice
and so is \(F(h)\), including above primes dividing a stabilizer
order.

Pair \(F(h)\) with the full unscaled kernel. The translation
by \(A\) imposes \(x_0,e_2,e_3,e_1\) at the old marks and reads
only block \(2\) at the new marks. The polynomial and source
characters also have the prescribed product over the four
blocks. Rational positions in the direct-sum polarization
sum independently, with absolute convergence, as in
Theorem~\ref{thm:arithmeticclass824}. We get
\[
   \langle F(h),K_\beta(g)\rangle
    =m_\alpha K_{0,\chi,\alpha}(g;x_0)
                         \prod_{i=1}^3 K_i(g,h_i).
\]

For each flat block apply to \(F\), in its \(h_i\) variable,
the weighted class sum
\[
       \frac{B_{r_i}b_{r_i}^2}{Q_{\beta_i}^{\,r_i-1}}
         \sum_{h_i}\frac{\phi_{\psi_i,\beta_i}(h_i)
                                      F(h)}{d(h_i)} ,
\]
where \(d(h_i)\) is the stabilizer order at that block level.
Use block \(1\) first. The summands descend to the stated
classes. A rational block change fixes the coefficient tensor.
The block level fixes its standard old column and, on block
\(2\) at a new mark, preserves the kernel of the selected
coordinate. Its action at the evaluation point therefore has
exactly the block level character. Since the point vector
occupies the conjugate-linear slot, its scalar under a level
translation is the inverse of the scalar on the flat form,
so their product is invariant. This verification persists
when the other two variables are fixed. By
Theorem~\ref{thm:flattorus853}, \(b_{r_i}\phi_{\psi_i,\beta_i}\)
is integral away from \(Q_{\beta_i}\), and \(d(h_i)\mid B_{r_i}\).
Hence each of the three sums preserves integrality above every
\(\ell\nmid Q_\beta\).

For fixed \(h_2,h_3\) and a false form \(Y\), apply the lemma
with \(h_+=\operatorname{diag}(1,h_2,h_3)\) and \(v=v_\alpha\).
The integrand is constant in the compact real variable of
block \(1\), because both \(v_\alpha\) and the flat coefficient
are fixed there. Its adelic integral is the finite class sum
with weights \(1/d(h_1)\) times the positive level volume.
Taking the conjugate and summing the resulting identities over
the source classes gives the vanishing of
\(\sum_{h_1}\overline{\phi_{\psi_1,1}(h_1)}
 \langle F(h),Y\rangle/d(h_1)\).
This is the pairing of the first weighted class sum with \(Y\)
up to its fixed real scalar. The other two finite weighted
sums remain orthogonal to \(Y\). This proves the required
testing for the resulting \(\Xi\), simultaneously for every
\(w\).

Finally, on block \(i\) one has for every \(g\)
\[
     \langle\phi_{\psi_i,\beta_i},K_i(g)\rangle
         =\|\phi_{\psi_i,\beta_i}\|_{\rm cl}^2/(a_{r_i}|C|).
\]
The normalized central translations of
Lemma~\ref{lem:centraltest736} permute the kernels transitively
and fix their sum. Thus the left side is independent of \(g\).
Sum it over \(C\) and use
\(\phi_{\psi_i,\beta_i}=a_{r_i}\sum_gK_i(g)\) to obtain the
displayed value. Conjugate-linearity in the first slot now
replaces each \(K_i(g,h_i)\), under its weighted sum, by
\(J_{r_i}(\psi_i,\beta_i)/(a_{r_i}|C|)\). The four-block
factorization gives the asserted identity.
\end{proof}
\clearpage

\subsection{The valuation of the twelve-row derivative table}
\label{sec:twelftransfer855}
Fix one of the main pairs from Section~\ref{sec:fourblock854}.
Write \(N=|C|\) for the number of source classes and
\(h=\sum_\tau d_\tau\) for the total derivative order. Use
the marked elliptic differential of Theorem~\ref{thm:vac242}
and put \(u_E=\pi/A\) in the notation of that theorem.
For a fixed nonzero algebraic number \(b\), define a table
indexed by the old residue data \(x_0\) and the multi-indices
\(\alpha\) of the previous subsection by
\[
   T_{\alpha,x_0}
     =b\,a_{12}u_E^{h/2}\,
          m_\alpha K_{0,\chi_0,\alpha}(1;x_0),\qquad
       \chi_0=\nu\upsilon_1^z/\psi_2 .
\]
Here \(\psi_2\) is any fixed original tame class character.
One may take the ramified order-\(Q\) character in the
simultaneous critical family for use of its height estimates.
The kernel in the display has identity source and target, so
the value there is independent of \(\chi_0\).
Let \(G_{x_0}(\mathbf Z)\) be the marked twelve-row theta section
in the matrix coordinates of Theorem~\ref{thm:vac242} with
old point \(x_0\). Then
\begin{equation}
 T_{\alpha,x_0}
     =b\left(\prod_\tau d_\tau!\right)
                      [\mathbf Z^\alpha]G_{x_0}(\mathbf Z).
 \label{eq:derivtable855}
\end{equation}
Indeed the marked oscillator identity sends the polynomial
\(\mathbf X^\alpha\) to its undivided derivative, with scalar
\(a_{12}u_E^{h/2}\). At zero this derivative is
\(\prod_{\tau,i}\alpha_{\tau i}!\) times the indicated coefficient.
This cancels the denominator of \(m_\alpha\).
Finite Fourier expansion of the old indicator and the marked
theta expansion in Theorems~\ref{thm:transport59} and~\ref{thm:vac242}
give algebraicity. The matrix variables are obtained from
product coordinates by the algebraic trace-basis matrix of
Lemma~\ref{lem:tracebasis821}.

Use compatible product coordinates in the four even blocks
and their orthogonal sum. In the notation of
Theorem~\ref{thm:torusnorm852} take
\[
       S_{\nu,z}=b\,a_{12}u_E^{h/2}\prod_{i=1}^3 a_{r_i}.
\]
With these coordinates \(a_n=a_{12}\prod_i a_{r_i}\).
Indeed the finite position boxes, lattice intertwiners,
product marked sections and real Gaussian factor on a
direct sum. The identity of the two normalized vacuum
sections holds as functions of their independent variables.
Both are nonzero sections, so their proportionality constant
is the product of those for the summands, even if a value
at the origin vanishes. The marked oscillator identity and
finite Fourier expansion therefore give algebraic entries
for the full scaled \(\phi_\beta\) as well. For general finite
arguments move to split representatives by
Theorem~\ref{thm:one61}, and then move back by rational
automorphy. The changes of polynomial coordinates and
finite determinant characters are algebraic. This proves
algebraicity also for \(\|\phi_1\|_{\rm cl}^2\), by its finite
class sum.

Set \(L>0\) equal to the positive root of
Corollary~\ref{cor:flatmean853} for \(r_1=n-20,\ r_3=4\).
Define
\[
   c_*=L J_4(\psi_2,1),\qquad
   \rho_\beta=J_4(\psi_2,\beta)/J_4(\psi_2,1),\qquad
   H_{\nu,z}=\frac{c_*^2}{N^4\|\phi_1\|_{\rm cl}^2}.
\]
The first and last quantities are independent of the twist.
All three are positive and algebraic. In particular \(c_*\)
is integral and its conjugates have positive logarithmic
size \(o(B_x^{\rm det})\) by Theorem~\ref{thm:flattorus853}.
For a normalized coefficient valuation \(v_\iota\) above
\(\ell\), put \(e_\iota=\min_{\alpha,x_0}v_\iota(T_{\alpha,x_0})\),
with \(v_\iota(0)=+\infty\). This is the valuation of the
fractional ideal generated by the table in any number field
containing it, using the zero ideal convention.

\begin{theorem}[Passing from the untwisted derivative table to the projected period]
\leavevmode\label{thm:twelftransfer855}
For each allowed twist \(\beta\) with \(\ell\nmid Q_\beta\)
and \(e_\iota<+\infty\), there is a vector \(U\) tested and
integral at all primes above \(\ell\) with
\[
        v_\iota(\langle U,\phi_\beta\rangle/\rho_\beta)
                    \le v_\iota(c_*)+e_\iota-2v_\ell(N).
\]
Let \(p_{\ell,\iota}\ge0\) be the bound for the positive part of
the negative valuation of the projected pairing for this main
pair in Corollary~\ref{cor:actualdenom803}. For every
prescribed twist in that simultaneous family,
\begin{equation}
  -v_\iota\bigl(H_{\nu,z}\rho_\beta^2 R_{\nu,z}(\beta)^{-1}\bigr)
                \le p_{\ell,\iota}+e_\iota+e_{\iota\circ c},
       \qquad c(a)=\bar a .
 \label{eq:tablevaluation855}
\end{equation}
In particular the two table valuations use identical finite
data at the identity, with no conductor primes adjoined.
\end{theorem}
\begin{proof}
Choose an index attaining \(e_\iota\), so the corresponding
kernel value at the identity is nonzero. On representatives
\(g\) of \(C\) containing the identity representative, put
\(s(g)=K_{0,\chi_0,\alpha}(g;x_0)\). For
\(\psi\in\widehat C\) form
\[
       A_\psi=N^{-1}\sum_{g\in C}\psi(g)s(g).
\]
Orthogonality of characters gives \(\sum_\psi A_\psi=s(1)\):
the sum of \(\psi(g)\) is zero if \(g\ne1\) in \(C\) and is
\(N\) otherwise. Thus at least one \(\psi\) satisfies
\(v_\iota(A_\psi)\le v_\iota(s(1))\). Equivalently this can be
applied after multiplying every term by the common nonzero
scalar \(b a_{12}u_E^{h/2}m_\alpha\), which makes the
quantities algebraic. This step holds also if \(\ell\mid N\).
Choose \(\psi_1\psi_3=\psi\) by Corollary~\ref{cor:flatmean853}.
It gives \(v_\iota(J_{r_1}(\psi_1,1)J_4(\psi_3,1))\le v_\iota(L)\).

Take \(U\) to be the actual vector of
Theorem~\ref{thm:fourblock854} for these data. With
\(\chi=\chi_0/\psi\), its twelve-row kernel at \(g\) is
\(\psi(g)s(g)\) because the source character occurs inversely.
Multiply the kernel pairing by \(S_{\nu,z}\) and sum on \(C\).
The identity of that theorem, in the chosen normalization, is
\[
 \frac{N^2\langle U,\phi_\beta\rangle}{\rho_\beta}
 = b a_{12}u_E^{h/2}m_\alpha A_\psi\,
           J_{r_1}(\psi_1,1)J_4(\psi_3,1)J_4(\psi_2,1).
\]
The three factors of \(N\) there become two since the
unweighted sum giving \(A_\psi\) is \(N A_\psi\). Substituting
the two valuation estimates just proved gives the first
assertion. In particular its pairing is nonzero.

For the second assertion the inequality is automatic if the
table is zero. Otherwise perform the first construction at
\(\iota\) and at \(\iota\circ c\) with the very same
\(\beta\), obtaining \(U,V\). Both meet the all-embeddings
integrality requirement of Theorem~\ref{thm:rankonepair736}.
Use them in Corollary~\ref{cor:torusvaluation852}. The
two upper estimates for the single pairings contribute at
most \(2v_\iota(c_*\rho_\beta)-4v_\ell(N)
+e_\iota+e_{\iota\circ c}\), since \(c_*\) and \(\rho_\beta\)
are real. The projected term contributes at most \(p_{\ell,\iota}\).
Moving the common scalars inside the negative valuation gives
\eqref{eq:tablevaluation855}. The first construction applied
no restrictions on primes beyond \(\ell\nmid Q_\beta\), so it
applies to both tests without changing the Selmer or critical
choices in that corollary.
\end{proof}

The arithmetic quantities still to be bounded in the directional
estimate are the upper valuations of the explicit Taylor table
in \eqref{eq:derivtable855} and the archimedean size of
\(1/\|\phi_1\|_{\rm cl}^2\) in this same normalization. Changing
\(b\) multiplies the table by the same scalar and the squared
norm by its squared modulus. The contributions of this change
to the two sides of \eqref{eq:tablevaluation855} therefore agree
at every coefficient embedding.
\clearpage

\subsection{The complete linear system of the twelve-row tests}
\label{sec:completesystem856}
Keep the two original marks \(p_0,p_1\) and the chosen position
component at every place above them. Put \(M=p_0p_1\), \(D=12f\).
The period module for the twelve-row vacuum is
\(\Lambda=(\mathfrak b^{-1})^{12}\) in the coordinates of
Theorem~\ref{thm:vac242}. Let \(\mathfrak m\) be the product of
the primes of \(E\) on the chosen position side at the old places,
and put \(\Lambda'=\mathfrak m\Lambda\).
The original marks split completely and are unramified, so
\(\mathfrak m\overline{\mathfrak m}=M\mathcal O_E\).
The inclusion of lattices induces the finite isogeny
\[
       A'=\mathbb C^D/v(\Lambda')
             \longrightarrow A=\mathbb C^D/v(\Lambda),
             \qquad v(a)=\delta\Phi(a).
\]
Let \(\mathcal L\) be the pullback to \(A'\) of the vacuum
theta line on \(A\), with its original multipliers. At the
origin use the same holomorphic coordinates and the same frame
as in \eqref{eq:derivtable855}.

\begin{theorem}[The residue sections and their Taylor coefficients]
\leavevmode\label{thm:completesystem856}
The sections \(G_{x_0}\) indexed by all choices of the twelve
old residue entries at every marked place form a basis of
\(H^0(A',\mathcal L)\). For every \(a\ge0\) with \(3a+2\le M\),
their Taylor polynomials through total degree \(a\) span
\(\mathbb C[\mathbf Z]/(\mathbf Z)^{a+1}\).
\end{theorem}
\begin{proof}
At an old split place the finite test defining \(G_{x_0}\)
is the indicator of the residue column \(x_0\) inside the
unit position lattice. A Weyl translation from \(\Lambda'\)
has position in the maximal ideal times this lattice and
frequency in its full integral dual. Its position shift
preserves the indicator, and its frequency character is one
on the unit lattice. Its scalar Weyl factor is the same as
on the unit vacuum. At the other places the two finite tests
coincide. Thus all the finite tests have the vacuum period
character restricted to \(\widehat{\Lambda'}\).
The translation formula of Theorem~\ref{thm:vac242} gives
the multipliers of \(\mathcal L\) for every \(G_{x_0}\).

Suppose a linear combination of these sections vanishes. Take
the same linear combination of the finite tests and tensor it
with the real vacuum. Its theta function under all real Weyl
translations then vanishes by the coherent-state formula in
Theorem~\ref{thm:vac242}. By Lemma~\ref{lem:zak242} the integral
of its squared absolute value over the real torus for
\(\Lambda'\) equals the squared norm of the tensor. The
distinct finite tests have disjoint supports, each of positive
measure, and the real vacuum has norm one. Every coefficient
of the linear combination therefore vanishes.

Write \(B\) for the integral alternating form of the vacuum.
The trace polarization of Theorem~\ref{thm:pol49} makes
\(\Lambda\) self-dual for \(B\), with the adjoint of
multiplication by \(u\in E\) given by multiplication by \(\bar u\).
Consequently \(v\) in the rational space pairs integrally
with \(\mathfrak m\Lambda\) precisely when
\(\overline{\mathfrak m}v\subset\Lambda\). The dual of
\(\Lambda'\) for \(B\) is
\(\overline{\mathfrak m}^{-1}\Lambda=M^{-1}\Lambda'\).
It follows that \(B/M\) is a positive integral principal
Riemann form on \(A'\). Choose a line bundle \(P_0\) with
this form by the Appell--Humbert theorem. The quotient
\(\mathcal L\otimes P_0^{-M}\) has trivial first Chern class.
Multiplication by \(M\) on the dual complex torus
\(\operatorname{Pic}^0(A')\) is surjective. Twisting \(P_0\)
by an \(M\)-th root of the quotient gives an ample principal
line bundle \(P\) with \(P^{\otimes M}\simeq\mathcal L\).
By Riemann--Roch and vanishing of higher cohomology for an
ample line bundle on a complex abelian variety,
\[
 h^0(A',\mathcal L)
   =M^D\frac{c_1(P)^D[A']}{D!}=M^D.
\]
The last fraction is one by principality. There are \(f\)
places of \(F\) above each \(p_i\), each with residue field
of size \(p_i\). The number of distinct \(x_0\) is thus
\(p_0^{12f}p_1^{12f}=M^D\), which proves the basis assertion.

Use the very ampleness of \(P^3\) on a complex abelian variety
(\href{https://www.jmilne.org/math/CourseNotes/AV.pdf}{Milne,
\emph{Abelian Varieties}, Chapter I, \S6}). Choose a section
\(s_0\) of \(P^3\) nonzero at the origin. The ratios of global
sections to \(s_0\) give coordinates in the affine chart of
the projective embedding. After subtracting their values at
zero, choose \(D\) such ratios \(s_i/s_0\) with independent
differentials. They form local holomorphic coordinates. For
\(|\gamma|\le a\) the section
\(s_0^{a-|\gamma|}\prod_i s_i^{\gamma_i}\) of \(P^{3a}\)
has, in the frame \(s_0^a\), the corresponding monomial in
these coordinates. These monomials span every jet of order \(a\).

Put \(u=M-3a\ge2\). There is a section of \(P^u\) nonzero
at zero. Indeed take a nonzero section of \(P\), which exists
by Riemann--Roch, with divisor \(D_0\). Choose \(u-1\)
translation parameters on \(A'\) and take the last to be
the negative of their sum. They may all be chosen outside
\(D_0\), since the coordinate projections and the negative
sum are surjective and finitely many proper closed subsets
do not cover \((A')^{u-1}\). The product of the translated
sections is then nonzero at zero and belongs to \(P^u\) by
the theorem of the square. Multiplication by this section
identifies the two local jet spaces. A change of holomorphic
coordinates with invertible differential and a change of
frame by a unit also act invertibly on finite jets. We
therefore have the asserted spanning property in the actual
frame and coordinates of the derivative table.
\end{proof}

\begin{corollary}[Nonzero tests at every prescribed derivative]
\leavevmode\label{cor:nonzerotable856}
On the sufficiently distant common intervals of
Theorem~\ref{thm:twelftransfer855}, the matrix
\((T_{\alpha,x_0})\), with rows of multidegree \((d_\tau)_\tau\),
has full row rank over any number field containing its entries.
In particular each such row has a nonzero entry. The valuation
\(e_\iota\) in \eqref{eq:tablevaluation855} is finite at every
coefficient place. For a fixed main pair it is zero outside
a finite set of places.
\end{corollary}
\begin{proof}
We check the order \(h=\sum_\tau d_\tau\) from
Section~\ref{sec:twelftransfer855}. The original weights
\(w_{j,\tau}=\kappa_\tau+(2j+1)s^8\) are at most \(2b_x\)
by the estimates of Theorem~\ref{thm:gluedcost811}. The
exponent of \(\alpha\) in the scalar exchange is
\(m-\mu_\tau\), where
\(\mu_\tau=w_{j,\tau}-j-(t+k-1)/2\). The exponent of
\(\sigma_c=\alpha^{-1}\upsilon_1^{3m+e'}\) has absolute
value at most \(|m-\mu_\tau|+3m+e'\). Passing to the inverse
has the same absolute value, and a tame twist has weight
zero. Thus \(d_\tau=c_\tau-z\ge0\) is \(O(b_x)\) uniformly,
since \(z\le e'\le s,\ t=5s,\ s\le3x^5\) and \(m\le2\sqrt x\).
Now \(f\ll_{q_0,d}x^{12}\), \(b_x\ge x^{400}\) and
\(M\ge b_x^{20}\) by the original prime choices
(Corollary~\ref{cor:actualmarks714}). Hence \(3h+2\le M\)
throughout all sufficiently distant intervals.

Project the space of \(h\)-jets onto the coordinates belonging
to the required multidegree. This projection is onto, so the
corresponding coefficient rows for the \(G_{x_0}\) are
linearly independent by the theorem. In \eqref{eq:derivtable855}
all entries have the same nonzero multiplier
\(b\prod_\tau d_\tau!\). A maximal row minor remains nonzero
after this scaling, and the entries are algebraic there.
Finally the table has finitely many entries. Outside the
finite set of primes dividing a numerator or denominator of
one of its nonzero entries, every such entry is a unit.
\end{proof}

This removes the nonvanishing hypothesis in the construction
of the tested vector in Theorem~\ref{thm:twelftransfer855}.
The directional prime sum also requires a quantitative bound
on the valuations of the table as the field and weights vary.
\clearpage

\subsection{Integral leading terms in the twelve-row linear system}
\label{sec:integralleading857}
Put \(D=12f\), \(M=p_0p_1\), and \(I=\mathfrak m\mathfrak b^{-1}\)
as in Section~\ref{sec:completesystem856}. For a fractional ideal
\(J\) of \(E\), write \(v_\iota(\tau J)\) for the valuation of a
local generator of its extension via \(\tau\) to the coefficient
ring at \(\iota\). Normalize \(v_\iota(\ell)=1\).
The following estimate uses the first nonzero homogeneous term
of an integral section. This makes the comparison independent
of the higher terms in the change to the marked analytic frame.

\begin{lemma}[Extracting one residue coefficient over a valuation ring]
\leavevmode\label{lem:residuecoeff857}
Let \(R\) be a local integer ring of a sufficiently large number
field, and let \(\mathcal L_R\) be an extension of the line bundle
in Theorem~\ref{thm:completesystem856} to the good abelian model
of \(A'\). Fix a residue column \(y\) with \(g=G_y(0)\ne0\).
Let \(\lambda_R\) be the value, in the marked frame, of a generator
of the fiber of \(\mathcal L_R\) at zero. If
\(S=\sum_x c_xG_x\) is an integral global section, then
\[
                  M^D c_x g/\lambda_R\in R\qquad\text{for every }x .
\]
\end{lemma}
\begin{proof}
The functions are identified with algebraic sections over a
common number field. Indeed the marked product theta has
algebraic divisor. By the torsion translation formula in
Theorem~\ref{thm:transport59}, the vacuum and its finite
translates are a common product theta times algebraic rational
functions on a finite isogenous cover. The finite Fourier
expansion of a residue indicator has algebraic coefficients.
Their ratios therefore descend to algebraic rational functions
on \(A'\), since on the cover the ratios are invariant under the
kernel of the covering isogeny. The common bundle is algebraic
as well. Its divisor can be obtained from a nonzero section
on this cover and descended: pullback of its complex divisor
is algebraic, and pushforward recovers the divisor times the
degree of the isogeny. We can use \(G_y\) as rational frame
near zero. The new frame has value \(g\) in the marked frame,
an algebraic number by the order-zero torsion formula. This
also defines the algebraic comparison at the fiber.

Denote the additive group of residue columns by \(X\), of order
\(M^D\). For each \(\xi\in\widehat X\) use multiplication by
\(\xi\) on the old unit position box. Its operator on sections
satisfies \(D_\xi G_x=\xi(x)G_x\). A unit position shift
gives an operator carrying \(G_x\) to any specified \(G_y\)
up to a root of unity. These formulas follow on the indicators
by multiplication and translation, respectively. They are
translations on the analytic theta line by the Weyl formula
of Theorem~\ref{thm:vac242}. More explicitly each corresponding
finite Weyl displacement pairs integrally with
\(\widehat{\Lambda'}\). Choose a rational representative modulo
\(\widehat{\Lambda'}\) by lattice approximation. It belongs
to the \(B\)-dual \(M^{-1}\Lambda'\), so the translation on
\(A'\) is by a torsion point. Replacing a representative
introduces only the vacuum period character and a rational
Weyl phase, both roots of unity. We take the scalar giving
the displayed action for \(D_\xi\). A positive power of each
chosen lift is the identity up to a root of unity.

We verify integrality of these operators even above \(M\).
Extend the field to define the torsion points. Each translation
extends to the abelian scheme by properness and the group law.
The isomorphism of the translated line bundle with the original
one exists algebraically on the generic fiber. An algebraic
choice differs from the analytic lift by a scalar, algebraic
since a positive power of each lift acts by an algebraic scalar.
After multiplication by an element of the fraction field we
obtain an isomorphism over \(R\). Indeed the divisor of the
generic isomorphism has only vertical terms, and the special
fiber is irreducible, reduced and cut out by a uniformizer.
A power of this integral isomorphism covering the identity
is a unit scalar, since the global invertible functions on
the abelian scheme are those of \(R\). The original lift has
a root of unity for that scalar after taking a common power.
Their ratio thus has valuation zero. In particular the original
operators act on integral sections.

The sum \(\sum_{\xi\in\widehat X}\xi(x)^{-1}D_\xi S\) is
integral and equals \(M^D c_xG_x\), by character orthogonality.
Apply a shift carrying \(x\) to \(y\), then evaluate at zero.
This gives \(M^Dc_xg\) in the integral fiber, up to a unit
root of unity, as asserted.
\end{proof}

\begin{theorem}[An upper bound from integral leading terms]
\leavevmode\label{thm:integralleading857}
At each sufficiently distant node there is a single nonzero
algebraic integer \(\gamma_x\), depending only on the original
field, marks and vacuum, with
\[
    \sup_\sigma \log^+|\sigma\gamma_x|
                     \ll_{q_0,d_{\rm K}} f\log(2b_x).
\]
For each permitted row \(\alpha\), every main pair, every
nonzero scalar \(b\) of \eqref{eq:derivtable855}, and every
coefficient place one has
\begin{equation}
 \min_{y}v_\iota(T_{\alpha,y})\le
 v_\iota(b\gamma_x)+\sum_\tau
   \left(v_\ell(d_\tau!)
           +d_\tau\bigl(v_\iota(\tau\mathfrak b)
                         -v_\iota(\tau\mathfrak m)\bigr)\right).
 \label{eq:integralleading857}
\end{equation}
\end{theorem}
\begin{proof}
Use the good model \((I^{12})\otimes_{\mathcal O_K}C\) for \(A'\),
where \(C\) is the marked elliptic curve. The construction by a
projective summand and its complex uniformization are those of
Theorem~\ref{thm:pol49}. After enlarging the number field
choose the root \(P^M\simeq\mathcal L\) constructed in
Theorem~\ref{thm:completesystem856} algebraically. Indeed the
fiber of the \(M\)-th power map in the Picard scheme has a complex
point and hence a point over \(\overline{\mathbb Q}\).
Locally at a finite place extend \(P\) by closure of a divisor
on the regular model. Its extension \(P_R\) is fiberwise ample.
To see this, the Hermitian form for \(P\) is \(1/M\) times
the positive trace form. Diagonalize this form over \(K\)
and clear denominators in both its diagonal entries and its
independent coordinate forms on \(I^{12}\). A positive integer
multiple is a sum of positive integral multiples of norms of
\(\mathcal O_K\)-linear coordinates \(I^{12}\to\mathcal O_K\).
Together these coordinates give an isogeny to a product of
copies of \(C\) on every fiber, since an integer multiple of
the inverse also preserves the lattices. The corresponding
product of positive powers of the elliptic polarization pulls
back to an ample bundle fiberwise. Its polarization homomorphism
agrees with the appropriate power of \(P_R\) on the generic
fiber, hence on the whole model by separatedness, as in the
proof of Theorem~\ref{thm:pol49}. Fiberwise their quotient
lies in \(\operatorname{Pic}^0\), proving ampleness.

Take \(\mathcal L_R=P_R^M\). Its sections map onto all
order-\(h\) jets at the zero section over \(R\), with
\(h=\sum_\tau d_\tau\). On a geometric special fiber the
cube of an ample bundle is very ample and its powers of
degree at least two are generated by sections. The coordinate
monomials in the proof of Theorem~\ref{thm:completesystem856}
therefore give surjectivity there as \(3h+2\le M\). These
statements on powers, and vanishing of higher cohomology of
an ample bundle on an abelian variety, hold in every
characteristic (Mumford, \emph{Abelian Varieties}, Chapter III,
\S\S16--17). Cohomology and base change lifts sections of
the special fiber. Jets form a finite free module over the
local ring by smoothness, so Nakayama's lemma gives the
asserted surjectivity over \(R\).

The marked elliptic differential generates the invariant
differentials in good reduction by the unit-discriminant
calculation in Theorem~\ref{thm:elliptic312}. In the \(i\)-th
copy of \(I\otimes C\), the differential \(dZ_{\tau i}\)
in matrix coordinates is the linear form \(a\mapsto\tau(a)\)
on the module \(I\), in that elliptic differential. This
follows by pulling back along \(z\mapsto(\tau(a)z)_\tau\)
for \(a\in I\). If \(u_\tau\) generates the inverse of the
fractional ideal \(\tau(I)R\), each \(u_\tau dZ_{\tau i}\)
is thus an integral cotangent vector. The top term of the
jet module with all lower terms zero is the \(h\)-th symmetric
power of the cotangent module tensored with the fiber of the
bundle. Prescribe there the monomial
\(\prod_{\tau,i}(u_\tau dZ_{\tau i})^{\alpha_{\tau i}}\)
times a fiber generator, and lift it to an integral section
\(S\). In marked coordinates it satisfies
\[
            [\mathbf Z^\alpha]S
                       =\lambda_R\prod_\tau u_\tau^{d_\tau}.
\]
Only the differential of the change of parameters and the
value of the frame at zero enter this homogeneous term. This
also holds at \(h=0\) by evaluation.

Fix once and for all a column \(y_*\) with \(g=G_{y_*}(0)\ne0\);
it exists by Theorem~\ref{thm:completesystem856} at order zero.
Expand \(S=\sum_y c_y G_y\). By the lemma each
\(M^Dc_yg/\lambda_R\) belongs to \(R\). Taking the indicated
Taylor coefficient in this expansion shows that \(M^Dg
\prod_\tau u_\tau^{d_\tau}\) is an integral linear combination
of the numbers \([\mathbf Z^\alpha]G_y\). Multiply by
\(b\prod_\tau d_\tau!\) and use \eqref{eq:derivtable855}.
The valuation of the ideal on the right is at most the valuation
of this nonzero combination. Here
\(v_\iota(u_\tau)=v_\iota(\tau\mathfrak b)
                  -v_\iota(\tau\mathfrak m)\).

Finally take \(\gamma_x=(Mt_E)^D g\), with \(t_E\) from
Lemma~\ref{lem:tracebasis821}. We check its size and integrality
in this same marked frame. The finite Fourier expansion of the
old point indicator has denominator \(M^D\) and its coefficients
have absolute sum one at each embedding. The passage to the
product theta has normalization \(I_0^{-1/2}\) with
\(I_0\mid t_E^{2D}\) by that lemma. Each remaining summand
at zero is a product of torsion values with a root of unity
factor and is integral by Theorem~\ref{thm:torsion312}.
It follows that \(\gamma_x\) is integral. At any complex
embedding the squared magnitude of each translated vacuum
at zero is at most
\([e\max(1,24\operatorname{rd}(E)/\pi)\max(1,\nu_\sigma)]^D\)
by Theorem~\ref{thm:transport59}. The finitely many elliptic
norms are bounded for fixed \(d_{\rm K}\). The bounds
\(\log M\le22\log b_x\), \(\log t_E\ll_{q_0,d_{\rm K}}\log(2x)\)
and \(\operatorname{rd}(E)\le q_1d_{\rm K}\) from the proof of
Theorem~\ref{thm:flattorus853} give the asserted estimate
for every conjugate. Since \(t_E\) is an integer, replacing
\(M^D g\) by \(\gamma_x\) preserves the valuation upper bound.
\end{proof}

\begin{corollary}[The projected period above the derivative weights]
\leavevmode\label{cor:largeprimetable857}
Let \(c\) denote complex conjugation on the coefficient numbers.
For an allowed twist with \(\ell\nmid Q_\beta\), suppose that
\(\ell>\max_\tau d_\tau\) and that \(\ell\) is prime to the
absolute norm of the relative different of \(E/K\). Then on
the same field and interval as Corollary~\ref{cor:actualdenom803},
\[
 -v_\iota\!\left(
  H_{\nu,z}\,b\bar b\,\gamma_x\bar\gamma_x\,
             \rho_\beta^2 R_{\nu,z}(\beta)^{-1}\right)
       \le p_{\ell,\iota}.
\]
This holds in particular at all primes larger than \(C b_x\)
for a fixed sufficiently large \(C=C(q_0,d_{\rm K})\) and
all sufficiently distant nodes, subject to \(\ell\nmid Q_\beta\).
\end{corollary}
\begin{proof}
All factorials in \eqref{eq:integralleading857} are units at
\(\ell\). The relative different is
\(\mathfrak b\overline{\mathfrak b}\), and \(\mathfrak m\)
is integral. Hence \(e_\iota\le v_\iota(b\gamma_x)\), also
at \(\iota\circ c\) with the same \(\gamma_x\). Substitution
into \eqref{eq:tablevaluation855} proves the display.
The weights are \(O(b_x)\) by Corollary~\ref{cor:nonzerotable856}.
The cyclotomic extension defining \(E/K\) ramifies only above
the primes dividing its conductor \(q_1\ll_{q_0}x^{11}\log(2x)\),
which is smaller than \(b_x\) on these intervals. This gives
the last assertion.
\end{proof}

At primes within the derivative weights the factorials in
\eqref{eq:integralleading857} can contribute. The directional
estimate requires finer control of these contributions together
with the archimedean size of the inverse untwisted norm.
\clearpage

\subsection{Formal theta frames and the split primes}
\label{sec:twocopy858}
Write \(B=I\otimes_{\mathcal O_K}C\) for one of the twelve
factors of the abelian scheme of Section~\ref{sec:integralleading857}.
Use its marked matrix coordinates \(Z=(Z_\tau)_\tau\) on the
complex universal cover and, by extension of the coefficients,
as formal logarithm coordinates at a finite place. In these
coordinates the group law is addition. Let \(R\) be a completed
local integer ring as in that section. We may make a further
extension to a complete discrete valuation ring, keeping
\(v_\iota(\ell)=1\).

\begin{lemma}[Cancellation of the theta frame along a formal diagonal]
\leavevmode\label{lem:diagonalframe858}
There are \(\delta_1,\ldots,\delta_6\in\mathbb Z_\ell\), with
\(\delta_1=1\) and \(\sum_i\delta_i^2=0\), with the following
property. Let \(L\) be the one-factor theta line whose external
twelfth power on \(B^{12}\) is \(\mathcal L\).
Take an integral local frame for an extension \(L_R\) along the
zero section. Write \(F(Z)\) for its scalar in the marked
analytic chart divided by its value at zero, so \(F(0)=1\).
Then \(\prod_{i=1}^6 F(\delta_i Z)\) and its inverse belong
to the integral completed local ring of \(B\) along zero.
\end{lemma}
\begin{proof}
The standard trace Hermitian form is the orthogonal sum over
the twelve entries. The period sign is
\((-1)^{\operatorname{Tr}_{E/K}\sum_i a_i\bar a_i}\) on
\(I^{12}\) by Lemma~\ref{lem:basis59}, so the factor of
automorphy is the product of the identical one-entry factors.
This identifies \(L\) by restriction to one factor and gives
the asserted external product in the actual analytic chart.
The algebraicity argument in Lemma~\ref{lem:residuecoeff857}
applies to the restriction. Concretely a residue section nonzero
at zero restricts to an algebraic local frame, and its Taylor
series after division by its value there has algebraic
coefficients. The ratio of another algebraic local frame to
this one is an algebraic function. Hence all the finite jets
of \(F\) can be read in the local coefficient field.

For each nonnegative integer \(n\), pullback by \([n]\) gives
the same analytic multipliers as \(L^{n^2}\). The exponential
part of the factor is quadratic under simultaneous scaling of
the period and variable, and its sign is unchanged on replacing
the exponent \(n\) by \(n^2\). The resulting identification
of rigidified bundles is algebraic by the theorem of the cube.
Normalize it to send the chosen fiber generator to its
\(n^2\)-th power at zero. It extends integrally: as in the
proof of Lemma~\ref{lem:residuecoeff857} a generic isomorphism
extends up to a power of the uniformizer, and normalization
on the fiber at the zero section fixes that power. Therefore,
in integral parameters \(T\) at zero,
\[
                   F(nZ)/F(Z)^{n^2}\ \in R[[T]]^\times .
\]
Indeed this is the ratio of the pulled-back integral frame
and its tensor power under that isomorphism. The normalization
at zero cancels their constant factors.

For each fixed jet order, the coefficients of this ratio in
\(T\) depend continuously on \(n\in\mathbb Z_\ell\).
In the logarithm coordinates substitution multiplies a term
of degree \(j\) by \(n^j\). For the power \(F(Z)^{-n^2}\)
use the finite binomial expansion in \(F-1\) through this
order. Both calculations have coefficients polynomial in \(n\)
over the fraction field. The change from logarithm coordinates
to \(T\) is fixed and invertible in characteristic zero.
Consequently \(F(\delta Z)/F(Z)^{\delta^2}\) has integral
coefficients and constant one for every \(\delta\in\mathbb Z_\ell\),
by closure of \(R\).

For each \(N\ge1\) write the nonnegative integer \(\ell^N-1\)
as four squares by Lagrange's theorem. Compactness of
\(\mathbb Z_\ell^4\) gives a convergent subsequence of the
four coordinates as \(N\) tends to infinity. Their limiting
squares sum to \(-1\). Prepend \(1\) and append \(0\) to
obtain the stated six scalars. Multiplying the six integral
ratios cancels the powers of \(F(Z)\), by the binomial law
for a formal unit of constant one. The product has constant
one, so its inverse is integral as well.
\end{proof}

\begin{theorem}[Two scalar formal groups in the twelve-row tests]
\leavevmode\label{thm:twocopy858}
Let \(H\) be an integral scalar function on the formal completion
of \(B\times B\) at the identity. In the two sets of logarithm
coordinates \(U,V\), put
\(q=[U^\mu V^\eta]H\), where
\(\mu_\tau+\eta_\tau=d_\tau\) for every \(\tau\).
For \(q\ne0\) the derivative table of the same field and marks
satisfies
\begin{equation}
          e_\iota\le v_\iota(b\gamma_x)
                         + v_\iota\!\left(q\prod_\tau d_\tau!\right).
 \label{eq:twocopy858}
\end{equation}
The common algebraic integer \(\gamma_x\) is precisely the
one in Theorem~\ref{thm:integralleading857}.
\end{theorem}
\begin{proof}
Choose \(L_R\) as in the lemma, and on \(B^{12}\) use its
external product, with the identification on the generic
fiber just proved. Over a local integer ring this model is
isomorphic to the model of \(\mathcal L\) used in
Theorem~\ref{thm:integralleading857}: any discrepancy between
extensions over the same regular total space is a vertical
divisor, hence principal. The integral surjection onto
\(h\)-jets established there thus applies to this model,
also after flat extension of the local integer ring.

Identify the first six logarithm variables with
\((\delta_1 U,\ldots,\delta_6 U)\) and the last six with
\((\delta_1 V,\ldots,\delta_6 V)\). Write
\(J(U)=\prod_i F(\delta_i U)\).
The function \(H(U,V)/(J(U)J(V))\) is integral in the
formal group parameters of \(B^2\) by the lemma. Pull
that function back along the projection onto factors \(1,7\)
of \(B^{12}\). In the external product frame use its Taylor
jet through order \(h=\sum d_\tau\). Lift this integral jet
to a global section \(S\) by the surjection. Let \(\lambda_R\)
be the value of this product frame at zero in the marked
analytic chart. After the indicated linear substitutions,
the marked expansion of \(S/\lambda_R\) agrees through degree
\(h\) with \(H(U,V)\), because \(\delta_1=1\) on both selected
factors and the other frame factors restore \(J(U)J(V)\).

Now expand \(S=\sum_y c_y G_y\). By
Lemma~\ref{lem:residuecoeff857}, \(M^D c_yg/\lambda_R\) is
integral for the single nonzero \(g=G_{y_*}(0)\) chosen in
the previous theorem. This assertion commutes with the flat
base extension, or follows directly by the same integral
translations. The coefficient at \(U^\mu V^\eta\) of a
substituted \(G_y\) is a sum of coefficients of that section
in the original twelve variables, with integral multipliers
given by products of powers of \(\delta_i\). Each index in
this sum has total \(d_\tau\) at \(\tau\). Thus
\(M^Dg\,b(\prod_\tau d_\tau!)q\) is an integral linear
combination of the actual \(T_{\alpha,y}\).
Since \(\gamma_x=t_E^D M^Dg\) with \(t_E\) integral, taking
valuations proves \eqref{eq:twocopy858}.
\end{proof}

\begin{corollary}[The projected theta period at every split prime]
\leavevmode\label{cor:splitperiod858}
Suppose \(\ell\) splits in \(K\). For every main pair, with
no restriction on its derivative weights, one has
\[
 e_\iota\le v_\iota(b\gamma_x)
       -\sum_\tau d_\tau\,v_\iota
                      \bigl(\tau(\overline{\mathfrak b}\mathfrak m)\bigr).
\]
In particular at every such \(\ell\nmid Q_\beta\), for the
same twists and integral local cohomological bound as in
Corollary~\ref{cor:actualdenom803},
\[
 -v_\iota\!\left(H_{\nu,z}\,b\bar b\,\gamma_x\bar\gamma_x\,
               \rho_\beta^2 R_{\nu,z}(\beta)^{-1}\right)
                         \le p_{\ell,\iota}.
\]
\end{corollary}
\begin{proof}
We first produce a multiplicative character of the integral
formal group of \(C\) whose logarithmic derivative at zero
is a unit. At a split prime
\(\mathcal O_K\otimes\mathbb Z_\ell=\mathbb Z_\ell\oplus\mathbb Z_\ell\)
splits the \(\ell\)-divisible group into two height-one
factors. The heights follow on the generic Tate module
from the lattice of \(C\), which has rank one over
\(\mathcal O_K\). Exactly one idempotent acts on the tangent
line of the good model. The other factor is \'etale: its
finite torsion levels have zero tangent space on the special
fiber and are finite flat. The elliptic principal polarization
interchanges the two idempotents under Cartier duality, since
its Rosati involution on \(K\) is conjugation. After passage
to the completed strict henselization the \'etale factors are
constant, with compatible cyclic generators. Their Cartier
duals identify the formal torsion of \(C\) at every level
with \(\mu_{\ell^n}\).

For completeness take an integral parameter \(T\) at zero.
The compatible multiplicative coordinates on these torsion
groups define a unit \(X(T)\) of constant one over the completed
ring. Indeed at level \(n\) the coordinate is defined modulo
\([\ell^n](T)\). These ideals are nested and lie in
\((\pi,T)^{n+1}\), where \(\pi\) is a uniformizer, because
\([\ell](T)=\ell T+O(T^2)\). Compatible polynomial or power
series lifts therefore converge in the completed local ring.
They give the prescribed coordinates modulo every such ideal.
The group law for \(X\) follows by the same argument in two
variables from the laws on all torsion levels. Its linear
coefficient is a unit, since on the special fiber the
isomorphism at level one induces an isomorphism on tangent
lines, and the differential of \([\ell]\) there is zero.
The form \(dX/X\) is invariant and is \(u\,dz\) for a unit
\(u\), by the integral differential calculation in
Theorem~\ref{thm:elliptic312}. Consequently in the marked
characteristic-zero logarithm coordinate
\(X(z)=\exp(u z)\), by the differential equation and \(X(0)=1\).

The trace dual of \(I=\mathfrak m\mathfrak b^{-1}\) over
\(\mathcal O_K\) is
\(\mathfrak D_{E/K}^{-1}I^{-1}
       =(\overline{\mathfrak b}\mathfrak m)^{-1}\).
Choose \(w\) in this ideal generating it at all primes above
\(\ell\). For example, modulo each such prime times the ideal
choose a nonzero class, and use the Chinese remainder theorem
on the ideal. The trace map \(a\mapsto\operatorname{Tr}_{E/K}(wa)\)
then takes \(I\) into \(\mathcal O_K\) and defines a morphism
\(B\to C\) over the good model. Its logarithm coordinate is
\(\sum_\tau \tau(w)U_\tau\) by the trace formula. Pull back
\(X\) along this map on the first factor of \(B^2\), taking
the second variable constant. The result is an integral
function with
\[
   \left(\prod_\tau d_\tau!\right)[U^{\mathbf d}V^0]H
                        =u^{h}\prod_\tau \tau(w)^{d_\tau}.
\]
The first inequality follows from \eqref{eq:twocopy858} and
the local generator property of \(w\). Both \(\mathfrak b\)
and \(\mathfrak m\) are integral, so their displayed correction
can be discarded for an upper bound, at \(\iota\) and at
\(\iota\circ c\). Substitution into
\eqref{eq:tablevaluation855} gives the period inequality.
\end{proof}

The large-prime bound of Corollary~\ref{cor:largeprimetable857}
and this split-prime estimate use the same scalar. For
nonsplit primes within the weights, \eqref{eq:twocopy858}
reduces the remaining valuation problem to integral functions
on the two indicated formal groups.
\clearpage

\subsection{The ramified prime and the pairs of larger weights}
\label{sec:shortpairs859}
At a prime \(\ell\) nonsplit in \(K\), let \(e_0\) and \(Q\) be
the ramification index and the residue cardinality of \(K_v\).
Write
\[
           \lambda_\ell=\frac1{\ell-1}-\frac1{e_0(Q-1)} .
\]
We first bound the Taylor table using the marked elliptic
formal group alone. All valuations have \(v_\iota(\ell)=1\).

\begin{theorem}[An elliptic derivative bound at a nonsplit prime]
\leavevmode\label{thm:scalartrace859}
For \(h=\sum_\tau d_\tau\) and the same \(\gamma_x\) as in
Theorem~\ref{thm:integralleading857},
\[
 e_\iota\le v_\iota(b\gamma_x)+h\lambda_\ell
         -\sum_\tau d_\tau\,v_\iota
                 \bigl(\tau(\overline{\mathfrak b}\mathfrak m)\bigr).
\]
\end{theorem}
\begin{proof}
Take an integral parameter \(T\) of the good elliptic formal
group as in Lemma~\ref{lem:unit41}. Write \(\partial_z\) for
differentiation in the marked logarithm coordinate. Define \(a_n\)
as the minimum valuation of \((\partial_z^n T^j)(0)\) for
\(0\le j\le n\). This is the minimum over all integral functions,
because their Taylor series in \(T\) can be truncated through
degree \(n\) for this derivative. It is finite by \(j=n\) and
nonnegative since \(\partial_z\) is an integral invariant
derivation by Theorem~\ref{thm:elliptic312}. Also \(a_0=0\).

We have \(a_{n+m}\ge a_n+a_m\). To see this, pull an integral
function back by formal group addition to a function integral
in the two separate parameters. Its derivative of bidegree
\((n,m)\) has valuation at least \(a_n+a_m\), by expanding in
monomials of the two parameters. In the two logarithm coordinates
this is exactly the derivative of order \(n+m\) of the original
function at zero.

For each \(r\ge0\), the coefficient of \(z^{Q^r}\) in \(T(z)\)
has valuation \(-(Q^r-1)/(e_0(Q-1))\). Indeed in
Lemma~\ref{lem:unit41} the parameter \(T\), after the rescaling
by \(\lambda\) with \(\lambda^{Q-1}\) a uniformizer of \(K_v\),
has a nonzero reduced coefficient at every \(Q^r\). The linear
change from that logarithm coordinate to the marked elliptic
one is a unit because both differentials generate. Since \(Q^r\)
is an \(\ell\)-power, Legendre's formula gives
\(v_\ell((Q^r)!)=(Q^r-1)/(\ell-1)\). Hence
\(a_{Q^r}\le(Q^r-1)\lambda_\ell\). Fix \(n\ge1\), divide
\(Q^r\) by \(n\) with quotient \(j_r\), and use superadditivity
and nonnegativity to get \(j_r a_n\le (Q^r-1)\lambda_\ell\).
Letting \(r\) increase gives \(a_n\le n\lambda_\ell\).

Choose \(j\) realizing \(a_h\) (take the constant for \(h=0\)).
Pull \(T^j\) back to \(B\) by the integral trace morphism
associated to the generator \(w\) of the trace dual at the
primes above \(\ell\), as in Corollary~\ref{cor:splitperiod858}.
This construction of the trace morphism works for every
decomposition type of \(\ell\). In \eqref{eq:twocopy858}
use this function on the first factor with \(\mu_\tau=d_\tau\)
and \(\eta_\tau=0\). The multinomial formula gives
\[
       \left(\prod_\tau d_\tau!\right)[U^{\mathbf d}]H
           =(\partial_z^h T^j)(0)\prod_\tau\tau(w)^{d_\tau}.
\]
Using \(v_\iota(\tau(w))=-v_\iota
(\tau(\overline{\mathfrak b}\mathfrak m))\) proves the assertion.
\end{proof}

For the shorter family put \(G=xb_x/x_{1,x}\), and write
\(\kappa_\tau=b_x\) on the specified set \(\mathcal T\) of
larger-weight embeddings and \(\kappa_\tau=G\) elsewhere.
Use the choice of \(\mathcal T\) in Lemma~\ref{lem:spacing781}.
Recall that the oriented embeddings are indexed by
\((\mathbb Z/q_1)^\times\). Every nonsplit \(\ell\) is
prime to \(q_1\), since the primes of its conductor split
in \(K\) by the proof of Theorem~\ref{thm:fixedconductor781}.
Write \(P_a\tau=-\ell^a\tau\bmod q_1\).
Let \(J_\ell\) be the set of positive odd integers at an
inert prime and all positive integers at the ramified prime.
The directional local sum is
\[
     S_\ell=\frac1{e_0}\sum_\tau\sum_{a\in J_\ell}
                  \bigl(\kappa_\tau/\ell^a-\kappa_{P_a\tau}\bigr)_+ .
\]
All the nonnegative series below are bounded by geometric series.

\begin{theorem}[The ramified prime and small inert primes]
\leavevmode\label{thm:shortpairs859}
On the same nodes and intervals as Theorem~\ref{thm:twelftransfer855},
put \(n_x=\lceil4b_x/k\rceil\), \(A_x=(n_x!)^f\).
For a nonsplit \(\ell\le k\) let \(L\) be the least positive
integer with \(\ell^L\ge k\ell\), and set
\[
 E_\ell=\frac1{e_0}\!
   \sum_{\substack{\tau\in\mathcal T,\ a\in J_\ell,\ a\le L\\
                           P_a\tau\in\mathcal T}}\frac{b_x}{\ell^a}.
\]
At every coefficient place on sufficiently distant intervals,
\[
             e_\iota\le v_\iota(b\gamma_x A_x)+S_\ell+E_\ell .
\]
The sum \(E_\ell\) is zero at the ramified prime. At an inert
prime its terms come from disjoint pairs
\(\{\tau,P_a\tau\}\) of distinct embeddings. In particular it
vanishes if \(-\ell^a\not\equiv1\bmod q_0\) at each of the
odd \(a\le L\).
For every prescribed twist at \(\ell\nmid Q_\beta\),
\begin{equation}
 -v_\iota\!\left(H_{\nu,z}\,b\bar b\,
              \gamma_x\bar\gamma_x A_x^2\,
                   \rho_\beta^2 R_{\nu,z}(\beta)^{-1}\right)
       \le p_{\ell,\iota}+2S_\ell+2E_\ell .
 \label{eq:shortpairs859}
\end{equation}
\end{theorem}
\begin{proof}
For an inert prime \(e_0=1,\ Q=\ell^2\), and for the ramified
prime \(e_0=2,\ Q=\ell\). Summing the respective geometric
series gives \(\lambda_\ell=e_0^{-1}\sum_{a\in J_\ell}\ell^{-a}\).
The identity \(u-(u-v)_+=\min(u,v)\) for \(u,v\ge0\) therefore
gives
\[
  \lambda_\ell\sum_\tau\kappa_\tau-S_\ell
   =e_0^{-1}\sum_\tau\sum_{a\in J_\ell}
                    \min(\kappa_\tau/\ell^a,\kappa_{P_a\tau}).
\]
For a fixed \(\tau\), all terms with \(a>L\) have total at
most \(b_x/((\ell-1)\ell^L)\). At an exponent \(a\le L\)
the minimum is at most \(G\) if either endpoint has weight
\(G\). If both endpoints have the larger weight the minimum
is \(b_x/\ell^a\), giving exactly its contribution to
\(E_\ell\).

The derivative orders satisfy \(|d_\tau-\kappa_\tau|
\le W=C_{q_0,d_{\rm K}}x^{42}\) uniformly. In the main-pair
window of Theorem~\ref{thm:negativeinterval740}, the absolute
character weight is \(w_{j,\tau}-h_\tau\) with
\(w_{j,\tau}=\kappa_\tau+(2j+1)s^8,\ 0\le h_\tau\le4s\).
The derivative order further subtracts \(z\le s\), so the
claimed bound follows from \(s\le3x^5,\ m\le2\sqrt x\).
Also \(L\le2+\log_2 k\). The preceding terms and the error
from replacing the weights are at most
\[
             LG+\frac{b_x}{(\ell-1)\ell^L}
                         +\frac{W}{\ell-1}
\]
per embedding. The first and third terms are each
\(o(b_x/(k\ell))\) uniformly for \(\ell\le k\): for the first
the ratio is at most \((2+\log_2 k)xk^2/x_{1,x}\), and for
the third it is at most \(2Wk/b_x\). Here \(x_{1,x}\ge x^{10^6}\)
and \(b_x\ge x^{400}\). The middle term is at most
\(b_x/(k\ell)\). Since \(b_x/k^2\) tends to infinity, the
sum of the three terms is at most
\(\lfloor n_x/\ell\rfloor\le v_\ell(n_x!)\) at every such prime
on sufficiently distant intervals. Applying
Theorem~\ref{thm:scalartrace859} and discarding the
nonpositive ideal correction gives the claimed bound for
\(e_\iota\).

The members of \(\mathcal T\) all have the same unit class
modulo \(q_0\). A pair in \(E_\ell\) thus requires
\(-\ell^a\equiv1\bmod q_0\). At the ramified prime
\(\ell=d_{\rm K}\equiv q_0/2-1\bmod q_0\), whose square is
\(1\bmod q_0\) since \(4\mid q_0\). Neither \(1\) nor
\(q_0/2-1\) is \(-1\bmod q_0\). This proves the ramified
assertion (and \(d_{\rm K}\le k\) eventually for fixed \(K\)).

To see disjointness use \(2L<l_0\), valid here since
\(L\le2+\log_2 k\), whereas \(l_0\ge27s_2\) with
\(2^{s_2}\ge x^{11}\) in Lemma~\ref{lem:spacing781}.
That lemma forbids three visits to \(\mathcal T\) with positive
gaps and total gap less than \(l_0\), including repeated
visits, when the ratios are the indicated powers of \(\ell\)
up to sign. Two arrows from the same start at distinct
exponents would give such a triple. For arrows to the same
end with exponents \(a<a'\), start at the source of the
longer arrow and use times \(a'-a,a'\). Equal exponents
and equal ends imply equal starts because \(P_a\) is a
permutation. Two successive arrows of lengths \(a,a'\)
give visits at times \(a,a+a'\), including when a single
arrow is a loop and can be repeated. These alternatives
exclude all repeated vertices. Finally \(S_\ell,E_\ell\)
are independent of the coefficient embedding. Insert the
bounds at \(\iota,\iota\circ c\) in \eqref{eq:tablevaluation855}
to get \eqref{eq:shortpairs859}.
\end{proof}

The positive scalar \(A_x\) has
\(\log A_x\ll f b_x\log(2b_x)/k\), by \(n_x!\le n_x^{n_x}\).
Even summing this cost over \(mk\) critical pairs and
dividing by \(sB_x^{\rm det}\) gives
\(O((\log b_x)^{-2})\): here
\(B_x^{\rm det}=fxb_x(\log b_x)^3/(mk)\), \(k\le s\) and
\(m^2\ll x\). In particular the ramified-prime comparison
has the directional bound with an explicitly controlled
common multiplier. At a small inert prime the terms
\(E_\ell\) require further savings whenever these pairs occur.
\clearpage

\subsection{The logarithmic derivative at an inert prime}
\label{sec:inertgradient860}
Let \(p\) be inert in \(K\) and use the marked elliptic curve,
differential \(dz\), and even theta function \(\Theta=\Theta_{t_*}\)
of Theorem~\ref{thm:hurwitz62}. In the shorter family \(2,3\)
split in \(K\), so \(p\ge5\). Work over a completed local
integer ring \(R\), extended as needed, with residue field
\(\mathbf k\). Throughout this subsection reduction of an
integral series means coefficientwise reduction into \(\mathbf k\).
Take \(\lambda^{p^2-1}=p\) and put \(c_p=\lambda^{p+1}/p\).
Thus \(v_p(\lambda)=1/(p^2-1)\) and \(v_p(c_p)<0\).

\begin{lemma}[The reduced elliptic logarithmic derivative]
\leavevmode\label{lem:cartiergradient860}
The series \(c_p^{-1}(d/dZ)\log\Theta(\lambda Z)\) belongs
to \(R[[Z]]\). There are \(u_j\in\mathbf k^\times\), for all
\(j\ge0\), such that its reduction is
\[
                      \sum_{j\ge0}u_j Z^{p^{\,2j+1}} .
\]
\end{lemma}
\begin{proof}
Let \(T\) be the integral elliptic parameter of
Lemma~\ref{lem:unit41}, multiplied by a unit so that its
differential at zero is \(dz\). That lemma gives
\[
 Y(Z)=\lambda^{-1}T(\lambda Z)\in R[[Z]],\qquad
             \bar Y(Z)=\sum_{j\ge0}h_j Z^{p^{2j}},\quad
                    h_0=1,\quad h_j\ne0.
\]
In fact the linearized parameter makes the first nonlinear
term of \([p]\) in degree \(p^2\) a unit. Substitution of
\(\lambda Y\) in its exponential identity gives
\(\bar Y+\bar b_{p^2}\bar Y^{p^2}=Z\), as in that lemma.

Let \(g(z)=(\log\Theta)''(z)=-u_\Lambda-x(z+t_*)\), using
Theorem~\ref{thm:germ57}. Write
\(g\,dz=\sum_{i\ge0}a_iT^i\,dT\) at zero. These coefficients
lie in \(R\): \(u_\Lambda\) is integral here and the mark
stays away from the origin at \(p\). We claim that
\[
                a_{p-1}\in R^\times,\qquad
                a_{p^2-1}\in pR .
\]
For the first assertion work on the geometric special fiber.
If \(a_{p-1}\) vanished there, the terms through degree
\(p-1\) could be integrated in the parameter \(T\). For a
polynomial \(v\) this would give \(g-\partial_z v=O(T^p)\).
Consequently \((\partial_z^{p-1}g)(0)=(\partial_z^p v)(0)=0\)
by Lemma~\ref{lem:kappa62}. The left side is the nonzero
\(\kappa_{t_*}\) of that lemma.

For the second assertion work over \(S=\operatorname{Spec}(R/pR)\).
On a smooth relative curve the relative Cartier operator takes
\(\sum b_i t^i\,dt\) to
\(\sum_j b_{pj+p-1}(t^{(p)})^j\,dt^{(p)}\) in a parameter and
its Frobenius twist. This follows by writing forms over the
relative Frobenius in the basis \(1,t,\ldots,t^{p-1}\):
all powers except \(t^{p-1}dt\) have primitives in that basis,
and its class corresponds to \(dt^{(p)}\). The operation is
linear in base scalars and functorial for morphisms of smooth
curves. This is the relative Cartier isomorphism over an
arbitrary base of characteristic \(p\).
\href{https://numdam.org/articles/10.1007/BF02684688/}{N.~Katz,
\emph{Nilpotent connections and the monodromy theorem}, Theorem~7.2}.

The meromorphic differential \(g\,dz\) has only the double
pole at \(-t_*\) and has zero residue there. In characteristic
zero this follows from the even Laurent expansion of \(x\) in
the elliptic logarithm about its pole. On the good model the
pole has order at most two, and the residue in an integral
parameter is the same zero scalar, also after reduction
modulo \(pR\). Cartier therefore takes it to a global regular
differential on the Frobenius twist: neither its double-pole
term nor any regular term gives a pole by the displayed local
rule, and its residue is zero. Cartier kills all global regular
differentials of the elliptic curve over \(S\). Indeed a CM
endomorphism with tangent scalar \(a\) generating the quadratic
residue extension acts by \(a\) on the invariant differential
and by \(a^p\) on its twisted counterpart. Naturality and
\(S\)-linearity give \((a^p-a)C(dz)=0\), with unit multiplier.
The same argument on the twist uses the unit \(a^{p^2}-a^p\).
Regular differentials on an elliptic scheme are invariant.
Thus applying relative Cartier twice to \(g\,dz\) gives zero
over \(R/pR\). Its constant coefficient in the twice twisted
parameter is \(a_{p^2-1}\bmod pR\). This proves the claim with
the asserted precision for ramified coefficient rings as well.

The marked function is even, so \((\log\Theta)'(0)=0\).
Termwise integration in characteristic zero now gives
\[
 c_p^{-1}\frac{d}{dZ}\log\Theta(\lambda Z)
       =\sum_{m\ge1} a_{m-1}\frac{p}{m}\lambda^{m-p}Y(Z)^m .
\]
The scalar in this sum before \(a_{m-1}\) has valuation
\(1-v_p(m)+(m-p)/(p^2-1)\). It is positive when \(p\nmid m\),
and for \(v_p(m)=1\) it is positive except at \(m=p\), where
it is zero. For \(v_p(m)=2\) it is positive once \(m\ge2p^2\).
At \(m=p^2\) the additional factor \(a_{p^2-1}\in pR\)
makes it positive. For \(v_p(m)=r\ge3\), use \(m\ge p^r\):
the bound \(1-r+(p^r-p)/(p^2-1)\) is positive at \(r=3\)
and increases thereafter. Thus reduction leaves precisely
\(\bar a_{p-1}\bar Y^p\), with the asserted nonzero coefficients.
\end{proof}

\begin{theorem}[The odd Frobenius powers in matrix coordinates]
\leavevmode\label{thm:inertgradient860}
Use \(B=I\otimes C\), its theta line \(L\), and \(P_a\tau=-p^a\tau\)
from Sections~\ref{sec:twocopy858}--\ref{sec:shortpairs859}.
There is an integral formal frame of \(L\) at the origin
whose normalized scalar \(F(Z)\) in the marked chart has
\(F(0)=1\) and, for every label \(\sigma\),
\begin{equation}
 c_p^{-1}\partial_{Z_\sigma}\log F(\lambda Z)\in R[[Z]],
 \qquad
 \overline{c_p^{-1}\partial_{Z_\sigma}\log F(\lambda Z)}
    =\sum_{j\ge0}u_j
          \sum_{\tau:\,P_{2j+1}\tau=\sigma} Z_\tau^{p^{2j+1}} .
 \label{eq:inertgradient860}
\end{equation}
Moreover, every integral scalar formal function on \(B\)
has integral expansion after \(Z\mapsto\lambda Z\).
\end{theorem}
\begin{proof}
At \(p\) the different of \(E/K\) and the two old marks are
units, so \(I_p=\mathcal O_{E,p}\). The trace Hermitian form
\(H(a,b)=\operatorname{Tr}_{E/K}(a\bar b)\) on this lattice
is perfect since \(p\nmid q_1\). Its reduction is a
nondegenerate Hermitian space over \(\mathbb F_{p^2}\).
It has an orthogonal basis of vectors of nonzero norm:
if all norms vanished, polarization with multiples of a
vector by every residue scalar and nondegeneracy of the
finite-field trace pairing would force the form to vanish.
Split off a nonzero norm line and use induction.
Lift this basis to \(I\) and orthogonalize over \(K\).
All pivots are \(p\)-units by the chosen reductions.
Clear the resulting denominators away from \(p\) to obtain
orthogonal \(w_1,\ldots,w_f\in I\) forming a basis at \(p\).
Their norms \(n_i=H(w_i,w_i)\) are positive rational integers
prime to \(p\): \(I\subset\mathfrak b^{-1}\) and
\(\mathfrak b\bar{\mathfrak b}=\mathfrak D_{E/K}\) give
integrality by trace duality.

The homomorphism \(C^f\to B\) with columns \(w_i\) is an
isogeny of degree prime to \(p\), by the lattice index.
It is an isomorphism on formal neighborhoods of zero over
\(R\). In logarithm coordinates the inverse linear forms are
\[
               l_i(Z)=n_i^{-1}\sum_{\tau}\overline{\tau(w_i)}Z_\tau.
\]
Indeed the matrix \(V_{\tau i}=\tau(w_i)\) satisfies
\(V^*V=\operatorname{diag}(n_i)\). All these entries and
the inverse entries are integral at \(p\). On the product
lattice the trace form and the original period sign split as
the elliptic form and sign with multiplicities \(n_i\), by
the sign calculation of Lemma~\ref{lem:basis59}. Thus
\[
                    F(Z)=\prod_i\Theta(l_i(Z))^{n_i}
\]
is a normalized scalar of a formal frame for \(L\). In detail,
on \(C^f\) this product is the section of the external
product of the powers of the marked elliptic divisor bundles,
and the divisors miss the zero section on the special fiber.
It gives an integral frame there. The generic identification
of the pulled-back line with this product extends over the
local base up to a scalar, since discrepancies are vertical
divisors. Normalization by the value at zero removes that
scalar. Transport through the formal isomorphism gives the
frame on \(B\). An integral scalar germ pulls back to a series
in the \(f\) integral parameters of \(C\). After the indicated
rescaling these become \(T(\lambda l_i(Z))=\lambda Y(l_i(Z))\),
which are integral with zero constant term. This proves the
last assertion.

Apply the lemma to each factor and differentiate the product
logarithm. For \(a=2j+1\), the coefficient of \(u_j Z_\tau^{p^a}\)
in the reduced derivative at \(\sigma\) is the reduction of
\[
       \sum_i \overline{\sigma(w_i)}
                     \bigl(\overline{\tau(w_i)}/n_i\bigr)^{p^a}.
\]
Frobenius to the odd power \(a\) acts on \(K\) by conjugation
and on \(\mu_{q_1}\) by the \(p^a\)-power map. Composing with
conjugation is therefore the oriented label \(P_a\tau\).
This gives
\(\overline{\tau(w_i)}^{\,p^a}\equiv(P_a\tau)(w_i)\) in the
residue field, with the outer bar on the left denoting complex
conjugation before reduction. Since \(n_i\) is a rational
\(p\)-unit, the sum reduces to the \((P_a\tau,\sigma)\) entry
of \(V\operatorname{diag}(n_i^{-1})V^*\), which is the identity
matrix. This proves \eqref{eq:inertgradient860}.
\end{proof}

The formula retains the coefficient of each individual odd
power. In particular it distinguishes two ends having different
preimages under the same Frobenius power, as required for the
pairs in Theorem~\ref{thm:shortpairs859}.
\clearpage

\subsection{Coefficients from prescribed Frobenius pairs}
\label{sec:pairedtable861}
Fix an inert prime \(p\) on the same node, and let \(F,\lambda,c_p\)
be as in Theorem~\ref{thm:inertgradient860}. Choose a set \(A\)
of oriented embeddings. For each \(\tau\in A\) choose a positive
odd \(a_\tau\), and write \(t_\tau=P_{a_\tau}\tau\). Assume
that for each \(\tau\) the only solution to
\[
     P_a\sigma=t_\tau,\qquad \sigma\in A,\quad
                    a\text{ positive odd},\quad a\le a_\tau
\]
is \((\sigma,a)=(\tau,a_\tau)\). In particular the \(t_\tau\)
are distinct, since two coincident ends would violate this
condition for the longer exponent. The condition holds for all
embeddings with \(a_\tau=1\), by bijectivity of \(P_1\).
It also holds for the starts of the pairs occurring in
\(E_p\) of Theorem~\ref{thm:shortpairs859}. Indeed any competing
start with the specified end at a smaller or equal exponent
would itself give a term of \(E_p\), and the pairs there
are disjoint with unique exponents.

Attach nonnegative integers \(j_\tau\) to \(A\), put
\(J=\sum_{\tau\in A}j_\tau\), and form two multidegrees
\[
   r_\tau=p^{a_\tau}j_\tau\quad(\tau\in A),\qquad
   s_{t_\tau}=j_\tau\quad(\tau\in A),
\]
with the unspecified entries zero. For a multidegree \(u\)
write \(|u|=\sum_\tau u_\tau\), \(u!=\prod_\tau u_\tau!\).
In particular \(|s|=J\).

\begin{lemma}[An exact mixed coefficient of the actual frame]
\leavevmode\label{lem:paircoefficient861}
Use \(\delta_1=1,\ldots,\delta_6\in\mathbb Z_p\) of
Lemma~\ref{lem:diagonalframe858}. Put
\[
      D(U,V)=F(U+V)\prod_{i=2}^6 F(\delta_i U),\qquad
                      C_{r,s}=[U^rV^s]D(U,V).
\]
Then \(C_{r,s}\ne0\) and
\begin{equation}
     v_\iota(C_{r,s})
        = Jv_p(c_p)-v_p(s!)-(|r|+J)/(p^2-1).
 \label{eq:paircoefficient861}
\end{equation}
\end{lemma}
\begin{proof}
Use \(D_\lambda=D(\lambda U,\lambda V)\), a unit in the
power series ring over the local field, with constant one.
Its restriction \(D_\lambda(U,0)\) is integral and has
constant one. Indeed Lemma~\ref{lem:diagonalframe858} gives
an integral function in formal group parameters before
rescaling, and Theorem~\ref{thm:inertgradient860} makes
that function integral in the rescaled logarithm variables.
The cancellation lemma applies also to the formal frame
used here. Replacing an algebraic frame by an integral formal
one multiplies its normalized value by a formal unit of
constant one. Pullback of this unit by any \(\delta_i\) is
integral, since integer multiplication is integral and its
finite jets converge coefficientwise when the integers tend
to \(\delta_i\) in \(\mathbb Z_p\). This last convergence
follows by linear multiplication in the logarithm coordinates
and a fixed change of parameters through each jet order.

For a multi-index \(v\) put
\[
       L_v=c_p^{-|v|}(\partial_V^v D_\lambda)/D_\lambda .
\]
Starting with \(L_0=1\), the product rule gives
\[
 L_{v+\mathbf e_t}=c_p^{-1}\partial_{V_t}L_v+
           L_v\,c_p^{-1}\partial_{V_t}\log D_\lambda .
\]
The last logarithmic derivative is just the one of
\(F(\lambda(U+V))\), as the other factors are constant in
\(V\). It becomes integral after multiplication by \(c_p^{-1}\)
by \eqref{eq:inertgradient860}. Since \(v_p(c_p^{-1})>0\),
induction proves integrality of \(L_v\), and the first term
of this recurrence disappears after reduction.
Consequently the reduction of \(L_s(U,0)\) is
\[
   \prod_{\tau\in A}\left(
    \sum_{\substack{a\ge1\ {\rm odd},\ \sigma\\P_a\sigma=t_\tau}}
              u_{(a-1)/2} U_\sigma^{p^a}\right)^{j_\tau}.
\]

Set \(U_\sigma=0\) for \(\sigma\notin A\). In each factor
before exponentiation the unique term of smallest total degree
is \(u_{(a_\tau-1)/2}U_\tau^{p^{a_\tau}}\). Their product
through the indicated powers has lowest term exactly in degree
\(r\), with nonzero coefficient. Multiplication by the
reduction of \(D_\lambda(U,0)\) preserves this coefficient,
since its constant is one. Thus
\([U^r](L_s D_\lambda)(U,0)\) is a unit. Coefficients supported
on \(A\) are unchanged by the specialization outside \(A\).
By definition this last coefficient is
\(c_p^{-J}s!\lambda^{|r|+J} C_{r,s}\).
This proves the identity.
\end{proof}

\begin{theorem}[Valuation savings in the integral twelve-row tests]
\leavevmode\label{thm:pairedtable861}
Let \(d=(d_\tau)_\tau\) be the actual derivative orders and
suppose that \(\eta=d-r-s\) is componentwise nonnegative.
Write \(h=|d|\) and
\(\lambda_p=1/(p-1)-1/(p^2-1)\), as in
Theorem~\ref{thm:scalartrace859}. At the fixed coefficient place
the derivative table satisfies both bounds
\begin{align}
 e_\iota-v_\iota(b\gamma_x)&\le v_p(d!)-J,
                 \label{eq:firstpair861}\\
 e_\iota-v_\iota(b\gamma_x)&\le
        \lambda_p h-J+v_p\!\left(\frac{d!}{r!s!\eta!}\right).
                 \label{eq:secondpair861}
\end{align}
\end{theorem}
\begin{proof}
First record two consequences of \eqref{eq:paircoefficient861}.
Since \(v_p(c_p)=1/(p-1)-1\) and \(|r|\ge pJ\),
one has \(v_\iota(C_{r,s})\le-J\). More precisely
\(v_\iota(C_{r,s})=-J-v_p(s!)-(|r|-pJ)/(p^2-1)\).
Also \(v_p(r!)\le |r|/(p-1)\), so
\[
           v_\iota(r!s! C_{r,s})\le
                             \lambda_p(|r|+J)-J.
\]

For each integral formal function \(H(W)\) of \(B\), we
claim that the actual table satisfies
\[
 e_\iota\le v_\iota(b\gamma_x)+
           v_\iota\bigl(d!\,C_{r,s}[W^\eta]H\bigr)
\]
whenever the last coefficient is nonzero. In the first six
factors of \(B^{12}\) make the logarithm substitutions
\(U+V,\delta_2 U,\ldots,\delta_6 U\); in the last six
make the substitutions \(\delta_1W,\ldots,\delta_6W\).
Write \(J_0(W)=\prod_iF(\delta_i W)\), an integral unit in
the formal group parameters by the cancellation lemma.
In the external product frame prescribe the integral jet of
\(H/J_0\) on the seventh factor, pulled back by projection
to that factor. Lift through order \(h\) to an integral
global section \(S\) as in Theorem~\ref{thm:integralleading857}.
If \(\lambda_R\) denotes the generator value of the integral
frame, the substituted marked expansion of \(S/\lambda_R\)
through that degree is \(D(U,V)H(W)\). The first six frames
supply \(D\) and the last six restore \(J_0\).
Write \(S=\sum_y c_yG_y\). Multiplying by \(M^{12f}g\) makes
each coefficient \(c_y/\lambda_R\) integral by
Lemma~\ref{lem:residuecoeff857}, with the \(g\) in the definition
of \(\gamma_x\). Substitution in a monomial uses only integer
binomial coefficients and powers of the integral \(\delta_i\).
All monomials contributing to \(U^rV^sW^\eta\) have the same
per-embedding degrees \(d_\tau\) in the original variables.
Multiplication by \(b\,d!\) yields an integral combination
of the table entries. As \(M^{12f}g\) divides \(\gamma_x\)
in the sense that their quotient is the integer \(t_E^{12f}\),
the claim follows.

For the first bound take \(H\) with leading homogeneous part
\(W^\eta\) in marked coordinates. Such an integral function
exists here: \(I_p=\mathcal O_{E,p}\) makes each \(dW_\tau\)
an integral cotangent vector by the differential computation in
Theorem~\ref{thm:integralleading857}. Lift these vectors to the
ideal of zero in the smooth formal group and multiply the lifts
to the prescribed powers. The coefficient at \(\eta\) is one.
The first bound follows. For the second bound use the scalar
trace construction of Theorem~\ref{thm:scalartrace859} with
total degree \(|\eta|\) and the trace morphism \(w=1\).
The trace of \(I\) is integral and all coefficients of the
linear trace coordinate are one. The construction therefore
gives \(v_\iota(\eta![W^\eta]H)\le\lambda_p|\eta|\).
Combining this with the inequality for \(r!s!C_{r,s}\) proves
the second bound, including \(\eta=0\).
\end{proof}
\clearpage

\section{Integral positions in the terminal tests}\label{nav:g-VIII-4}

\subsection{Original-level projection of terminal covector translates}\label{nav:VIII-13}
\label{sec:terminal199}

Fix one of the fields \(E/F\), \(F=E^+\), and its terminal prime
\(T=\varpi\bar\varpi\) from Section~\ref{sec:middle131}.
There are \(f=[F:\QQ]\) places \(v\mid T\), all with
\(F_v=\QQ_T\). Put \(O=\ZZ_T\), \(q=T\), and \(r=2s+k+1\).
At these places use the first split factor of that section. The
row group is \(\GL_r(F_v)\). The compact column group
\(\mathcal U_v=\GL_k(O)\) acts on the middle \(k\) coordinates
and fixes all other coordinates. For a nonzero residue covector
\(\ell\in X_v=(\mathbb F_q^k)^\vee\setminus\{0\}\), put
\[
 \Lambda_\ell=O^s\oplus
       \{(w,z)\in O^k\oplus O:z=\ell(\bar w)\bmod T\}
                  \oplus O^s .
\]
Each \(\Lambda_\ell\) is a sublattice of index \(q\) in the
same lattice \(\Lambda_*=O^r\). The compact row group \(K_v\)
in the original period is the stabilizer of
\(\Lambda_{\ell_{0,v}}\), for the prescribed \(\ell_{0,v}\).
The paired lattice in the other split factor is its dual.

Write \(R\) for right translation. If the row transforms by
a smooth unitary character \(\kappa_v\) of \(K_v\), use the
following convention for translations at \(v\). In a basis of
its lattice \(K_v=\GL_r(O)\). Every character is trivial on
\(\mathrm{SL}_r(O)\). Row reduction over the local ring gives
generation of that group by elementary matrices, and each
\(1+tE_{ij}\) is the commutator of \(1+tE_{il}\) and
\(1+E_{lj}\) for a third index \(l\). Thus \(\kappa_v\)
factors through the determinant. Extend this unit character
to \(F_v^\times\) by assigning \(T\) the value one and write
\(\widetilde R(g)=\kappa_v(\det g)^{-1}R(g)\), using the
same letter for the extended character. Under this action
the original level space is fixed by \(K_v\). All factors
introduced on compact translations are roots of unity. We use
the same twist on the column action through its block embedding.
For trivial level character, \(\widetilde R=R\).

\begin{lemma}[A uniform bound for the neighboring lattices]
\label{lem:coefficient199}
In lattice coordinates put \(d=\operatorname{diag}(T,T^{-1},1,\ldots,1)\)
and let \(e_v\) be orthogonal projection to the \(K_v\)-fixed
vectors for \(\widetilde R\). On the cuspidal row space the operator
\(M_v=e_v\widetilde R(d)e_v\) is self-adjoint and
\begin{equation}
        \|M_v\|\le c_q,\qquad c_q=\frac{3q-1}{q(q+1)} .
                                                \label{eq:constant199}
\end{equation}
\end{lemma}
\begin{proof}
We first exclude vectors invariant under \(\mathrm{SL}_r(F_v)\).
The rational row is isotropic. Let \(N_0\) be the unipotent
radical of the proper rational parabolic stabilizing the first
isotropic line. Its rational points are dense in its adeles
away from \(v\), including the real places. Here is a proof of
the approximation used in this assertion. For the additive
group, a finite list of nonarchimedean conditions away from \(v\)
can be met by the Chinese remainder theorem after clearing
denominators, with integrality at the other unspecified finite
places. A fractional ideal \(\mathfrak a\) preserves all these
conditions by addition, as does \(\mathfrak a\mathfrak p_v^{-j}\)
for any \(j\ge0\). The union of these ideals is dense at infinity.
Indeed, a character of the real torus \(F_\infty/\mathfrak a\)
annihilating the union would have frequency in every trace-dual
ideal \(\mathfrak d_F^{-1}\mathfrak a^{-1}\mathfrak p_v^j\).
Their intersection is zero by the valuation at \(v\).
Characters separate closed subgroups of a torus, proving density.
The finite exponential and logarithm polynomials identify \(N_0\)
with its Lie algebra as varieties in characteristic zero. For
unitary matrices the logarithm is skew adjoint since taking the
adjoint takes the matrix to its inverse. The polynomials preserve
integrality at almost all places. Additive approximation in a
rational basis of the Lie algebra proves the claim for \(N_0\).

Suppose a smooth cusp vector \(A\) is invariant under
\(\mathrm{SL}_r(F_v)\). For \(g\in G(\mathbb A_F)\) and
\(n\in N_0(\mathbb A_F)\), approximate \(n\) by products
\(\gamma u_v\) with \(\gamma\in N_0(F)\), \(u_v\in N_0(F_v)\),
correcting the \(v\)-component exactly. Then
\(A(\gamma u_vg)=A(g)\) by automorphy and by
\(g_v^{-1}u_vg_v\in\mathrm{SL}_r(F_v)\). Continuity gives
\(A(ng)=A(g)\). Its constant term along \(N_0\) is therefore
\(A(g)\) times a positive covolume and vanishes by cuspidality.
This also excludes invariant vectors in the cuspidal Hilbert space.
The invariant subspace is stable under the whole right action
by normality, so smoothing on the group and taking compact
types gives a dense set of smooth cuspidal vectors there.
The determinant twist is trivial on \(\mathrm{SL}_r(F_v)\).

In any unitary representation of \(\GL_r(F_v)\), a vector
fixed by a full elementary root subgroup is fixed by
\(\mathrm{SL}_r(F_v)\). To see this, write \(x(b),y(t)\)
for the upper and lower elementary matrices in the two
associated coordinates. For \(a\in F_v^\times\), take
\(b=(a-1)/t\), \(t\ne0\). Then
\[
           x(b)y(t)x(-b/a)=
                    \begin{pmatrix}a&0\\t&a^{-1}\end{pmatrix}.
\]
On an \(x\)-fixed vector the norm of the change on the left is
the norm of its change under \(y(t)\), which tends to zero
with \(t\). Hence every such diagonal \((a,a^{-1})\) fixes
the vector. For \(a=T\), conjugation by powers of this diagonal
or its inverse contracts each root subgroup having at least one
index in this pair to the identity. Invariance and continuity give
invariance under all those subgroups. Taking their commutators
gives the remaining elementary subgroups, and row reduction
generates \(\mathrm{SL}_r(F_v)\).

Let \(U\) consist of the two entries in rows \(1,2\) of column
\(3\) added to the identity. In the spectral resolution of the
unitary action of \(U\simeq F_v^2\), the projection at the zero
character vanishes by the preceding argument. We use the spectral
theorem for unitary representations of a locally compact abelian
group. In these coordinates the projection for a dual ball
\(\xi_0+T^jO^2\) is also the Haar probability average on
\(T^{-j}O^2\) against the conjugate character of \(\xi_0\),
using an additive character of conductor \(O\).
Project each nonzero frequency to its direction in
\(\mathbb P^1(F_v)\). For every \(K_v\)-fixed unit vector, the
squared lengths of its spectral projections on directions give
the same probability \(m\). Indeed the upper \(2\)-by-\(2\)
copy of \(\mathrm{SL}_2(O)\) preserves the vector and acts
transitively on directions of primitive pairs. A primitive pair
extends to a determinant-one basis over \(O\). Transitivity
and compact Haar measure give uniqueness of the probability.

Choose the affine coordinate on the dual projective line so that
\(d\) acts by \(t\mapsto T^2t\). Write the three parts of the line as
\(A=T^2O\), \(B=TO\setminus T^2O\), and \(C=\mathbb P^1(F_v)
\setminus TO\). Their masses for \(m\) and \(d_*m\) are
\[
 \begin{array}{c|ccc}
       & A&B&C\\ \hline
 m     &1/[q(q+1)]&(q-1)/[q(q+1)]&q/(q+1)\\
 d_*m  &q/(q+1)&(q-1)/[q(q+1)]&1/[q(q+1)] .
 \end{array}
\]
Reduction modulo \(T\) gives \(q+1\) equal direction classes,
each splitting into \(q\) equal classes modulo \(T^2\), as is
seen by lifting a primitive projective pair in either affine
chart. In particular \(m(O)=q/(q+1)\). Scaling and inversion
give the second row. For two fixed unit vectors, project the
first and the \(d\)-translate of the second onto the three sets.
Cauchy--Schwarz on each set bounds their inner product in
absolute value by the sum of the three square roots of the
products of corresponding masses. The sum is \(c_q\).
Finally \(d^{-1}\) is conjugate to \(d\) by interchanging the
first two basis vectors in \(K_v\). The compressed operator
is therefore self-adjoint. This proves \eqref{eq:constant199}.
The general root-subgroup bound is due to H.~Oh \cite{XVI-XIII-IV-oh199}.
\end{proof}

Write \(X=\prod_{v\mid T}X_v\), \(\ell_0=(\ell_{0,v})_v\),
and give \(\CC^X\) the inner product in which the point masses
are orthonormal. Use the action \(h\cdot\ell=\ell h^{-1}\).
Choose \(h_{\ell,v}\in\mathcal U_v\) sending \(\ell_{0,v}\)
to \(\ell\), and take \(h_{\ell_{0,v},v}=1\).
For an original-level cusp vector \(A\) and \(a\in\CC^X\) write
\[
       D(a,A)=\sum_{\ell\in X} a(\ell)\,
                  \widetilde R\big((h_{\ell_v,v})_v\big)A,
       \qquad e=\prod_{v\mid T}e_v .
\]
The definition is independent of the chosen lifts because the
stabilizer of \(\ell_{0,v}\) preserves \(\Lambda_{\ell_{0,v}}\).
Let \(U_X\subset\CC^X\) consist of arrays whose sum over each
single coordinate is zero with the others fixed. It contains
the refined columns whose sum on each local scalar orbit is zero.

\begin{theorem}[The Gram matrix and its original-level projection]
\label{thm:hec199}
On original-level cusp vectors put \(L_T=\prod_{v\mid T}(1-M_v)\).
For \(a,b\in U_X\) and such \(A,A'\), with inner products
conjugate-linear in the first variable, one has
\begin{align}
 \langle D(a,A),D(b,A')\rangle
     &=\langle a,b\rangle\,\langle A,L_T A'\rangle ,
                                                 \label{eq:gram199}\\
 eD(a,A)&=a(\ell_0)L_T A,                       \label{eq:compress199}\\
 (1-c_q)^f\,1 &\le L_T\le (1+c_q)^f\,1 .       \label{eq:lower199}
\end{align}
Set \(p_r=(q^r-1)/(q-1)\), \(n_r=p_r(p_r-1)\).
The operator \(n_r^f L_T\) on this level is a sum of
finite translations with algebraic integer coefficients.
\end{theorem}
\begin{proof}
Each lattice is the inverse image of its hyperplane in
\(\Lambda_*/T\Lambda_*\). For two distinct covectors the
hyperplanes are distinct, even if one covector is a scalar
multiple of the other, since the coefficient of \(z\) is fixed.
Their sum is \(\Lambda_*\), and their intersection has index
\(q\) in each. The elementary divisors of the second lattice
relative to the first therefore have one \(T\), one \(T^{-1}\),
and \(r-2\) units. Indeed the positive exponents measure the
quotient of the first lattice by the intersection and the
negative ones measure the corresponding quotient for the
second. Smith reduction over \(O\) then identifies the
relative double coset with \(K_vdK_v\), independently of
the distinct covectors. For equal covectors it is \(K_v\).

The compressed translation between two lifts is consequently
\(1\) on the diagonal and \(M_v\) off it. With \(J_v\) the
matrix of ones on \(X_v\), the single-place Gram matrix is
\(1_{X_v}\otimes(1-M_v)+J_v\otimes M_v\).
At different places translations and projections commute.
The matrix on the product is the product of these factors
acting on the corresponding indices and the common row space.
Each \(J_v\) kills \(U_X\), giving \eqref{eq:gram199}.
Its row at \(\ell_0\) evaluates the \(K_v\)-compression at
all terminal places. On \(U_X\) this is \eqref{eq:compress199}.
The commuting \(M_v\) are self-adjoint contractions with the
bound in Lemma~\ref{lem:coefficient199}. The spectral bounds
on their product give \eqref{eq:lower199}.

To count the right \(K_v\)-cosets in \(K_vdK_v\), start at
the original lattice and adjoin an overlattice of index \(q\).
There are \(p_r\) choices, the lines of its inverse-\(T\)
quotient by itself. In each overlattice take a sublattice of
index \(q\) different from the starting lattice. There are
\(p_r-1\) choices, by counting hyperplanes modulo \(T\).
The intermediate lattice is recovered as the sum of the first
and last ones. The elementary-divisor computation shows that
this counts the entire double coset exactly once. On fixed
vectors \(M_v\) is the average over these \(n_r\) translates
with equal weights, by transitivity of left \(K_v\) on the
right cosets. Multiplication by \(n_r\) clears the denominator.
The character twists on the representatives are roots of unity.
Multiplying the resulting integral sums over the places proves
the last assertion.
\end{proof}

\begin{corollary}[Integral character sums in the first theta test]
\label{cor:theta199}
Let \(Z=\prod_{v\mid T}\mathbb F_q^\times\) be the scalar
subgroup of \(\prod_{v\mid T}\GL_k(\mathbb F_q)\).
Choose a character \(\lambda=\prod_v\lambda_v\) nontrivial in
each coordinate. In the permutation action on \(X\) form
\begin{equation}
        w_\lambda=\sum_{z\in Z}\overline{\lambda(z)}
                      \,\delta_{z\cdot\ell_0}.
                                                  \label{eq:word199}
\end{equation}
Its coordinates are algebraic integers. For the iterated
complex theta operator \(\mathcal T=\mathcal T_{\mathbf a,0}\)
of Theorem~\ref{thm:norm177}, using its fixed first-edge test
\(\phi_*\) at each \(v\mid T\), one has
\begin{equation}
 \begin{split}
  \mathcal T D(w_\lambda,A)&=\mathcal T L_T A,\\
 e^{-\Delta_x}(1-c_T)^{2f}\|A\|^2
       &\le \frac{\|\mathcal T D(w_\lambda,A)\|^2}{C_xP(\mathbf a)}
        \le e^{\Delta_x}(1+c_T)^{2f}\|A\|^2 .
 \end{split}                                       \label{eq:theta199}
\end{equation}
Here \(A\) belongs to the original-level space of that theorem.
On the common intervals, \(\|L_T-1\|=O(f/T)\) and
\(\log(n_r^f)=O(fr\log T)=o(b/x^B)\) for each fixed \(B>0\),
where \(b=\exp(x^\eps)\).
\end{corollary}
\begin{proof}
The scalar action on a nonzero covector is free. Its orbit
sum has coefficient one at \(\ell_0\) and sum zero along
each coordinate, by nontriviality of that character.
Thus \eqref{eq:compress199} gives \(eD(w_\lambda,A)=L_TA\).
The first theta kernel is formed from the vector \(\phi_*\)
which transforms by \(\kappa_v\) in its source variable,
by Theorem~\ref{thm:tests165}. In its pairing with the source
row, translate the integration variable by \(k\in K_v\).
The two compact characters cancel on the matched isotypic
space. Averaging the source with the projector to that
character leaves the pairing unchanged because the kernel belongs
to that character space in the integrated variable.
This argument also defines the first lift on the finite sum
of smooth translates in \(D\). All subsequent tests stay fixed,
so \(\mathcal TD=\mathcal T eD\). The operator \(L_T\) acts at
the terminal places, preserves their original level, and commutes
with all other level and real type conditions on \(A\).
Apply \eqref{eq:common177} at depth zero, where \(I=1\), to
\(L_TA\), and use \eqref{eq:lower199} to bound its squared norm.

The ranks \(f,r\) and \(\log T\) are bounded by fixed powers
of \(x\). Also \(\log T\ge(16-o(1))m x^\eps\) on these
intervals, with \(m\to\infty\), by
Lemma \ref{lem:size97} and Theorem \ref{thm:common111}.
In particular \(fc_T\to0\), and \((1+c_T)^f-1=O(f/T)\).
The lower endpoint in \eqref{eq:lower199} differs from one
by at most \(fc_T\). Finally \(p_r\le 2T^{r-1}\), so
\(f\log n_r=O(fr\log T)\), which is polynomial in \(x\)
and has the asserted comparison with \(b/x^B\).
\end{proof}

\begin{corollary}[Mass of the full zero-sum projection]
\label{cor:whole199}
At each terminal place let \(P_{1,v}\) be averaging under the
twisted action of the residue scalars, using their multiplicative
unit lifts. Put
\[
        P_{\rm nz}=\prod_{v\mid T}(1-P_{1,v}),\qquad
        \delta_T=\left(\frac{T-2}{T-1}\right)^f .
\]
On the original-level cuspidal space one has
\begin{align}
        eP_{\rm nz} A&=\delta_T L_T A, \notag\\
        \|(1-P_{\rm nz})A\|^2
            &\le \frac{5f}{T}\|A\|^2 .               \label{eq:all199}
\end{align}
For the individual sum \(D_\lambda=D(w_\lambda,A)\) and
\(A\ne0\), the fraction of squared norm retained at the
original level lies in the interval
\begin{equation}
 \frac{(1-c_T)^f}{(T-1)^f}\ \le\
       \frac{\|eD_\lambda\|^2}{\|D_\lambda\|^2}
         \ \le\ \frac{(1+c_T)^f}{(T-1)^f}.             \label{eq:frac199}
\end{equation}
The operator \((T-1)^fP_{\rm nz}\) is integral over the ring
of finite character values. Its coefficients on the scalar
translations in the twisted action are ordinary integers.
\end{corollary}
\begin{proof}
For odd \(T\), the \(T-1\) nonzero residues have unique lifts
solving \(z^{T-1}=1\) in \(\ZZ_T^\times\), by Hensel's lemma.
They form a subgroup. Its finite character projections are
pairwise orthogonal, so \(P_{\rm nz}\) is an orthogonal projection.
On \(A\) it is \(|Z|^{-1}\sum_\lambda D(w_\lambda,A)\),
summed over the \((T-2)^f\) characters nontrivial at every
place. Equations~\eqref{eq:word199} and \eqref{eq:compress199}
give its asserted compression. Since \(eA=A\),
\[
 \|(1-P_{\rm nz})A\|^2=\|A\|^2-\delta_T\langle A,L_TA\rangle .
\]
By \eqref{eq:lower199} the relative loss is at most
\(1-[(1-1/(T-1))(1-c_T)]^f\).
This is at most \(f(1/(T-1)+c_T)\le5f/T\) for \(T\ge3\).

For one character, the squared norm before projection is
\(|Z|\langle A,L_TA\rangle\) by \eqref{eq:gram199}.
After projection it is \(\langle A,L_T^2A\rangle\).
Their quotient gives \eqref{eq:frac199} by the spectral bounds
for \(L_T\). Finally \((T-1)^fP_{\rm nz}\) is the product of
\((T-1)\,1-\sum_{z_v}\widetilde R(z_v)\), proving integrality.
The same fixed first-edge test sends \(P_{\rm nz}A\) to
\(\delta_T\mathcal T L_TA\). Thus the estimate for the combined
characters also applies before the remaining theta integrations.
\end{proof}

We next apply the same finite symmetry to the Bessel restriction
itself. On a smooth row let
\[
    \mathcal F_\beta A(h)
       =\int_{N(F)\backslash N(\mathbb A_F)}
                    A(nh)\overline{\chi_\beta(n)}\,dn,\qquad
          \beta=\varpi .
\]
One may follow this map by orthogonal projection to any of the
chosen real column types in \(L^2([H])\), and write \(\mathcal F\)
for the resulting map or for \(\mathcal F_\beta\) itself. Both
quotients in \eqref{eq:period94} are compact, so on these smooth
vectors \(\mathcal F A\) has finite norm. Its pairing with a column
in the specified type is precisely the period in that equation.
Let \(p_{\lambda,H}\) be the projection to scalar character
\(\lambda\) on the terminal finite orbit of \(\mathcal F A\),
using the twisted column action. In \(\CC^X\) denote the
corresponding character space by \(V_\lambda\).

\begin{theorem}[One Bessel vector for each terminal scalar character]
\label{thm:bessel199}
For \(a,b\in V_\lambda\) and smooth original-level cusp rows
\(A,A'\), there is an exact identity
\begin{equation}
 \begin{split}
 \langle\mathcal F D(a,A),\mathcal F D(b,A')\rangle
    ={}& (T-1)^f\langle a,b\rangle\\[-2pt]
       &{}\times\langle p_{\lambda,H}\mathcal F A,\,
                         p_{\lambda,H}\mathcal F A'\rangle .
 \end{split}                                           \label{eq:bessel199}
\end{equation}
There are \(((T^k-1)/(T-1))^f\) independent directions in
\(V_\lambda\). In each tuple of residue lines the vector
\eqref{eq:word199} transported to those lines is an integral
generator of the corresponding one-dimensional character space
over the ring of integers containing the character values.
\end{theorem}
\begin{proof}
Right translation on \(H\) commutes with taking the unipotent
coefficient, by substitution in the displayed definition of
\(\mathcal F_\beta\). The real projection and the compatible
determinant twists preserve this equivariance. The orbit map
\(\delta_\ell\mapsto\mathcal F D(\delta_\ell,A)\) is therefore
equivariant for the product of the compact column groups.
The principal congruence kernel in each group fixes the base
vector since it fixes the initial covector. It is normal in
the compact column group, so it fixes the entire orbit.
Consequently this is an equivariant map for the finite product
group \(\prod_{v\mid T}\GL_k(\mathbb F_T)\).

We verify irreducibility on a single character space for one
factor. The functions of that character on a fixed residue line
form a line supported on its nonzero vectors. Distinct residue
lines give orthogonal lines in the function space. For a pair
of different residue lines choose an element of \(\GL_k(\mathbb F_T)\)
acting by one on the first and by any specified scalar on the
second. This is possible in a basis of the dual space extending
representatives of the two lines. Since \(\lambda_v\ne1\), one
may choose a scalar whose action on the second character line
is nontrivial. An equivariant endomorphism has zero entry between
these lines. Its diagonal entries agree by transitivity on residue
lines. Hence its commutant consists of the scalars, proving
irreducibility by orthogonal projection. The same proof works on
tuples. For two different tuples use a differing place to kill
the off-diagonal entry, and use product transitivity for the
diagonal. It also gives the stated dimension since there are
\((T^k-1)/(T-1)\) residue lines at each place.

The Gram form of the orbit map on \(V_\lambda\), using \(A\) in
its first variable and \(A'\) in its second, thus multiplies the
ordinary inner product by a single scalar. To compute it, insert
\(p_\lambda\delta_{\ell_0}=w_\lambda/|Z|\) in both variables.
Its squared norm is \(1/|Z|\). Its image for \(A\) is precisely
\(p_{\lambda,H}\mathcal F A\), by equivariance and the character
average. This gives \eqref{eq:bessel199}, including its
normalizing factor \(|Z|=(T-1)^f\). On each tuple of lines all
coordinates of an eigenvector are determined by one coordinate
times roots of unity. Conversely arbitrary integral values of
that coordinate give integral eigenvectors, proving the last
assertion.
\end{proof}

For the actual integral Fourier sum, \eqref{eq:bessel199} gives
\(\mathcal F D(w_\lambda,A)=|Z|p_{\lambda,H}\mathcal F A\)
directly by equivariance. The lower bound \eqref{eq:theta199}
is relative to the Petersson norm of \(A\). In particular,
putting \(A=S_0\xi\) leaves its lower bound proportional to
\(\|S_0\xi\|^2\), the original projected Bessel Gram form.
For the arithmetic comparison one still needs a lower estimate
on the designated Bessel pairings, with their integral
normalizations. A choice of a nonzero row in the complex
theta correspondence does not supply that estimate.

\subsection{The original row projection on integral terminal characters}\label{nav:VIII-28}
Keep the field \(E\), the interval and the prime \(T\) of Theorem~\ref{thm:common111}. In this paragraph \(r\) is the rank of the indefinite row,
\(k\) the rank of the definite column and \(f=[E^+:\QQ]\).
At each \(v\mid T\) use the original row lattice
\(\Lambda_{\ell_{0,v}}\) of
Section \ref{nav:VII-00}.
Its middle block consists of pairs \((w,z)\) satisfying
\(z\equiv\ell_{0,v}(\bar w)\pmod T\).
For \(h\in\mathrm{GL}_k(\ZZ_T)\), the translated lattice uses the
residue covector \(\ell_{0,v}h^{-1}\).

\begin{theorem}[A terminal character sum surviving the original row projection]
\label{thm:original200}
Let \(A\) be a cusp row of the specified real type and original
finite level. Write \(e\) for projection to its prescribed
compact type at all \(v\mid T\). Let
\(Z=\prod_{v\mid T}\mathbb F_T^\times\), acting by scalar
matrices on the column coordinates. For each \(z\in Z\), lift
these scalars to units and write \(\widetilde R(z)A\) for the
corresponding row translate, with its compact determinant
character removed. For any character \(\lambda\) of \(Z\)
nontrivial at each place put
\[
      A_\lambda=\sum_{z\in Z}\overline{\lambda(z)}
                                \widetilde R(z)A .
\]
There is one positive operator \(L_T\) on the original-level
cuspidal row, independent of \(\lambda\), such that
\begin{equation}
        e A_\lambda=L_TA,\qquad
        \|L_T-1\|\ll f/T .                         \label{eq:lift200}
\end{equation}
More precisely, if \(c_T=(3T-1)/(T(T+1))\), its spectrum
lies in \([(1-c_T)^f,(1+c_T)^f]\).
Writing \(p_r=(T^r-1)/(T-1)\), multiplication of \(L_T\)
by \((p_r(p_r-1))^f\) gives an integral linear combination
of translations over the ring of integers containing the
finite character values. The logarithm of this denominator
is \(o(b/x^B)\) for every fixed \(B>0\).
\end{theorem}
\begin{proof}
At one place, two distinct covectors give distinct hyperplanes
in the reduction of the same rank-\(r\) lattice. Their
intersection has index \(T\) in each lifted lattice. The
relative elementary divisors are \(T,T^{-1},1,\ldots,1\).
Hence every translate with covector different from the base
has the same original-level compression \(M_v\).
The unitary matrix coefficient estimate for this double coset
on cuspidal rows gives \(\|M_v\|\le c_T\). Its proof uses the
three sets of directions of an abelian unipotent subgroup in
Lemma \ref{lem:coefficient199}.
There are \(p_r(p_r-1)\) right cosets in the double coset.
The scalar element one gives the identity compression, and the
sum of the other character coefficients is \(-1\). Thus the
single-place projection is \(1-M_v\). Multiplication over the
places gives \(L_T=\prod_{v\mid T}(1-M_v)\), with the
claimed spectrum and denominator by
Theorem \ref{thm:hec199}.
Lemma~\ref{lem:size97} gives \(f/T\to0\) and bounds \(f,r,\log T\)
by fixed powers of \(x\). The assertions about the norm
and \(b=\exp(x^\eps)\) follow.
\end{proof}

\begin{corollary}[Two period norms for the same terminal character]
\label{cor:original200}
Let \(\mathcal T\) be the iterated complex theta operator in
Theorem \ref{thm:norm177}
with its original first-edge test at \(T\), and let \(A\)
satisfy the other real and finite type conditions of that
theorem. Write \(\mathcal F A\) for the Bessel coefficient of
frequency \(\varpi\) restricted to the definite column group,
optionally projected to the specified real column type. Then
\begin{align}
  \frac{\|\mathcal T A_\lambda\|^2}{C_xP(\mathbf a)}
       &\ge e^{-\Delta_x}(1-c_T)^{2f}\|A\|^2 ,
                                                    \label{eq:theta200}\\
  \|\mathcal F A_\lambda\|^2
       &=(T-1)^{2f}\|p_{\lambda,H}\mathcal F A\|^2 .
                                                    \label{eq:bessel200}
\end{align}
Here \(p_{\lambda,H}\) is the orthogonal scalar-character
projection on the finite orbit of the Bessel coefficient.
On the entire character space of dimension
\(((T^k-1)/(T-1))^f\), the Bessel Gram form for this row is a
scalar times the ordinary inner product on covectors. Its
scalar is \((T-1)^f\|p_{\lambda,H}\mathcal F A\|^2\).
\end{corollary}
\begin{proof}
The first theta kernel has the prescribed source character
at \(T\), so its integral depends only on \(eA_\lambda\).
Apply its norm formula to \eqref{eq:lift200}, retaining its
real constant \(C_x\), its finite factor \(P(\mathbf a)\)
and its error \(\Delta_x\).
For the second identity, right translation on the column
group commutes with the unipotent integral defining
\(\mathcal F\). The finite sum therefore gives
\(|Z|p_{\lambda,H}\mathcal F A\). The finite character
space is irreducible under \(\prod_{v\mid T}\mathrm{GL}_k(\mathbb F_T)\).
Its Gram scalar follows by evaluating on the character
sum on the chosen scalar orbit, whose squared coordinate
norm is \(|Z|\). Full proofs, including irreducibility, are in
Corollary \ref{cor:theta199} and Theorem \ref{thm:bessel199}.
\end{proof}

\begin{proposition}[The mass retained by all nontrivial terminal characters]
\label{prop:all200}
Let \(P_{\rm nz}\) remove the trivial scalar character at each
terminal place, using the twisted translations on the cusp row.
Then
\[
   eP_{\rm nz}A=
      \left(\frac{T-2}{T-1}\right)^f L_TA,\qquad
   \|(1-P_{\rm nz})A\|^2\le \frac{5f}{T}\|A\|^2 .
\]
Clearing \((T-1)^f\) in \(P_{\rm nz}\) gives an integral sum
of translations. For an individual character and \(A\ne0\),
one has uniformly on the common intervals
\[
     \frac{\|eA_\lambda\|^2}{\|A_\lambda\|^2}
           =\frac{1+O(f/T)}{(T-1)^f}.
\]
\end{proposition}
\begin{proof}
There are \((T-2)^f\) allowed characters. Their orthogonal
projections on \(A\) are \(A_\lambda/(T-1)^f\).
Summing \eqref{eq:lift200} gives the first formula and
multiplying out the complementary scalar averages proves
integrality. The lost squared norm is
\(\|A\|^2-((T-2)/(T-1))^f\langle A,L_T A\rangle\).
The spectrum in Theorem~\ref{thm:original200} bounds it as
claimed. The unprojected squared norm of a single sum is
\((T-1)^f\langle A,L_T A\rangle\), by the Gram identity in
Theorem \ref{thm:hec199}.
Its projected squared norm is \(\|L_T A\|^2\), giving the ratio.
\end{proof}

For \(A\ne0\) the theta norm in \eqref{eq:theta200} is positive.
For the designated Bessel column \(A=S_0\xi\), the bound still
contains \(\|S_0\xi\|^2\). Likewise \eqref{eq:bessel200} leaves
the projected global coefficient \(p_{\lambda,H}\mathcal F A\)
to be estimated on the integral row and paired with the
designated integral column. These quantities enter the arithmetic
comparison at the original finite level. That comparison and the
uniform estimate for the cohomology sum in \eqref{eq:region198}
remain open in the manuscript.

\subsection{Mixed Gram determinants along the theta lift}\label{nav:VIII-33}
\begin{theorem}[Exterior Gram forms at the original level]
\label{thm:minors239}
For a fixed translate and the data of Theorem~\ref{thm:range239},
sum the maps \(S_c\) on the direct sum of their domains to
obtain \(S:D_{\rm all}\to\mathcal H\).
Use one theta operator on their common real line and at their
common finite level. Put \(\widehat\Theta=\alpha^{-1/2}\Theta\).
For any two lists \(A=(a_1,\ldots,a_s)\) and \(B=(b_1,\ldots,b_s)\)
in \(D_{\rm all}\), write
\[
 G_{A,B}=(\langle Sa_i,Sb_j\rangle)_{ij},\qquad
 H_{A,B}=(\langle\widehat\Theta Sa_i,\widehat\Theta Sb_j\rangle)_{ij}.
\]
The Hermitian forms for one list satisfy
\(e^{-\Delta_x}G_{A,A}\le H_{A,A}\le e^{\Delta_x}G_{A,A}\)
and have the same kernel. Moreover,
\[
 \begin{split}
 e^{-s\Delta_x}\det G_{A,A}
   &\le\det H_{A,A}\le e^{s\Delta_x}\det G_{A,A},\\
 \big|\det H_{A,B}-\det G_{A,B}\big|
   &\le (e^{s\Delta_x}-1)
              \sqrt{\det G_{A,A}\,\det G_{B,B}} .
 \end{split}
\]
For \(1\le s\le\dim D_{\rm all}\) the factor \(e^{s\Delta_x}-1\)
is at most \(\exp(-(\log 2)x^{377}/3)\) on every sufficiently
distant common interval. All determinants of higher order vanish.
\end{theorem}
\begin{proof}
Set \(F_0=S D_{\rm all}\) and restrict \(\widehat\Theta\) to it.
Let \(L\) be the product of the adjoint of this restriction
with the restriction itself, acting on \(F_0\). By
Theorems~\ref{thm:range239} and~\ref{thm:norm177} its quadratic
form lies between \(e^{-\Delta_x}\) and \(e^{\Delta_x}\) times
the squared norm. Applying this to each linear combination of
the \(Sa_i\) proves the Hermitian inequalities and kernel equality.
On \(\bigwedge^s F_0\) the eigenvalues of \(\bigwedge^s L\)
are the \(s\)-fold products of distinct entries in an
orthonormal eigenbasis of \(L\), with multiplicities.
They lie in \([e^{-s\Delta_x},e^{s\Delta_x}]\).

Use the exterior inner product for which the squared norm of a
simple wedge is its Gram determinant. For
\(v_A=Sa_1\wedge\cdots\wedge Sa_s\), and likewise \(v_B\),
the two mixed determinants are respectively
\(\langle v_A,(\bigwedge^s L)v_B\rangle\) for \(H\) and
\(\langle v_A,v_B\rangle\) for \(G\). Taking \(A=B\) proves
the first determinant inequality. Subtracting the identity
and applying Cauchy--Schwarz proves the second, including
when either list is dependent. If \(s>\dim F_0\) every such
wedge is zero.

Here \(\Delta_x=Cmfr\,2^{-(t_{\rm r}+1)/2+10(r+1)^3}\), with
\(t_{\rm r}\ge x^{377}\) by the choice of the rank step.
Lemma~\ref{lem:dimension239} gives \(\log s\le x^{40}\),
\(r=k+m-1=O(x^5)\), and \(mfr\) is polynomial in \(x\).
Thus \(\log(s\Delta_x)\le-(\log 2)x^{377}/2+O(x^{40})\)
uniformly in \(s\). The elementary bound \(e^y-1\le2y\)
for \(0\le y\le1\) gives the claimed estimate.
\end{proof}

\subsection{The prescribed columns at the ordinary real places}\label{nav:VIII-46}
Use the weights in \eqref{eq:letters111}, the spacing of
Theorem~\ref{thm:ordinary164} and the core of
Theorem~\ref{thm:realcore253}. Call the real places outside
the \(h\) designated high places ordinary. Write \(J\) for
all the ordinary places. The desired local highest weights
are the column parameters \(u_1,\ldots,u_k\) minus the half
sum of positive roots of \(U(k)\). The conclusions also
hold in the contragredient convention.

\begin{theorem}[All ordinary places and the obstruction at high places]\leavevmode
\label{thm:word275}
On every sufficiently distant common interval, for each given word,
bottom row and terminal character of Theorem~\ref{thm:small274},
one can impose its prescribed real column type at all the
ordinary places simultaneously on a restricted row after a
determinant twist unramified at every finite place. Its coefficient
on \(H_0(F_\infty)\) over the identity finite component has norm
at least \(\mathcal C_w/(2\sqrt{d_\rho})\), with \(w\) as in
Theorem~\ref{thm:extract275}. The field, terminal lattice,
frequency \(\varpi\) and finite scalar character are those of
the original word. At
each high place the prescribed type has multiplicity zero in
the restriction of the compact core type \(\rho\) under
every determinant twist.
\end{theorem}
\begin{proof}
Let \(n=r+1=k+m\), \(d=n+4\), \(N=k+m(r+6)\), and write \(g_*\)
for the gap of Theorem~\ref{thm:first250}. At the terminal step
the compact parameter, up to duality, has the \(k-1\) entries
\(v_i-1/2\) below all the others. The next lowest entry is
\(a_m=v_1+1/2+g_*\). Indeed the lowest added string is centered
at \(\lfloor(a_{m-1}+a_m)/2\rfloor\) and its radius is
\((d-1)/2\), with \(g_*>d+2\). The string centered at
\(\lfloor(a_0+a_1)/2\rfloor\) contributes \(d\ge k\) entries
at least \(a_m+(m-1)g_*\), since \(a_0=a_m+mg_*\).
Let \(\lambda_i\) be the decreasing list after subtraction of
\((N+1)/2-i\). We obtain
\[
       \lambda_k-\lambda_{N-k+1}\ \ge\ (m-1)g_*-N .
\]
Negating and reversing the list, or shifting it by a constant,
preserves this difference. This proves the same lower bound
for the actual highest weight of \(\rho_\tau\).

At an ordinary place the desired highest weights
\(\mu_i=u_i-((k+1)/2-i)\) have absolute values
\(O(k\Delta+k)\). In the common family \(f\ll_\psi k\) by
Theorem~\ref{thm:rangenew50}. The chosen
\(\Delta=8(f+1)n^3+8r^2+5\) gives \(k\Delta=O_\psi(n^5)\).
We have \(N=O(n^2)\) and \(g_*\ge n^7\).
Also \(w_E\le 2[E:\mathbb Q]^2=O_\psi(n^2)\).
For the last bound, \(\varphi(w_E)\le [E:\mathbb Q]\);
the prime-power formula gives \(p^a/\varphi(p^a)^2\le1\)
at odd primes and at most two at \(p=2\). Multiplication gives
\(w_E\le2\varphi(w_E)^2\).

At each ordinary place choose \(j_\tau\in w_E\mathbb Z\)
nearest the negative of the midpoint of
\([\lambda_{N-k+1},\lambda_k]\). For all sufficiently large
\(x\), half the length of this interval exceeds
\(\max_i|\mu_i|+w_E\), uniformly in the word. Thus every
\(\mu_i-j_\tau\) belongs to it and satisfies all the
inequalities of Proposition~\ref{prop:branch275}.
Set \(j_\tau=0\) at the high places. The tensor product of
these ordinary shifted types occurs in \(\rho|_{K_J}\).
The ratio \(q_\nu\) of Theorem~\ref{thm:extract275} is at
least \(1/d_\rho\), since its numerator is a positive integer.
Use that theorem and multiply the resulting row by the
global determinant character of Proposition~\ref{prop:branch275}.
It has modulus one, is trivial on every unipotent argument,
and takes value one on the terminal scalar operators.
Consequently the projected norm is unchanged and its ordinary
real type is the prescribed one.

At a high place the largest and smallest prescribed highest
weights are \(a\) and \(-b_*\), by \eqref{eq:letters111}.
Their difference is \(a+b_*\), with \(a=\lceil x\exp(x^\varepsilon)\rceil\).
Every entry of the terminal core parameter and its half-root
shift has polynomial size in \(x\) by Theorem~\ref{thm:first250}
and the string formulas above. Any restricted highest weight
has spread at most \(\lambda_1-\lambda_N\) by
Proposition~\ref{prop:branch275}. Taking a dual or a
determinant twist preserves spread, which proves the vanishing.
\end{proof}

\subsection{An integral column on the scalar orbit}\label{nav:VIII-47}
Set \(p(h)=F_0A_\lambda(h,1_f)\) for the row of
Theorem~\ref{thm:types276}. At \(v\mid T\) take the original
column group \(K_{\pi,v}(\ell)=\{u\in\mathrm{GL}_k(\mathbb Z_T):
\ell_v u\equiv\ell_v\bmod T\}\). At the remaining places use
a sufficiently fine level fixing the Fourier coefficient.
Let \(h_\circ\) count the finite classes at the combined level
\(K_\circ\) chosen in the proof. At the identity class take
finite integral generating lists \(v_i\) and \(l_j\) of the
word representation lattice and its dual, of sizes \(I\) and
\(J\). Use a unitary real norm. Denote the smallest eigenvalues
of the positive frame operators \(\sum_i v_i v_i^*\) and
\(\sum_j l_j^*l_j\) by \(b_V\) and \(b_L\), respectively.
They are positive because the lists span over \(\mathbb C\).

\begin{corollary}[An arithmetic column with the prescribed real and terminal characters]\leavevmode
\label{cor:column276}
There is an integral column \(Q_\lambda\) of the full real
word type and terminal character \(\lambda\), at the original
column level above \(T\), satisfying
\[
   \big|\langle Q_\lambda,F_0 A_\lambda\rangle_{[H_0]}\big|
    \ \ge\ \frac{|Z|}{h_\circ}
       \left(\frac{b_V b_L}{d_\mu I J}\right)^{1/2}\|p\|_2 .
\]
It is supported on \(|Z|\) distinct classes above the same
underlying column lattice. In the displayed pairing the
column level away from \(T\) may depend on the restricted row.
For any finite list of \(L\) word and character pairs on the
same field, row group and terminal graph, a single integral
linear combination of the constructed rows gives every
corresponding pairing at least half its individual lower bound.
The integer coefficients can all lie between zero and \(L\).
\end{corollary}
\begin{proof}
The initial row is fixed by the terminal graph stabilizers.
Their intersection with \(H_0\) is
\(K_{\pi,v}(\ell)\) by Theorem~\ref{thm:midlevel131}.
The small Fourier integral commutes with right column
translations. Since the terminal scalars are central on
\(H_0\), the character sum retains these levels on the
column. At every other place intersect its finite level of
smoothness with the stabilizer of the original column lattice.
At one place above the odd split auxiliary prime \(g_1\)
intersect further with the principal congruence subgroup of
depth one. These operations yield an open compact \(K_\circ\).
At every finite class its rational stabilizer at infinity is
a finite group of definite lattice isometries. At the chosen
place it embeds in a conjugate of the principal congruence
group, which is torsion-free by the binomial argument of
Proposition~\ref{prop:terminalcore255}. Thus all stabilizers
are trivial. Each class consequently has mass \(1/h_\circ\)
for quotient probability.

The translates of the identity class by \(Z\) are distinct.
Indeed, an equality of two such finite classes gives a
rational isometry \(\gamma\) relating their representatives
up to \(K_\circ\). Each representative acts by integral
units on the column lattice, so \(\gamma\) preserves that
global definite lattice. At the chosen place above \(g_1\)
it is in the torsion-free principal congruence group. Hence
\(\gamma=1\). Equality now requires the quotient of the two
scalars to lie in \(K_{\pi,v}(\ell)\) for every \(v\mid T\).
The reduction of each covector is nonzero, so each quotient
is one modulo \(T\). For the Teichm\"uller scalars this implies
equality. All these classes use the same underlying lattice.

At the identity class take \(Q_{ji}(s)=l_j(\mu(s)v_i)\)
on the real variable, extended by zero on the other classes.
The lattice here is the tensor of the integral Schur models
of the original lattice at the given embeddings with their
determinant factors. Its matrix coefficients in these lists
are integral class coefficients in the sense of
Proposition~\ref{prop:terminalcore255}. Write
\(p(s)=\operatorname{tr}(M\mu(s))\). Orthogonality gives
\(\|M\|_{\rm HS}^2=d_\mu\|p\|_2^2\). The squared pairings
on the real group with the \(Q_{ji}\) sum to
\(d_\mu^{-2}\sum_{ij}|v_i^* M l_j^*|^2\).
Applying the two lower frame bounds gives a sum at least
\(b_V b_L\|p\|_2^2/d_\mu\).
One pair attains at least \(1/(I J)\) times this squared
bound. The global pairing with that \(Q_{ji}\) is its real
pairing divided by \(h_\circ\).

Take \(Q_\lambda=\sum_{z\in Z}\overline{\lambda(z)}R(z)Q_{ji}\).
It belongs to the same finite level since \(Z\) is central.
On each translated class use \(z^{-1}\) as its representative.
Its integral lattice and real coefficient are the original
ones there, up to the root of unity in the sum. Thus all
coefficients are integral after adjoining the character values.
The sum over \(Z\) is self-adjoint and acts on \(F_0A_\lambda\)
by \(|Z|\). This proves the inequality. Its right side is
positive by Theorem~\ref{thm:types276}.

For the simultaneous assertion fix a witness row \(A_i\) and
column \(Q_i\) just constructed for each pair. Take the common
refinement of their levels away from \(T\). Their pairings are
integrals on the same adelic quotient and are unchanged by
this refinement. Let \(f_i\) be pairing against \(Q_i\) after
the small Fourier map and the indicated character sum. This
is a complex linear functional of the input row. The real
word projector can also be inserted, since it fixes \(Q_i\).
Choose independently and uniformly \(c_1,\ldots,c_L\) from
\(\{0,\ldots,L\}\). With all but \(c_i\) fixed, successive
values of \(f_i(\sum_j c_j A_j)\) on their line have spacing
\(|f_i(A_i)|>0\). At most one lies in the open disk of radius
\(|f_i(A_i)|/2\) about zero. The probability that any of the
\(L\) lower bounds fails is therefore at most \(L/(L+1)<1\).
Some choice satisfies them all. The common terminal row
invariance is preserved under the linear combination.
\end{proof}

The comparison for Dirichlet GRH requires control with the
prescribed row and the original finite level at every place.
The class count \(h_\circ\), the integral frame constants and
the raw factor \(M_B\) in this result enter its normalization.
A uniform comparison on the common intervals, the upper bound
for the full theta period, and the strict cohomology bound
remain open.

\subsection{An integral trace to one fixed theta test}\label{nav:VII-30}
At every \(v\mid T\) take the split matrix model whose row
factor gives positive valuation to \(\varpi\), with the actual
terminal lattice \(\Lambda\) of Theorem \ref{thm:midlevel131}. In
the first theta edge write
\(Y=\binom{X}{Z}\), where \(X\) is square of size \(r\).
Take the coarse test with \(X\) integral invertible and with
the lower block \(Z\) divisible by \(T\). Multiply its values
by the level character divided by the Weil splitting character,
evaluated on \(X\) in the chosen bases. Both take unit values.
At depth \(j\ge1\), retain the part where \(X\bmod T^j\) is
in the image of \(B_v\cap K_v^0\). Depth zero means the whole
coarse test. All depths thus use the same coarse theta datum.

\begin{theorem}[The integral trace and the original period]
\label{thm:trace166}
Write \(D_B=(r(r+1)-k(k-1))/2\), \(I_0=1\), and set
\[
       I_j=T^{(j-1)D_B}(T^k-1)
                     \prod_{h=k+1}^{r}(T^h-1)\quad(j\ge1).
\]
For arbitrary depths \(\mathbf j=(j_v)_{v\mid T}\) let
\(I=\prod_{v\mid T} I_{j_v}\). The normalized average of the
fine finite theta test over the original levels has exact
integer denominator \(I\). Multiplying by \(I\) gives the
coarse integral test. Let \(F_{\mathbf j}\) and \(F_*\) denote
their rows for the same partner input, whenever both lifts
converge and belong to automorphic \(L^2\). Let \(e\) project
onto the original finite level characters at the refined places.
With \(S_{\mathbf j}\) the corresponding fine Bessel map, for
each input \(\xi\) one has
\begin{align*}
  \langle F_*,S_0\xi\rangle
       &=I^2\langle eF_{\mathbf j},e S_{\mathbf j}\xi\rangle,\\
  \|I S_{\mathbf j}^*S_{\mathbf j}-1\|_{\rm op}
       &\le C_x T^{-\eta\sum_v j_v},\qquad \eta=L-4fr^2 .
\end{align*}
In particular the depth-one denominator is prime to \(T\) at
every refined place. The exact \(T\)-valuation of the total
denominator is \(D_B\sum_v\max(0,j_v-1)\).
\end{theorem}
\begin{proof}
The split radical has \(a=(r^2-(k+1)^2-2s)/2\) free matrix
entries in the basis of \(\Lambda\). On the middle block an
integral column isometry \(h\) must preserve the fixed nonzero
residue covector \(\ell\). Its action modulo \(T^j\) also
retains \((\ell-\ell h)/T\bmod T^j\). Including this last
row gives \(T^k|\mathrm{GL}_k(\ZZ/T^j\ZZ)|/(T^k-1)\) possible middle
blocks. Dividing the full row group order by this number and
\(T^{ja}\) gives exactly \(I_j\), as proved in
Theorem \ref{thm:index165}.

The translates of the fine theta test partition the coarse
support by left cosets. Their corrected values can be recovered
separately by evaluating at a representative top block with
zero lower block. Products of these evaluations give a unit
divided by \(I\) for the average, proving the exact denominator.
Theorem \ref{thm:tests165} also realizes the corresponding finite
subgroup by a Weil matrix coefficient on the same pair.
Right translation of a theta row acts on its test by this
representation, so finite averaging gives \(eF_{\mathbf j}
=F_*/I\). For the Bessel seed the same fraction of the original
right coset survives, giving \(e S_{\mathbf j}=S_0/I\).
The two identities prove the period equality. The nonidentity
rational positions satisfy the product of the congruence
divisibilities at all the refined places. Taking the norm of
one nonzero frame coordinate gives their height separation
by \(T^{\sum_v j_v}/C_x\). Applying the integral tail estimate
of Theorem~\ref{thm:depth164} as in Corollary \ref{cor:trace165}
proves the norm bound. Finally every \(T^h-1\) is a unit at
\(T\), which gives the asserted valuation.
\end{proof}

For total depth polynomial in \(x\),
\(\log I^2\le2D_B(\sum_v j_v)\log T=o(b/x^A)\) for every fixed
\(A>0\). The preceding proof controls this explicit part of the
finite normalization. At fixed \(x\), the norm estimate concerns
the whole fine image. The period equality selects its projection
at the original level and its pairing with a specified theta
row. A uniform nonzero, arithmetically normalized pairing of
this kind along the full chain, as well as the remaining
strict cohomology bound in \eqref{eq:tail114}, is still required
to prove \eqref{eq:window}.

\begin{proposition}[The nonvanishing test on an original word]
\label{prop:word166}
Fix a word \(c\), a unit real vector \(v\in U\), and write
\(S\ell=S_0(v\otimes\ell)\) on its full class space \(D_c\).
Let \(R\subset D_c\) be its subspace of ramified columns.
For each rational double coset outside the stabilizer of the
entire isotropic space let
\(O_\gamma\) be its operator on \(D_c\) given by the integral
in equation \ref{eq:geom163}, with this common \(v\). Then
\[
 E_c=\sum_{\gamma,\delta}
       \operatorname{tr}(O_\gamma^*O_\delta)
       =\|S^*S-1\|_{\rm HS}^2\ge\dim\ker S .
\]
The double series is absolutely convergent. If
\(E_c<\dim R\), some ramified column has nonzero image at
the prescribed level. For the closed span \(\mathcal T\) of
any specified collection of convergent theta rows there,
let \(p_{\mathcal T}\) denote orthogonal projection. One has
\[
 \|p_{\mathcal T} S|_R\|_{\rm HS}^2
  =\dim R+\sum_\gamma\operatorname{tr}_R(O_\gamma)
                   -\|(1-p_{\mathcal T})S|_R\|_{\rm HS}^2 ,
\]
where the trace on \(R\) uses orthogonal compression.
Strict positivity of this expression is equivalent to a
nonzero pairing of a ramified column with one of those rows.
\end{proposition}
\begin{proof}
Use the absolute integral bounds of Theorem~\ref{thm:ordinary164}
on every pair of a fixed orthonormal basis of \(D_c\).
There are finitely many pairs, so
\(\sum_\gamma\|O_\gamma\|_{\rm HS}<\infty\).
The terms from its stabilizer give the identity. Indeed their
absolutely convergent subseries unfolds to the intermediate
pairing evaluated at \(K^0\) in Theorem \ref{thm:midnorm131}.
The remaining cosets have two distinct isotropic spaces. Thus
\(S^*S-1=\sum_\gamma O_\gamma\). Squaring the Hilbert--Schmidt
norm gives the first identity since each trace product is
bounded by the product of the two Hilbert--Schmidt norms.
On an orthonormal basis of the kernel the difference from
the identity acts by minus one, giving the dimension bound.
If all of \(R\) vanished its dimension would be at most
\(E_c\). This is the dimension argument at
\srcloc{Development/AbelianVarietyTheory/Proofs/DecoratedLatticeProofs/NonvanishingFromTruePack}{26}.
The source asks for a surviving ramified column for every
word in \srcloc{Development/AbelianVarietyTheory/Definitions/LatticeFrames/PairForm}{143}.
For the second identity sum Pythagoras for \(S\ell\) over
an orthonormal basis of \(R\) and use the parabolic expansion
for \(\|S|_R\|_{\rm HS}^2\). The projected operator is zero
exactly when every such image is orthogonal to the whole
collection. This proves the equivalence. On the varying
fields the comparison of \(E_c\) with \(\dim R\), and control
of the component orthogonal to the specified theta rows,
are additional estimates.
\end{proof}

At \srcloc{Development/AbelianVarietyTheory/Proofs/DecoratedLatticeProofs/DiagonalRowOfSiegelWeilUnitary}{50}
the proposed bound is assembled from estimates for degenerate
terms, collisions and two corank-one contributions. Absolute
convergence in Proposition~\ref{prop:word166} justifies the
expansion to which such estimates would apply at the original
level.

\subsection{A second moment on the terminal integral coefficient lattices}\label{nav:VI-39}
Fix an original field, finite levels and ascents as in
Theorem~\ref{thm:coregram258}. Use probability measure on the
definite terminal quotient. Denote the whole ascent with the
selected tests by \(T\), and its restriction to \(H(v)\) by
\(T_v\). For a terminal form \(e\), write \(T^\dagger e\) for
the reverse composition of the adjoint theta integrals with
the same tests and measures. Write \(P_v\) for orthogonal
projection onto the bottom subspace \(H(v)\). Fix a complex
linear functional \(\ell\) on that finite-dimensional Hilbert
space. Its dual norm is \(\|\ell\|_*\).

For each terminal class \(c\), in the notation of
Proposition~\ref{prop:terminalcore255}, take a basis \(v_{c,b}\)
over the coefficient field in the integral lattice \(W_c\) and
a basis \(l_{c,a}\) in its integral dual. Such bases can be
extracted from any finite module generating lists because
the lattices span their spaces over the field. No freeness
over the ring of integers is required. Let \(e_{c,a,b}\) be
zero off \(c\) and on \(c\) the function
\(g\mapsto l_{c,a}(\rho(g)v_{c,b})\).
Use a Hermitian inner product linear in its second argument.
Let \(G\) be the Gram matrix of this ordered family of
\(M=h_m d_\rho^2\) functions and let
\(y_{c,a,b}=\overline{\ell(P_v T^\dagger e_{c,a,b})}\).

\begin{theorem}[A second moment of the descending periods]\leavevmode
\label{thm:moment258}
With \(C_{\rm fix}\) and \(\delta\) of Theorem~\ref{thm:coregram258}
for the selected measures,
\[
  e^{-\delta} C_{\rm fix}\|\ell\|_*^2
       \ \le\ y^*G^{-1}y\
       \le\ e^\delta C_{\rm fix}\|\ell\|_*^2 .
\]
The matrix \(G\) is positive definite. It is zero between classes
and its within-class entries are
\[
 G_{(c,a,b),(c,a',b')}=
 \frac{\langle l_{c,a},l_{c,a'}\rangle_*
       \langle v_{c,b},v_{c,b'}\rangle}
      {h_m d_\rho}.
\]
In particular at least one of the given integral terminal
coefficients satisfies
\[
 |\ell(P_vT^\dagger e_{c,a,b})|
 \ \ge\
 \left(\frac{e^{-\delta} C_{\rm fix}}
 {M\|G^{-1}\|_{\rm op}}\right)^{1/2}\|\ell\|_* .
\]
These conclusions apply to the first restriction functional
\(\ell(A)=\mathcal B(A,T_0w)\) of Theorem~\ref{thm:allrows251}
on \(H(v)\). For each prescribed nonzero smooth row vector
\(T_0w\), some pure compact-finite \(v\) makes this functional
nonzero.
\end{theorem}
\begin{proof}
At the prescribed level, Proposition~\ref{prop:terminalcore255}
identifies the quotient with \(h_m\) copies of the compact
group, each of mass \(1/h_m\). On a copy, the coefficients of
two bases as above are a basis for the full right
\(\rho\)-isotypic functions. This follows by the compact
matrix coefficient decomposition. Schur orthogonality applied
to a pair of coordinates gives the stated Gram entry, with
the dual inner product taken on complex linear functionals.
Different classes have disjoint supports. Hence these functions
are a basis of the full terminal \(\rho\)-space and their
Gram matrix is invertible.

Let \(I:\CC^M\to\mathcal A_m\) take coordinates to this basis.
Then \(I^*I=G\), and \(IG^{-1}I^*\) is the identity on its
range. Let \(a_\ell\in H(v)\) be the Riesz representative,
so \(\ell(A)=\langle a_\ell,A\rangle\).
For a terminal vector \(e\), the adjoint calculation in
Theorem~\ref{thm:downcore254} gives
\(T_v^*e=P_vT^\dagger e\). Indeed one can test against any
element of \(H(v)\) and move each kernel across its Petersson
pairing. These integrals converge by cuspidal decay on both
sides at each noncompact step, as proved there. The last
orthogonal projection expresses restriction of the domain
to \(H(v)\). Consequently \(y=I^*T_v a_\ell\). The real type
of the ascent is \(\rho\) at the terminal group and it has
the given finite level, so \(T_v a_\ell\) is in the range
of \(I\). We obtain the identity
\[
                  y^*G^{-1}y=\|T_v a_\ell\|^2.
\]
Apply Theorem~\ref{thm:coregram258} to \(a_\ell\), whose norm
equals \(\|\ell\|_*\). The individual bound follows by using
\(y^*G^{-1}y\le M\|G^{-1}\|_{\rm op}\max_\nu|y_\nu|^2\).

For the first restriction, the given row form is smooth and
cuspidal and the bottom core forms in \(H(v)\) are cuspidal.
The restriction integral exists by Theorem~\ref{thm:allrows251}
and its proof. That theorem supplies a compact-finite core
partner with nonzero pairing. At the fixed finite level such
a partner is a finite sum of pure real compact-finite tensors
times finite multiplicity vectors. Indeed only finitely many
compact types occur in that vector, their spaces in the real
representation are finite-dimensional, and the real
Harish-Chandra module is the algebraic tensor product over
the real places. At least one tensor has nonzero pairing.
Normalizing its real part gives the required \(v\).
\end{proof}
The common real norm factor \(C_{\rm fix}\), the metric of the
integral terminal coefficients and the dual norm of the
restriction enter the individual lower bound. Section~\ref{sec:missing}
requires absolute arithmetic estimates after the specified
period projection on these fields.

\subsection{The valuation of a prescribed signed period coefficient}\label{sec:fixedadic332}
The normalized arithmetic column of Theorem~\ref{thm:column242}
is a finite signed sum of algebraic coefficients under the hypotheses
of that theorem. Keep its levels and its prescribed endpoint.
Section~\ref{sec:fixedcolumn332} writes the resulting scalar as
\(Z=\sum_{t\bmod L} C_t\zeta^t\), where \(\zeta\) is a primitive
\(L\)-th root of unity. The same notation applies to the prescribed
joint auxiliary row when its fiber sums are algebraic.
We give a finite computation at every coefficient prime using these
very \(C_t\), allowing ramification of any order.

\begin{theorem}[Fixed values in a ramified coefficient field]\leavevmode
\label{thm:fixedadic332}
Choose a finite place above \(p\) of a number field containing the
data, and write \(E\) for its completion. Normalize \(v_p(p)=1\).
Let \(K\) be the maximal unramified subextension of \(E/\QQ_p\).
Put \(e=[E:K]\) and choose a uniformizer \(\varpi\) of \(E\).
Write \(L=p^b m\), with \((m,p)=1\), and factor
\(\zeta=\omega\eta\) with \(\omega\) of order \(m\) and \(\eta\)
of order \(p^b\). Then \(\omega\in K\).

If \(b>0\), put \(M=p^b,\ D=p^{b-1},\ n=(p-1)D,\ h=e/n\).
Expand uniquely over \(K\)
\[
 C_t\omega^t=\sum_{i=0}^{h-1}\varpi^i
                  \sum_{j=0}^{n-1} c_{ijt}\eta^j .
\]
For \(v\bmod M\) and \(0\le i<h\) set
\[
 A_{i,v}=\sum_{j+t\equiv v\bmod M} c_{ijt}.
\]
Here \(j,t\) run over the full ranges in the expansion of \(Z\).
For \(0\le k<n\), writing \(k=v+uD\) with \(0\le v<D\),
put
\[
 B_{i,k}=A_{i,k}-A_{i,v+(p-1)D},\qquad
 D_{i,j}=\sum_{k=j}^{n-1}\binom{k}{j} B_{i,k}.
\]
The valuation of the prescribed scalar, with \(v_p(0)=+\infty\), is
\begin{equation}\label{eq:fixedadic332}
 v_p(Z)=\min_{\substack{0\le i<h\\0\le j<n}}
              \left(v_p(D_{i,j})+\frac{i+hj}{e}\right).
\end{equation}
In particular, for each integer \(a\), \(v_p(Z)\ge a\) if and only if
\(B_{i,k}\in p^a\mathcal O_K\) for every \(i,k\). Vanishing is equivalent
to equality of all \(p\) coefficients \(A_{i,v+uD}\) as \(u\)
runs from zero to \(p-1\), for every \(i,v\).

If \(b=0\), expand \(C_t\omega^t=\sum_{i=0}^{e-1}c_{it}\varpi^i\)
over \(K\). Then
\[
 v_p(Z)=\min_{0\le i<e}\left(v_p\Big(\sum_t c_{it}\Big)+i/e\right).
\]
\end{theorem}
\begin{proof}
The coprime orders give the factorization of \(\zeta\) in its
cyclic group. Reduction of a prime-to-\(p\) root of unity is
injective on the group of such roots. Its reduction has a unique
lift of that order by Hensel's lemma in \(K\), which has the same
residue field as \(E\). This proves \(\omega\in K\).
For \(b>0\), the polynomial \(\Phi_{p^b}(1+T)\) is Eisenstein
over \(\mathcal O_K\). Its constant coefficient is \(p\), and
modulo \(p\) it equals \(T^n\) by
\(\Phi_{p^b}(X)=\sum_{u=0}^{p-1}X^{uD}\). Consequently
\(K(\eta)/K\) is totally ramified of degree \(n\),
\(v_p(\eta-1)=1/n\) and \(n\mid e\).
The \(e\) elements \(\varpi^i(\eta-1)^j\) in the indicated
ranges have all the distinct valuations
\((i+hj)/e\), for \(0\le i+hj<e\).
The valuation group of \(K^\times\) is \(\ZZ\), so in a nonzero
\(K\)-linear combination there is a unique term of least valuation.
In particular these elements are linearly independent. Their number
is \([E:K]\), so they form a basis. Replacing \((\eta-1)^j\) by
\(\eta^j\) uses an integral triangular basis change with diagonal
one and proves existence and uniqueness of the asserted coordinates.

Substitution and collection modulo \(M\) gives \(\sum_i\varpi^i
\sum_v A_{i,v}\eta^v\). The identity
\(\sum_{u=0}^{p-1}\eta^{v+uD}=0\) removes the last term in each
such progression. The remaining coefficients are exactly
\(B_{i,k}\). Expanding \(\eta^k=(1+(\eta-1))^k\) gives
\(D_{i,j}\) in the basis with distinct valuations. The unique
least term proves \eqref{eq:fixedadic332}, including the zero
case by the stated convention. Since \(0\le(i+hj)/e<1\), an
integer threshold \(a\) is attained exactly when all the \(D_{i,j}\)
belong to \(p^a\mathcal O_K\). The integral triangular change of
basis has integral inverse, so the same holds for the \(B_{i,k}\).
It also identifies vanishing with all the displayed progression
differences being zero. When \(b=0\), powers \(\varpi^i\) for
\(0\le i<e\) alone have distinct valuations modulo \(\ZZ\),
and the same argument applies.
\end{proof}

\begin{corollary}[A test before any average over frequencies]\leavevmode
\label{cor:unramtest332}
At one auxiliary place \(v\) of residue characteristic \(p\)
let the joint fiber weights in Theorem~\ref{thm:signedfibers331}
at a prescribed frequency \(\nu\) belong, under the chosen
coefficient embedding, to an unramified extension \(K/\QQ_p\).
For the Haar formulation take the fiber means as weights, removing
the common \({\cal V}\). Set \(T_y=\sum_b\chi(b)H_\nu(b,y)\).
After removing all factors independent of the position fibers,
the value is \(U=\sum_yT_y\zeta^{\nu y}\) with
\(\zeta=\psi(p^{-s})\). For any integer \(a\),
\[
 v_p(U)\ge a\quad\Longleftrightarrow\quad
 T_{y+p^{s-1}}-T_y\in p^a\mathcal O_K\quad
            \text{for every }y\bmod p^s .
\]
In particular this one frequency gives zero exactly when \(T\)
is constant on each coset of \(p^{s-1}\ZZ/p^s\ZZ\). If the \(T_y\)
are integral and at least one of these differences is a unit, then
\(0\le v_p(U)<1\).
\end{corollary}
\begin{proof}
The character \(\chi\) takes values in the \((p-1)\)-st roots of
unity in \(\QQ_p\), so all \(T_y\) belong to \(K\). Apply the
theorem in \(E=K(\zeta)\), with \(L=p^s\), after permuting the
indices by multiplication by \(\nu\). There is one index \(i\);
in the initial expansion only \(j=0\) is used. Multiplication by
\(\nu\) permutes the progressions modulo \(p^{s-1}\), and permutes
their members. The theorem tests equality modulo \(p^a\) to one
designated member of each progression. This is equivalent to equality
for every pair of successive members because the progression is cyclic.
Equality in \(K\) follows from equality modulo every \(p^a\).
Integrality and failure of the test with \(a=1\) prove the last
assertion. The other factors in the Haar integral are
\({\cal V}\eta_\chi p^{r-1-(s+1)d}G(\chi)^h\chi(\bar\nu)^{-1}\)
with the notation of Theorem~\ref{thm:signedfibers331}.
\end{proof}

For the normalized definite column take
\(Z=D_0(\prod_i A_i^{\rm jet})\mathcal A\), with its original
class tests as in Theorem~\ref{thm:finitecolumn313}. There \(Z\)
is an algebraic integer. Once the rational-array test of
Theorem~\ref{thm:scalarorbits332} detects \(Z\ne0\), the same
scalar satisfies
\[
 \frac{1}{[B:\QQ]}\sum_{w\ {\rm finite}}
       \operatorname{ord}_w(Z)\log Nw
 \ \le\ \tfrac12\log\left(\frac1{\varphi(L)}
                      \sum_a b(a)^*G(a)b(a)\right)
\]
in the notation of \eqref{eq:fixedmoment332}. Indeed the finite
sum is the logarithm of its absolute norm divided by \([B:\QQ]\).
That is the mean of \(\log|\sigma Z|\), bounded by half the logarithm
of the mean of \(|\sigma Z|^2\) by concavity. At a place \(w\mid p\)
the term uses \(\operatorname{ord}_w(Z)=e_w v_p(Z)\), with \(e_w\)
its absolute ramification index. Formula~\eqref{eq:fixedadic332}
and its \(b=0\) case compute these orders from the same inputs.
The uniform comparison in Section~\ref{sec:missing} still requires
control of the actual signed class coefficients as \(x\) varies
and their comparison to the prescribed analytic pairing.


\subsection{The norm and denominators of the split transfer}\label{sec:mixedfourier334}

Retain the notation and the additive Haar measures of
Section~\ref{sec:matchingframes334}. Write \(p_m\) for the measure of
the primitive locus in \(M_{m,n}(O)\). Choosing the rows over the
residue field gives \(p_m=\prod_{i=0}^{m-1}(1-q^{i-n})\).
Put \(L=e-c\ge0\) and let \(N_b(L)\) be the number of \(O\)-lattices
\(J\) between \(\pi^L O^b\) and \(O^b\), including both endpoints.

\Needspace{5\baselineskip}
\begin{theorem}[Exact norm at a ramified split place]\label{thm:frameoperator334}
On the full space of level-\(e\) tests, with additive \(L^2\) norms,
the transfer of Theorem~\ref{thm:matchingframe334} satisfies
\[
       \|\mathcal T\|_{\rm op}^2
       =\frac{p_a}{p_b}N_b(e-c),\qquad
       \frac{p_a}{p_b}=\prod_{j=b+1}^{a}(1-q^{-j}).
\]
On the primitive-supported covariant subspace its squared norm ratio
for every nonzero vector is \(p_a/p_b\).
\end{theorem}
\begin{proof}
Let \(dk\) denote Haar probability on \(K_n\). The frames \(X_k,Y_k\)
separately have the normalized additive distributions on the corresponding
primitive loci. Indeed \(K_n\) is transitive there by basis completion,
and preserves additive measure. Thus
\(\|\mathcal T\Psi\|_2^2=p_a\int_{K_n}|A_\Psi(Y_k)|^2\,dk\).
In the primitive covariant case \(A_\Psi=\Psi\) on these frames. The
formula for \(p_m\) gives the asserted ratio.

Let \(\Omega\) consist of invertible integral \(g\) with largest Smith
exponent at most \(L\). Lemma~\ref{lem:ramifiedsum334} permits us to
restrict \(A_\Psi(Y_k)\) to \(\Omega\). Two such matrices have the same
column lattice precisely when they belong to one coset \(g_JK_b\),
so \(\Omega\) has multiplicative measure \(N_b(L)\). Cauchy--Schwarz gives
\[
 |A_\Psi(Y_k)|^2\le N_b(L)
       \int_\Omega |\Psi(gY_k)|^2|\det g|^n\,d^\times g .
\]
For each lattice \(J\) fix \(g_J\) with \(g_J O^b=J\). A full-rank
matrix with column lattice \(J\) is uniquely \(g_JY\) with \(Y\)
primitive. Left multiplication by \(g_J\) has additive Jacobian
\(|\det g_J|^n\). In the double integral over \(k\) and \(g_JK_b\)
the factor from \(K_b\) can be absorbed into the invariant distribution
of \(Y_k\). Its contribution, including \(|\det g_J|^n\), is
\(p_b^{-1}\int_{XO^n=J}|\Psi(X)|^2\,dX\). Summing the disjoint
sets proves the upper bound.

For sharpness choose a nonzero covariant function \(f\) on primitive
\(b\)-frames modulo \(\pi^c\). One may give it value one at a chosen
basis of a single direct summand, extend by the unit character under
changes of basis, and set it to zero on the other summands. The
character is trivial on changes of basis congruent to the identity
modulo \(\pi^c\). Define a test on the matrices with allowed column lattice by
\[
     \Psi(gY)=\nu(\det g)^{-1}|\det g|^{-z}f(Y),
                \qquad g\in\Omega,\quad Y\ {\rm primitive},
\]
and set it to zero elsewhere. A second decomposition has the form
\((gk',k'^{-1}Y)\) with \(k'\in K_b\), giving the same value by
covariance. Perturbing \(gY\) by \(\pi^e U\), \(U\) integral, changes
\(Y\) by \(g^{-1}\pi^e U\in\pi^c M_{b,n}(O)\), so the new frame is
primitive with the same \(f\). The allowed set and its complement are
therefore unions of residue classes. This test has level \(e\). Its
integrand on \(\Omega\) equals \(f(Y)\), and its support uses all allowed
column lattices and no others. Both inequalities in the upper bound are
equalities, and the test has positive norm.
\end{proof}

\Needspace{5\baselineskip}
\begin{proposition}[The lattice count at excess depth]\label{prop:smindex334}
Set \(\gamma_0=1\). For integers \(0\le u_1\le\cdots\le u_b\le L\),
let \(m_j\) be the multiplicity of \(j\) among them. Then
\[
 N_b(L)=\sum_{u_1\le\cdots\le u_b}
  q^{\sum_{i>j}(u_i-u_j)}\frac{\gamma_b}{\prod_{j=0}^L\gamma_{m_j}}.
\]
In particular, with \(d_b=\lfloor b^2/4\rfloor\),
\(q^{Ld_b}\le N_b(L)\le (L+1)^b\exp(2b)q^{Ld_b}\), and \(N_b(0)=1\).
\end{proposition}
\begin{proof}
Each Smith type forms one \(K_b\)-orbit. The stabilizer of
\(\operatorname{diag}(\pi^{u_i})O^b\) consists of the integral
invertible matrices with \(ij\)-entry divisible by
\(\pi^{\max(0,u_i-u_j)}\). The unconstrained reduction is block upper
triangular and is invertible precisely when each equal-exponent diagonal
block is invertible. Its additive mass is therefore
\(q^{-\sum_{i>j}(u_i-u_j)}\prod_j\gamma_{m_j}\). Dividing \(\gamma_b\)
by this mass gives the index. The exponent is
\(\sum_i(2i-b-1)u_i\), at most \(Ld_b\). For \(L>0\) the type with
\(\lfloor b/2\rfloor\) entries \(L\) and all others zero attains this
exponent. Its ratio of \(\gamma\)'s is at least one by comparing the
last \(\lfloor b/2\rfloor\) factors of \(\gamma_b\) with the first ones.
For the upper bound there are at most \((L+1)^b\) possible types and at
most \(b\) nonempty blocks. The inequality
\(-\log(1-q^{-i})\le2q^{-i}\) gives \(\gamma_j\ge\exp(-2)\).
For \(L=0\) there is exactly one lattice.
\end{proof}

\Needspace{5\baselineskip}
\begin{theorem}[The transfer in the two partial Fourier models]\label{thm:mixedfourier334}
Let \(P\) be the set of columns Fourier transformed on the small side,
\(d=|P|\), and use the complementary columns on the large side. Suppose
the additive character \(\psi\) has conductor \(O\), meaning its dual
lattice to \(O\) is \(O\). The linear map between the original level-\(e\)
tests obtained by matching these two charts has the same squared operator
norm \((p_a/p_b)N_b(e-c)\), for every \(P\). If \(D_0\Psi\) has algebraic
integer values and \(D_1\nu(\pi)\) is an algebraic integer, where \(D_0,D_1\)
are positive integers, all coefficients of the output divided by
\(\sigma_P=(\nu(\pi)q^{-z})^{e(a-d)}\) become algebraic integers upon
multiplication by
\[
         D_0 q^{ebd} H_b(e)\,(D_1 q^\delta)^{b(e-c)} .
\]
\end{theorem}
\begin{proof}
For a level-\(e\) test the partial Fourier transform \(\widehat\Psi\)
integrates the entries in \(P\) against \(\psi\) of the entrywise
pairing. It vanishes outside the dual of the translation lattice in
those columns, by translating the integral by a period. Inside that
dual it is the sum over \(R_e^{bd}\) with coefficient \(q^{-ebd}\).
Let \(D'\) be diagonal with \(\pi^{-e}\) in the columns \(P\) and one
elsewhere. Then \(\Psi_1(Y)=\widehat\Psi(YD')\) has level \(e\).
Put \(D=\pi^{-e}D'^{-t}\), the matrix with \(\pi^{-e}\) on \(P^c\)
and one elsewhere. For a primitive frame \(X_k\), \(k\in K_n\), use
\(h=kD^{-t}\) in the definition of \(F\). This gives \(X_h=X_k D\)
and \(Y_h=\pi^e Y_kD'\). Since
\(\det(D^{-t})=\pi^{e(n-d)}\), small-side covariance gives
\[
 F_{\widehat\Psi}(X_k D)
   =(\nu(\pi)q^{-z})^{e(n-d-b)}
       \nu(\det k)\,A_{\Psi_1}(Y_k)
   =\sigma_P\,(\mathcal T\Psi_1)(X_k).
\]
Restrict \(F_{\widehat\Psi}\) to primitive rows relative to \(O^nD\),
extending it by zero. Every row space again has exactly one \(K_a\)-coset
of bases primitive in this lattice. The proof of
Theorem~\ref{thm:matchingframe334} now gives \(B=S\) in the transformed
charts at every \(h\). Apply the inverse partial Fourier transform for
\(\psi^{-1}\) on \(P^c\) to recover the large-side test. The result is
supported on integral matrices and is constant modulo \(\pi^e\).
Integration in each of its dual variables is an unnormalized sum indexed
by \(R_e\), since each coset of \(O\) has mass one. Finite character
orthogonality gives Fourier inversion in the two respective charts.

On the self-dual additive measures each partial Fourier transform
preserves the squared norm, by the same orthogonality. Converting the
small-side transform to the coordinates of \(\Psi_1\) multiplies its
squared norm by \(q^{-ebd}\). Returning from the standard large frame
to \(O^nD\) multiplies the squared norm by \(q^{ea(n-d)}\). Together
with \(|\sigma_P|^2=q^{-ne(a-d)}\) these factors have product one,
since \(a+b=n\). The Fourier transforms and diagonal changes biject
the corresponding level spaces, so the extremizer of
Theorem~\ref{thm:frameoperator334} also gives equality in these charts.

All Fourier phases and the factors \(\nu(\det k)\) are roots of unity.
In a contributing summand of Lemma~\ref{lem:ramifiedsum334} the
uniformizer exponent \(s\) satisfies \(0\le s\le b(e-c)\), and its only
additional factor is \(\nu(\pi)^s q^{-\delta s}\) times a unit character
value. The small Fourier coefficient, the group mass \(H_b(e)^{-1}\),
and these factors give the displayed common integer. Keeping
\(\sigma_P\) explicit also covers the case \(d>a\).
\end{proof}

\Needspace{5\baselineskip}
\begin{corollary}[The split factors on the period interval]\label{cor:splitproduct334}
For \(x\in[V,V^{3/2}]\) as in Section~\ref{sec:missing}, fix a
complementary-rank square as above on \(F=E_x^+\). At its marked split
places \(v\in T=S_0\cup Q(\xi)\) use the local characters, Fourier blocks
and levels of that square. Write \(q_v\) for the residue cardinality and
\(c_v\) for the conductor exponent of the quotient character. For
\(e_v\ge c_v\) at every marked place, the squared norm of the tensor
transfer on their full level spaces is
\[
 \prod_{v\in T}
  \left(N_{b_v}(e_v-c_v)
        \prod_{j=b_v+1}^{a_v}(1-q_v^{-j})\right),
\]
where the lattice count is taken over \(O_v\). On any prescribed subspace
this product is an upper bound. At minimal levels the map on the
primitive covariant data in the linear charts has unit elementary
divisors. The negative logarithm of the product of the primitive measure
ratios is \(O([F:\QQ]\,2^{-(b_*+1)})\), with \(b_*=\min_{v\in T}b_v\)
when \(T\) is nonempty.
\end{corollary}
\begin{proof}
The character quotient is ramified at the marked places of this square,
and their additive characters give self-dual integral lattices. Apply
Theorem~\ref{thm:mixedfourier334} place by place. Singular values multiply
under Hilbert tensor products, or one can apply the bound one factor at
a time and tensor the extremizers. At \(e_v=c_v\) only invertible
residue matrices contribute. The isomorphism on covariant frames in
Theorem~\ref{thm:matchingframe334} is integral in both directions.
If an input is specified at \(e_v<c_v\), a unit in \(U_{e_v}\) with
nontrivial character can multiply one row of the integration matrix
without changing its residue. This makes the entire small section and
the matched output zero at that level. Thus all levels can be tested on
the fields where they occur.

Above a rational prime \(p\) there are at most \([F:\QQ]\) places and
\(q_v\ge p\). Summing \(-\log(1-q_v^{-j})\le2q_v^{-j}\) first in
\(j\ge b_*+1\) gives at most \(4[F:\QQ]\sum_{m\ge2}m^{-(b_*+1)}\).
For \(u\ge2\) the latter sum with exponent \(u\) is at most
\(2^{-u}(1+2/(u-1))\), by integration from 2 to infinity. This proves
the bound uniformly on the original interval.
\end{proof}

At the other split places, Proposition~\ref{prop:sphericalframe334}
supplies the common normalized spherical value. Almost all its local
values on an adelic group element are one, since that element is integral
there and the splitting characters are unramified. The resulting products
of split sections are therefore equal pointwise. Any common finite linear
combination of their translates, in particular the specified finite
character average acting on those sections, preserves that equality.
For the global period in Section~\ref{sec:missing} one still needs to
compare the two global functionals on these matched sections and to
bound the signed value on the designated integral column. A norm on the
full space above gives no lower bound on that value.


\subsection{The kernel on refined terminal columns}\label{sec:projectiveperiod337}
Use the class tables of Corollary~\ref{cor:rank111} with the
determinant character removed as in Theorem~\ref{thm:original200}.
Thus a local scalar acts by permuting the nonzero residue covectors.
In the original action the statements below use the corresponding
scalar character obtained by putting the determinant twist back.
Write \(a=k\) for the column rank of that corollary and \(s\ge1\)
for the flag length. Keep both auxiliary primes, the integral
lattices away from \(T\), the real components and the field fixed.
Use the first Bessel position with the marked halves of
Section~\ref{sec:exchange336}.

\begin{theorem}[Factorization through projective residue covectors]\label{thm:projectiveperiod337}
At every terminal place give the row the graph compact type of
Theorem~\ref{thm:midlevel131}. On the refined column type of
multiplicity \(R_c\), the pairing with all such rows factors through
summing the column coordinates over scalar orbits at each place.
The rank in the column variable is at most
\[
             B_c=d_c\left(\frac{T^a-1}{T-1}\right)^f .
\]
It is zero on the subspace of dimension \(R_c^\circ\) in
Corollary~\ref{cor:rank111}. On the integral class table the
map that sums over the fibers is split surjective and every one
of its nonzero Smith factors equals one over \(\mathbb Z\).
\end{theorem}
\begin{proof}
Fix any row of the indicated type. At a terminal place, every
scalar unit in \(H(O)\) preserves \(J_\ell\). Theorem~\ref{thm:graphtransport337}
makes its pairing invariant under these units on the column.
The other exchanges and all conditions at other places commute
with this action. Thus two basis vectors with covectors in the
same product of scalar orbits have equal values for every row.
Conjugate linearity on the column gives the conjugate of the
sum of its coordinates on each such fiber.

The unrefined classes have trivial stabilizers, by
Proposition~\ref{prop:column110}. Above each class the refinements
are the independent choices of a nonzero covector at the \(f\)
places. The image of each coordinate under the fiber sum is its
tuple of projective covectors. Choose one representative in every
fiber and send the corresponding projective basis element to it.
This is an integral section. Within each fiber the differences
from the representative form an integral basis of the kernel.
The same computation applies to each multiplicity coordinate of
the unrefined table. There are \(B_c\) image coordinates. A
function with zero sum separately on every local scalar orbit
has zero sum on the whole fiber. This proves all assertions,
including the integer Smith calculation before extension of scalars.
\end{proof}

\begin{corollary}[Loss in the proposed Gram estimate]\label{cor:rankobstruction337}
Let \(D_c\) be the map from this refined column into the dual row
space furnished by the first-position pairing on graph rows, using
the Hilbert space identification with conjugate duals as needed.
The same assertions apply after projecting arbitrary rows to this
graph compact type before evaluation. For any scalar \(u\) one has
\[
 \|(uD_c)^*(uD_c)-I\|_{\rm HS}^2\ \ge\ R_c-B_c .
\]
On the separately zero-sum subspace the squared discrepancy from
the identity is exactly \(R_c^\circ\). For a row \(A\) of the original
graph type and its signed sum \(A_\lambda\) of
Theorem~\ref{thm:original200}, the raw first-position pairing
\(\mathcal J_\phi(A_\lambda)\) vanishes for every column whenever
\(\lambda\) is nontrivial at at least one terminal place.
\end{corollary}
\begin{proof}
An orthonormal basis beginning with a basis of the kernel
contributes at least \(\dim\ker D_c\) to the squared Hilbert--Schmidt
norm. The preceding theorem gives the stated bound and full
vanishing on the indicated subspace. For a scalar lift \(v\)
the change of variables in Theorem~\ref{thm:graphtransport337}
gives \(\mathcal J_\phi(\widetilde R(v)A)
=\mathcal J_{\widetilde R(v^{-1})\phi}(A)=\mathcal J_\phi(A)\).
Summing a nontrivial character over its finite group gives zero.
\end{proof}

Here \(B_c/R_c=(T-1)^{-f}\). Thus even an arbitrary scalar
normalization cannot give a mean squared Gram error tending to
zero on these refined tests. If the prescribed row projection
\(e\) is applied to \(A_\lambda\) first and \(\lambda\) is nontrivial
at every place, equation~\eqref{eq:lift200} gives
\(\mathcal J_\phi(eA_\lambda)=\mathcal J_\phi(L_T A)\).
This is again invariant under terminal scalars on the column.
Equation~\eqref{eq:theta200} gives a positive lower bound for
the norm of the complex theta image of the raw signed test for
every nonzero \(A\) under the hypotheses there. The comparison
still requires the specified projection of the theta image.
Equation~\eqref{eq:bessel200} concerns the symmetric Bessel coefficient.
The exchange integral used above occurs only at the first position.

\begin{proposition}[A translation retaining the residue directions]\label{prop:firstcolumn337}
In the chosen split half, interchange \(e\) and the first flag
vector \(e_1\), using the inverse adjoint on the opposite half.
For the resulting translated graph lattices, the
\(q-1\) nonzero scalar multiples of any fixed residue covector
give distinct orbits under the full exchange group
\(\prod_{i=1}^s X_i(A_k)\). The scalar stabilizer of one of
these orbits inside \(O^\times\) is \(1+\pi O\).
\end{proposition}
\begin{proof}
Put \(B=A_k\oplus ke_1\). The common exchange group fixes \(B\)
pointwise. For any invertible \(P\) in it, one has
\(P\Lambda\cap B=\Lambda\cap B\) since \(P^{-1}\) fixes \(B\).
In the translated graph this intersection is precisely
\[
     \{w+t e_1:w\in O^a,\ t\in O,\ t\equiv\ell(w)\pmod\pi\}.
\]
Two such lattices in \(B\) coincide only if their covectors agree,
as is seen by testing every \(w\). A scalar \(v\in O^\times\)
changes the covector to \(\bar v^{-1}\bar\ell\); when \(v\equiv1\)
it fixes the entire lattice, which is a direct sum with the other
coordinate lines. This proves the assertion.
\end{proof}

This translation is available at each of the same terminal places.
Distinct exchange orbits distinguish the local lattice tests.
Their integrals against a global automorphic row may still agree
or vanish. A lower bound on the projected global pairing needs
information about those actual values.


\subsection{The first theta test at the translated terminal graph}\label{sec:pintable340}
Use the chosen half over $T$ of Theorem~\ref{thm:terminalwalk338}.
Write $q=T$, $\pi$ for its uniformizer and $O$ for its valuation ring.
The row lattice $\Lambda=\Lambda_{\ell_0}$ has rank $r=a+2s+1$.
Here $s\ge1$. The middle coordinates satisfy $t\equiv\ell_0(w)\bmod\pi$ for
$w+te\in O^a\oplus Oe$. Write $g$ for the first-position translate
in that theorem. We work in a basis of this particular $\Lambda$.

\begin{proposition}[Elementary divisors in the terminal graph frame]\leavevmode\label{prop:graphframe340}
The elementary-divisor exponents of $g\Lambda$ relative to $\Lambda$
are $-1$ with multiplicity $s+1$, $1$ with multiplicity one and
zero with multiplicity $r-s-2$. Put $h=s+1$, $K=\mathrm{GL}(\Lambda)$,
$S=K\cap gKg^{-1}$. Then
\[
       [K:S]=q^h {r\brack h}_q {r-h\brack 1}_q ,
\]
where ${n\brack k}_q$ counts $k$-dimensional subspaces of
$\mathbf F_q^n$. The $h$-dimensional subspace
$W=(\pi g\Lambda+\pi\Lambda)/\pi\Lambda$ contains the whole first
flag $E_s=\langle\bar e_1,\ldots,\bar e_s\rangle$.
\end{proposition}
\begin{proof}
Change the integral basis in the column space to make $\ell_0$ a
coordinate modulo $\pi$. Choose its last basis vector $w_0$ with
coordinate one. Alongside the other column vectors the graph has
basis $u=w_0+e$, $v=\pi e$. The translate interchanges $e_1,e$ and
then multiplies the first $s$ flag vectors by $\pi^{-1}$. On the
ordered triple $e_1,u,v$ its matrix is
\[
 \begin{pmatrix}0&\pi^{-1}&1\\0&1&0\\
                 \pi^{-1}&-\pi^{-1}&0\end{pmatrix}.
\]
The least valuation of an entry is $-1$, the least valuation of a
minor of size two is $-2$, and the determinant has valuation $-1$.
Thus the three exponents are $-1,-1,1$. The other $s-1$ first flag
vectors have exponent $-1$. All remaining vectors are fixed. The
reductions of $\pi g e_1$ and $\pi g u$ span $\bar v,\bar e_1$,
and $\pi g e_i=e_i$ for $2\le i\le s$. This proves the description
of $W$. In a Smith basis the reduction of $S$ stabilizes the flag
of dimensions $h,r-1$. Every element of this residue parabolic
lifts. Inside its full inverse image the last row of the first
$h$ columns must vanish also modulo $\pi^2$. There are $h$
independent extra residue conditions because the lifts of the
entries are free at the first congruence level. Counting the two successive
subspaces in the flag gives the index.
\end{proof}

Let $w\ge r$ be the next rank and $d'=w-r$. Write $Y=\binom{X}{Z}$
in the matrix model of Theorem~\ref{thm:tests165}. At depth $j$ put
$L=(B_v\cap K)K_r(j)$ for $j\ge1$, and $L=K$ for $j=0$. The local
tests $\phi_i$ are supported on $X\in Li$, $Z\in\pi M_{d',r}(O)$,
where $i\in L\backslash K$. Their nonzero values are $\theta(X)$,
with $\theta$ the compact character specified there. Use the partner
translate $b=\operatorname{diag}(g,I_{d'})$ in the same matrix
frame. Let $S=K\cap gKg^{-1}$ as above, and set
\[
         J=L\cap S,\qquad H=gLg^{-1}\cap S,\qquad C=J\cap H.
\]
Use the inner product linear in the first variable. All entries below
are divided by the same original squared norm $\|\phi_i\|^2$.

\begin{theorem}[The mixed Gram matrix on the original finite tests]\leavevmode\label{thm:pintable340}
Every nonzero entry of the matrix
$\langle\omega(g,b)\phi_k,\phi_i\rangle/\|\phi_i\|^2$ equals
\[
             m=\eta(g,b)\,q^{-d'(s+2)/2}[L:C]^{-1} ,
\]
where $\eta$ is the scalar character in the chosen splitting relative
to the matrix splitting. The active rows are indexed by $J\backslash S$
and the active columns by $H\backslash S$. Their entry is nonzero
exactly when the two cosets meet. In particular the fractional ideal
generated by all these entries is $(m)$ in any valuation ring of a
coefficient field containing $\sqrt q$ and the specified splitting value.
With the standard Euclidean norms on the coordinate vectors the matrix
has operator norm $|m|\sqrt{[J:C][H:C]}$ and squared Hilbert--Schmidt
norm $|m|^2[S:C]$.

At every depth $j\ge1$ its rank divided by $[K:L]$ is at most
\[
          \frac{{s+1\brack s}_q}{{r\brack s}_q} .
\]
At depth zero the unique entry has absolute value
$q^{-d'(s+2)/2}/[K:S]$. Translating the row by $g$ with the identity
on the partner gives a zero matrix at every depth.
\end{theorem}
\begin{proof}
In these coordinates the matrix action is
\[
 (\omega(g,b)\phi_k)(Y)=\eta(g,b)|\det g|^{w/2}|\det b|^{-r/2}
                                              \phi_k(b^{-1}Yg).
\]
In the normalization of \eqref{eq:weil161} the Jacobian of $Y\mapsto b^{-1}Yg$
is $|\det g|^w|\det b|^{-r}$, and its square root gives the unitary
matrix action. The top block changes by $X\mapsto g^{-1}Xg$ and
the lower block by $Z\mapsto Zg$. The character $\theta$ on $K$
factors through determinant. Indeed the elementary matrices generate
$\mathrm{SL}_r(O)$ by elementary row reduction over the valuation
ring. Each elementary matrix is a commutator with an integral diagonal
unit whose difference from one is a unit, since $q>2$. Thus the top
character cancels on the common support.

For an exponent $t$ of $g$, the normalized overlap on a single lower
coordinate is $q^{-|t|/2}$. In fact the support intersection has
relative measure $q^{-\max(0,-t)}$ and the square root of the Jacobian
contributes $q^{-t/2}$. Changes of integral frames preserve the
lattice and measure. By the proposition, over all lower rows the
absolute factor is $q^{-d'(s+2)/2}$. On the top block the common
support is $Li\cap gLkg^{-1}$ inside $S$. Every nonempty intersection
of $S$ with $Li$ is one coset of $J$. The corresponding intersections
for the translated column cosets are exactly those of $H$. When the
two cosets meet they meet in one coset of $C$. On the matrices of
unit determinant absolute value, additive measure is proportional
to multiplicative Haar measure. The ratio to the volume of $L$ is
therefore $[L:C]^{-1}$. The identity cosets meet, giving the ideal
assertion including at a prime dividing $[L:C]$.

Each active row meets $[J:C]$ active columns, and each active column
meets $[H:C]$ active rows. Applying Cauchy--Schwarz to each row and
then summing bounds the operator norm by the stated expression.
The vector constant on the active columns and zero elsewhere attains
the bound. The nonzero entries are indexed by $C\backslash S$,
which gives the squared Hilbert--Schmidt norm.

The compact Bessel group preserves the first isotropic flag, and
$K_r(j)$ fixes its reduction for $j\ge1$. Consequently $i^{-1}E_s$
is well defined on $Li$. If $Li$ meets $S$, represent it by an element of $S$. It preserves
$W$, so the inclusion in the proposition gives $i^{-1}E_s\subset W$. As $Li$ ranges over the finite table, $i^{-1}E_s$
takes each $s$-plane equally often by the transitivity of $K$ modulo
$\pi$. Only ${s+1\brack s}_q$ of the ${r\brack s}_q$ planes lie in
$W$. This bounds the number of active rows and hence the rank. For
$L=K$ one has $J=H=C=S$. Finally a row translation alone requires
both $X$ and $Xg$ to lie in $K$. Since $g\notin K$, these two support
conditions are disjoint.
\end{proof}

For tensor products over the $f$ places above $T$, the local matrices
tensor, their ideals multiply, and the rank fractions multiply. The
result concerns the arithmetic matrix of overlaps of the actual
prescribed local tests. Estimating the global signed projection in
Section~\ref{sec:missing} also requires its values on the automorphic
forms obtained by integrating the theta kernels.


\subsection{The auxiliary character on the original signed projector}\label{sec:signedreturn342}
At a half over $g_1$ or $g_2$ use the signed subgroup $\Gamma$,
character $\lambda_\Gamma$ and ordinary Bessel character $\chi$ of
Lemma~\ref{lem:auxiliaryroots339}. Put $P_\lambda=\int_\Gamma
\overline{\lambda_\Gamma(k)}R(k)\,dk$ with probability Haar measure.
Assume $s\ge1$ so that there is a simple string subgroup $Z_j$ as
in that lemma.

\begin{theorem}[Vanishing on returning to the original signed type]\leavevmode\label{thm:signedreturn342}
For the unmodified absolute-trace character the ordinary global
Bessel period of every row in the image of $P_\lambda$ is zero at
every column. Furthermore, for either auxiliary diagonal $t$ of
Theorem~\ref{thm:auxiliaryflag339} or Proposition~\ref{prop:centeredflag341},
\[
                        P_\lambda R(t)P_\lambda=0 .
\]
This identity holds in every unitary representation of the row
group, including on projected automorphic theta forms. It allows
arbitrary simultaneous operations commuting with this local group.
\end{theorem}
\begin{proof}
By Lemma~\ref{lem:auxiliaryroots339}, for any simple string root
one can choose $w\in O$ with $u=Z_j(w)\in\Gamma$ and
$\lambda_\Gamma(u)\ne1$. The absolute-trace character has conductor
ideal $O$ on this root, so $\chi(u)=1$. The column group acts on
$A$ and commutes with $u$, which involves only the string. In the
ordinary Bessel integral translate the unipotent variable on the
right by $u$ at this place. The quotient measure is right invariant
and this change multiplies the character by one. Thus replacing a
row $F$ by $R(u)F$ leaves its period fixed. For a signed row it
also multiplies it by $\lambda_\Gamma(u)$. This proves the first
assertion, with the other local and archimedean factors kept fixed.

Both diagonals satisfy $v(l_{j+1}/l_j)=1$. Consequently
$t^{-1}ut=Z_j((l_{j+1}/l_j)w)$ is in the first principal congruence
subgroup and has trivial compact character. It follows that
$R(u)R(t)P_\lambda=R(t)P_\lambda$. On the other hand
$P_\lambda R(u)=\lambda_\Gamma(u)P_\lambda$, by changing the variable
in the compact average. Applying these two equalities to the same
vector proves the operator identity. A determinant twist has no
effect on $u$ or its conjugate. The argument works in the actual
automorphic representation before any subsequent linear period is
applied. Thus this projector annihilates the translated signed row
even when the local principal matrix indicators in
Theorem~\ref{thm:auxoverlap341} have positive overlap.
\end{proof}


\subsection{The integral map at the last hyperplane}\label{sec:lowerlattice347}

At one embedding let \(N=n+1\), with \(n\ge1\). Write the two
dominant weights as \(\lambda\) and \(\mu\), in decreasing order.
Suppose \(\lambda_i\ge\mu_i\ge\lambda_{i+1}\) for \(1\le i\le n\)
and \(s=\lambda_N=\mu_n\). Retain the determinant lines separately.
In the resulting polynomial models the partitions are
\(l_i=\lambda_i-s\) and \(m_i=\mu_i-s\). For a partition \(v\)
let \(p_v\) be its product of leading minors and put
\(H_v=\|p_v\|_{\rm F}^2\). Use the rational bilinear form
\[
 B_v(P,Q)=H_v^{-1}\sum_\alpha \alpha! P_\alpha Q_\alpha .
\]
Monomials have their usual (undivided) coefficients. The
corresponding positive Hermitian norm on complex polynomials uses
complex conjugates in the second slot. The lattice in a frame is
the intersection of the flag polynomial space with the integral
monomial lattice, as in Theorem~\ref{thm:terminaldual325}.

Let \(D_i\) be the leading minor of order \(i\) of an
\(N\)-by-\(N\) variable matrix. Let \(D_i'\) use the first
\(i\) columns and rows \(1,\ldots,i-1,N\), in that order. Put
\[
 a_i=\mu_i-\lambda_{i+1},\qquad b_i=\lambda_i-\mu_i,\qquad
 X=\prod_{i=1}^n D_i^{a_i}(D_i')^{b_i},\qquad
 \kappa=B_l(X,X).
\]
For an \(N\)-square matrix put its last column equal to zero and
extract the coefficient of \(\prod_{i=1}^n X_{Ni}^{b_i}\) in the
last row. Denote this linear operation by \(C_b\), using \(X_{Ni}\)
just for the matrix entries in this expression. Let \(P_m\) be
Fischer orthogonal projection onto the smaller flag polynomial
space of partition \(m\). Define
\[
 d_0=\sum_i a_i,\quad
 I_i=[\lambda_{i+1},\min(\lambda_i,\lambda_{i+1}+d_0)]\cap\mathbb Z,
 \quad
 c=\prod_i\prod_{v\in I_i\setminus\{\mu_i\}}\prod_{j=1}^n
                 (v-i-\mu_j+j).
\]
With the normal line in its standard frame, the prescribed
restriction on polynomials is \(L=c\kappa P_m C_b\).
These are the constants in
\srcloc{Formalism/LocalToGlobal/Columns/FiniteSumNonzeroLargeDepth}{53}
and lines 113--138 there. We compute their effect on the lattice.

\begin{lemma}[The squared length of the interlacing line]\leavevmode
\label{lem:highestline347}
Put \(h=\lambda_1-\lambda_N\). The rational number \(\kappa\)
is positive, and both \(c\) and \(\kappa\) are units in
\(\mathbb Z_{(p)}\) for every rational prime \(p>h+n\).
Moreover \(C_bX=p_m\).
\end{lemma}
\begin{proof}
In extracting the last-row power, the variable in column \(n\) must
come from every modified minor of size \(n\). Their cofactors are
\(D_{n-1}\) with positive sign. Continuing in decreasing order
forces the last-column term of each remaining modified minor.
The result is \(\prod_iD_i^{a_i}D_{i-1}^{b_i}=p_m\), with
\(D_0=1\). The diagonal monomial of \(p_m\) has coefficient one.
Consequently \(X\) is primitive in the integral polynomial lattice.

The flag minors place \(X\) in the large module of partition \(l\).
Under row changes on the first \(n\) coordinates it is a highest
vector. Each factor is unchanged by adding an earlier row to a
later one among these coordinates. Its row weight is \(m\), as is
also seen by counting the rows in the factors. Restriction has
multiplicity one by the interlacing rule of
Proposition~\ref{prop:branch275}. Thus the subspace of weight \(m\)
in the \(m\)-isotypic summand is this one line.

Over \(\mathbb Z_{(p)}\) the large flag lattice is self-dual for
\(B_l\) by Theorem~\ref{thm:terminaldual325}. Indeed every hook
of \(l\) has length at most \(h+n\).
The isotypic projector for the row block is integral there. To
see this directly on polynomials, use the single block of size \(n\)
in \eqref{eq:sieve}. Its matrix generators are sums of operators
of the form \(X_{ij}\partial_{X_{kj}}\), up to transpose and
the separate integral determinant shift. They preserve integral
polynomials. The denominator has prime factors at most \(h+n-1\).
The isotypic projector is orthogonal for the compact unitary
action and rational by that formula. Its matrix therefore is
self-adjoint also for the rational bilinear form. Extraction of
the row weight \(m\) is an integral self-adjoint projection
by monomials. These projections commute. Their product splits
off the highest line orthogonally. An orthogonal direct summand
of a self-dual lattice is self-dual, since a functional on it
extends by zero on its complement. Its lattice is generated
by \(X\), by primitivity and its coefficient one. Its Gram
entry \(\kappa\) is consequently a unit at \(p\).
Positivity follows from the real Fischer form and \(X\ne0\).

For the integer \(c\), the element \(\mu_i\) belongs to \(I_i\)
because \(a_i\le d_0\). If \(j<i\), then
\(\mu_j\ge\lambda_i\ge v\), and if \(j>i\), then
\(\mu_j\le\lambda_{i+1}\le v\). In either case
\(v-i\ne\mu_j-j\). The factor with \(j=i\) is nonzero by
the exclusion of \(\mu_i\). All absolute values are at most
\(h+n-1\), proving the assertion also for the empty product.
\end{proof}

\begin{theorem}[A split surjection in the prescribed normalization]\leavevmode
\label{thm:lowerlattice347}
Let \(L^*\) be the adjoint between the two normalized bilinear
spaces. On the smaller space
\[
                 LL^*=c^2\kappa\,1 .
\]
The nonzero singular values of \(L\) in the two normalized
Hermitian norms are all \(|c|\sqrt{\kappa}\).
At every prime \(p>h+n\) the map on the indicated polynomial
lattices is a split surjection with integral section
\(L^*/(c^2\kappa)\). These assertions tensor over embeddings
and extend to the local rings of a coefficient number field.
\end{theorem}
\begin{proof}
Write \(A=P_m C_b\). Row changes in the first \(n\) coordinates
commute with coefficient extraction, and the projection commutes
with their compact group. Its complexification gives equivariance
for \(\mathrm{GL}_n\). The two separated determinant twists agree
on this block. The adjoint also intertwines, first for the compact
group in the Hermitian metrics, then for its complexification.
All matrices in the polynomial bases are real and rational, so
the bilinear adjoint gives the same matrix. The highest line
computed in Lemma~\ref{lem:highestline347} forces
\(A^*p_m=tX\). Since \(AX=p_m\) and \(B_m(p_m,p_m)=1\),
pairing with \(X\) gives \(t=1/\kappa\).
The endomorphism \(AA^*\) is scalar on the irreducible target.
Its value on \(p_m\) is \(\kappa^{-1}p_m\).
Multiplication by \(c\kappa\) proves the identity and the
assertion on singular values.

For the lattice assertion, \(C_b\) takes integral coefficients
to integral coefficients. Pairing such an output with any integral
target polynomial by the Fischer form divided by \(H_m\) is
integral at \(p\), since \(H_m\) is a unit there. Projection
preserves these pairings. Self-duality of the target now shows
that the projected output is integral. Both source and target
are self-dual, so the adjoint preserves their lattices as well.
The scalar in the identity is a unit by the lemma. It gives the
stated section. Taking products of the sections and extending
scalars proves the last assertions.
\end{proof}


\subsection{One integer for the terminal normalizations of all selected words}\label{sec:wordlower348}

Use the common fields, weights and intervals of
Theorems~\ref{thm:common111} and~\ref{thm:markedstrict337}.
At a fixed \(x\), a translate fixes the small terminal row and
its weights. A binary word \(w\) determines the large weights
by \eqref{eq:letters111}. Denote by \(U_w\) the nonzero rational
number in Theorem~\ref{thm:terminalformula347}, and by
\(Z_w\) its normal-coordinate factor for the chosen \(\pi\).

\begin{theorem}[Finite and archimedean variation in the last restriction]\leavevmode
\label{thm:wordlower348}
There is an integer \(J_x\ge1\), with \(\log J_x\ll x^C\)
for a fixed \(C\) uniformly on the common intervals, such that
for every translate and every two words at that translate both
the numerator and denominator of \(|U_w/U_{w'}|\), in lowest
terms, divide \(J_x\). At the three marked primes the ratio is
a unit. In the selected family at its common prime \(T\) and
generator \(\pi\), the same assertions hold for
\((Z_wU_w)/(Z_{w'}U_{w'})\).
\end{theorem}
\begin{proof}
First fix one translate and one embedding \(\tau\). Suppress
that index and write \(N=n+1\) for the terminal rank.
In the explicit lists of Corollary~\ref{cor:normaldegree347},
the repeated tail, the first and last core entries, the rank
shifts, and the entire small list \(\mu\) are fixed.
Only \(\lambda_2,\ldots,\lambda_{k-1}\) change. Each has two
possible values differing by \(\Delta\), the spacing in
\eqref{eq:letters111}. In particular the common determinant
shift \(s=\lambda_N=\mu_n\) is fixed.

Form a box of the exponents
\(a_i=\mu_i-\lambda_{i+1},\ b_i=\lambda_i-\mu_i\) using their
coordinatewise minima and maxima over all words. There are at
most \(2k\) varying exponents, each of range at most \(\Delta\).
Consequently the difference of the total degrees of the two
extremal products is at most \(2kN\Delta\). Their upper degree
is at most \(N\max_w(\lambda_1-\lambda_N)+2kN\Delta\).
All of \(N,k,\Delta\) are polynomial in \(x\), and
\(\log(2+\max_w(\lambda_1-\lambda_N))\ll x^\eps+\log(2x)\)
by the explicit terminal lists. Theorem~\ref{thm:commonfischer348}
now gives one integer of logarithm \(O(x^C)\) whose divisors
represent every ratio of the squared norms \(q(X_w)\).
Apply the same theorem to the leading-minor products
\(p_{l_w}=\prod D_i^{\lambda_i-\lambda_{i+1}}\).
Here at most \(2k\) exponents change, by at most \(2\Delta\).
This provides the same type of estimate for
\(H_{l_w}=q(p_{l_w})\). Since \(\kappa_w=q(X_w)/H_{l_w}\),
the product of these two integers clears the numerator and the
denominator of every \(\kappa_w/\kappa_{w'}\).

It remains to treat the integer \(c_w\). For each \(i\) all its
intervals \(I_i(w)\) from Section~\ref{sec:lowerlattice347}
contain \(\mu_i\). Their lower endpoints vary by at most
\(\Delta\). Their upper endpoints vary by at most
\((k+1)\Delta\), since \(d_0=\sum_i(\mu_i-\lambda_{i+1})\)
varies by at most \(k\Delta\) and taking a minimum cannot
increase the largest of the two variations. Delete \(\mu_i\)
from every interval. The factors for its common intersection
cancel in all ratios. In the union outside the intersection,
at most \((k+2)\Delta+2\) integers \(v\) remain per index.
For each of these \(v\) and each \(1\le j\le n\) form the
absolute value of \(v-i-\mu_j+j\). They are positive and at
most \(n+\max_w(\lambda_1-\lambda_N)\) by the proof of
Lemma~\ref{lem:highestline347}. Their product accounts for
every remaining factor with its multiplicity and has logarithm
\(O(x^C)\). Both parts of each reduced ratio of the \(c_w\)
divide this common integer.

The target hook factor \(H_{m_\tau}\) is fixed within a
translate. Multiply the three integers just obtained, then
multiply over the CM embeddings and the translates. Both finite
index sets have size at most \([E_x:\mathbb Q]\), polynomial
in \(x\) by Proposition~\ref{prop:field12}. If the two parts
of each reduced fraction divide \(J_i\), both parts of their
product divide \(\prod_iJ_i\), by adding prime valuations.
This proves the claimed bound with a common \(J_x\) even when
the word family has exponential size. The marked-prime
assertion follows from Theorem~\ref{thm:terminalformula347}.
For the selected family \(Z_w=Z_{w'}\) by
Corollary~\ref{cor:normaldegree347}.
\end{proof}

\begin{corollary}[The same factor in every exterior minor]\leavevmode
\label{cor:wordplaces348}
At one translate fix a selected reference word \(w_0\). Take
any \(a\ge1\) selected columns labelled \(w_1,\ldots,w_a\),
and any \(a\) tests on their common small row. Let \(M\) be the
matrix of their prescribed periods, and \(S\) the matrix of the
weighted coefficient sums of Theorem~\ref{thm:terminalformula347},
including the same stabilizer weights. Then
\[
  \frac{\det M}{(Z_{w_0}U_{w_0})^a}
     = R\det S, \qquad R=\prod_{j=1}^a U_{w_j}/U_{w_0}.
\]
The rational scalar obeys
\[
 \frac1a\sum_p |v_p(R)|\log p\le\log J_x,\qquad
 \frac1a|\log|R||\le\log J_x.
\]
Its valuations at \(g_1,g_2,T\) vanish. All these identities
apply to the original signed columns and tests.
\end{corollary}
\begin{proof}
The period formula multiplies an entire column by its same scalar
for every test. Factoring the determinant gives the identity,
including for dependent columns. At any \(p\), every individual
valuation difference lies in \([-v_p(J_x),v_p(J_x)]\).
Summing this bound over the \(a\) factors and then the primes
gives the first estimate. A positive fraction whose two reduced
parts divide \(J_x\) lies between \(J_x^{-1}\) and \(J_x\),
giving the second estimate. The marked-prime assertion follows
from the theorem. In particular these relative logarithmic costs
per column are \(o(b/x^A)\) for every fixed \(A>0\), since
\(b=\exp(x^\eps)\).
\end{proof}
The scalar for one reference word and the determinant of the
actual coefficient sums still appear in an absolute estimate.
The latter depends on the incoming rows and on cancellation
between the prescribed classes at the same translate.

\section{Local divisors and archimedean scales of the projected determinant}\label{nav:terminal-determinants}

\subsection{The directional period estimate at every finite prime}
\label{sec:allprimetable862}
Keep the fields, old marks, larger-weight embeddings and main
pairs of Section~\ref{sec:shortpairs859}. Write \(S_p=0\) at a
split prime and use the sum of that section otherwise. For the
given \(d=(d_\tau)_\tau\), let
\[
 n=\lceil4b_x/k\rceil,\quad D_\tau=(d_\tau-n)_+,\qquad
 B_\tau=\prod_{i=1}^n(D_\tau+i),\qquad A_d=\prod_\tau B_\tau^2.
\]
Every \(B_\tau\) is divisible by \(n!\), by the binomial formula.
On all sufficiently distant intervals
\[
         \log A_d\ll_{q_0,d_{\rm K}} f b_x\log(2b_x)/k .
\]
Indeed \(d_\tau\le2b_x\) by the derivative window in the proof
of Theorem~\ref{thm:shortpairs859}, \(n=O(b_x/k)\) and every
factor in the products is at most \(2b_x+n\).

\begin{lemma}[A common normalization at the original marked places]
\leavevmode\label{lem:idealscale862}
There is a nonzero algebraic number \(b\) with absolute value
\(M^{h/2}\) at every complex embedding and with
\[
         v_\iota(b)=\sum_\tau d_\tau v_\iota(\tau\mathfrak m)
\]
at every finite coefficient embedding, where \(h=\sum_\tau d_\tau\).
\end{lemma}
\begin{proof}
View the oriented embeddings as automorphisms of the abelian
field \(E\) fixing \(K\). In \(E\) form
\(\mathfrak a=\prod_\tau\tau(\mathfrak m)^{d_\tau}\).
Since conjugation commutes with these automorphisms,
\(\mathfrak a\bar{\mathfrak a}=(M)^h\). Choose \(N\ge1\)
and \(w\in E^\times\) with \(\mathfrak a^N=(w)\) by
finiteness of the ideal class group. Take any root satisfying
\(b^{2N}=M^{hN}w/\bar w\). The ideal on the right is
\(\mathfrak a^{2N}\), so taking local valuations gives the
display. At any complex embedding of \(E\) the image of
\(w/\bar w\) has absolute value one. Taking absolute values
in the defining equation gives the stated value for every
conjugate of the root; in particular \(b\bar b=M^h\).
\end{proof}

\begin{theorem}[The derivative table at all finite primes]
\leavevmode\label{thm:allprimetable862}
Use the \(b\) of Lemma~\ref{lem:idealscale862} in
\eqref{eq:derivtable855}. Put \(T_d=\gamma_x A_d\), using
the nonzero algebraic integer of Theorem~\ref{thm:integralleading857}.
For every main pair on the sufficiently distant common intervals
and every coefficient embedding above every prime \(p\),
\begin{equation}
                     e_\iota\le v_\iota(T_d)+S_p .
 \label{eq:allprimetable862}
\end{equation}
In particular \(T_d\) is independent of \(p,\iota\) and of the
prescribed twist.
\end{theorem}
\begin{proof}
At a split prime use Corollary~\ref{cor:splitperiod858}.
Its correction from \(\mathfrak m\) cancels \(v_\iota(b)\)
by the lemma. The other ideal correction is nonpositive,
so \(e_\iota\le v_\iota(\gamma_x)\). At all nonsplit primes
\(b\) is a unit because its ideal is supported only above the
split marks. At the ramified prime use Theorem~\ref{thm:shortpairs859},
where \(E_p=0\) and \((n!)^f\) divides \(A_d\).
Here \(p=d_{\rm K}\le k\) once the interval is sufficiently
distant. We treat the two remaining ranges of inert primes.

Suppose first that \(p\le k\). In the notation of
Theorem~\ref{thm:shortpairs859} use every pair
\(\tau\mapsto t_\tau=P_{a_\tau}\tau\) occurring in \(E_p\)
and put \(j_\tau=\lfloor d_\tau/p^{a_\tau}\rfloor\).
The ends and starts are all disjoint, and the smallest-exponent
condition in Section~\ref{sec:pairedtable861} holds.
The power \(p^{a_\tau}j_\tau\) fits at its start, and
\(j_\tau\le(b_x+W)/p\le b_x-W\le d_{t_\tau}\) fits
at its end; here \(W=O(x^{42})\) bounds
\(|d_\tau-\kappa_\tau|\), and \(W/b_x\) tends to zero.
Thus \(\eta=d-r-s\ge0\) componentwise. With \(J=\sum j_\tau\)
we have \(E_p-J\le f(1+W/p)\), since each start occurs
at most once and has weight \(b_x\).

Legendre's digit formula applied to a three-part multinomial
at each of the \(f\) embeddings gives
\[
           v_p\bigl(d!/(r!s!\eta!)\bigr)
                  \le 3f\bigl(1+\log_p(2b_x)\bigr).
\]
In fact each of the three lower integers is at most \(2b_x\),
and its base-\(p\) digit sum is at most
\((p-1)(1+\log_p(2b_x))\). The upper integer's digit sum is
subtracted. From the geometric-series calculation in the
proof of Theorem~\ref{thm:shortpairs859} we have, before
absorbing any error into a factorial,
\[
  \lambda_p h-S_p-E_p\le f\left(LG+
          \frac{b_x}{(p-1)p^L}+\frac{W}{p-1}\right).
\]
Here \(L\) is as in that theorem. Substitution in
\eqref{eq:secondpair861} shows that the excess over
\(v_\iota(\gamma_x)+S_p\) is bounded per embedding by
\[
    LG+\frac{b_x}{(p-1)p^L}+\frac W{p-1}
                         +1+\frac Wp+3(1+\log_p(2b_x)).
\]
The middle geometric tail is at most \(b_x/(kp)\). The other
terms are together \(o(b_x/(kp))\) uniformly here. For the
first and third terms this was checked in the preceding
theorem on the small primes. For the additional terms use
\(p\le k\), \(Wk/b_x\to0\) and \(k^2\log(2b_x)/b_x\to0\).
The last limit follows from \(k\le3x^5,\ b_x\ge x^{400}\)
and monotonicity of \(\log(2y)/y\) for large \(y\).
Since \(b_x/k^2\to\infty\), the display is at most
\(\lfloor n/p\rfloor\le v_p(n!)\). This proves the stated
bound for all these inert primes, including those at which
\(E_p\ne0\).

Finally let \(p>k\) be inert. Use every oriented embedding
with exponent one, and put
\[
        i_\tau=\lfloor D_\tau/p\rfloor,\qquad
        j_\tau=\min(i_\tau,d_{P_1\tau}).
\]
Here \(i_\tau\le 2b_x/p\le n\). The power entries
\(p j_\tau\) at \(\tau\) are bounded by \((d_\tau-n)_+\),
and the entry arriving at \(\tau\) is bounded by
\(\min(n,d_\tau)\), since \(P_1\) is bijective. Their sum is
at most \(d_\tau\), even when the permutation fixes \(\tau\).
Apply \eqref{eq:firstpair861}. For each \(\tau\) one has
\[
    v_p(d_\tau!)\le v_p(B_\tau)+v_p(D_\tau!)
       =v_p(B_\tau)+i_\tau+v_p(i_\tau!)
       \le 2v_p(B_\tau)+i_\tau .
\]
The first inequality is equality when \(d_\tau\ge n\);
when \(d_\tau<n\), \(d_\tau!\) divides \(B_\tau=n!\).
The middle equality uses the multiples of \(p\) in \(D_\tau!\).
The last inequality uses \(i_\tau!\mid n!\mid B_\tau\).

The actual orders satisfy
\(\kappa_\tau\le d_\tau\le\kappa_\tau+W\).
For the lower bound the shift in the main-pair window is
\((2j+1)s^8-h_\tau-z\) with \(0\le h_\tau\le4s,\ z\le s\);
it is positive for these ranks. Also \(n\ge W\) eventually
by the scale bounds. Hence \(i_\tau\le\kappa_\tau/p\) and
\(d_{P_1\tau}\ge\kappa_{P_1\tau}\). It follows that
\[
        i_\tau-j_\tau
                  \le(\kappa_\tau/p-\kappa_{P_1\tau})_+ .
\]
The sum of the right sides is the first-exponent contribution
to \(S_p\), and all its other terms are nonnegative.
Summing the factorial inequalities after subtracting \(j_\tau\)
proves the required assertion in the last range.
\end{proof}

\begin{corollary}[The same-field comparison for the projected period]
\leavevmode\label{cor:allprimeperiod862}
On these same fields and intervals take the simultaneous local
cohomology and critical twists of Corollary~\ref{cor:actualdenom803},
and use the marked normalization of Section~\ref{sec:twelftransfer855}
with \(b\) from Lemma~\ref{lem:idealscale862}. For every
allowed \(p,\iota,\beta\), with \(p\nmid Q_\beta\),
\begin{equation}
  -v_\iota\!\left(H_{\nu,z}\,T_d\overline{T_d}\,
                    \rho_\beta^2 R_{\nu,z}(\beta)^{-1}\right)
                          \le p_{p,\iota}+2S_p .
 \label{eq:allprimeperiod862}
\end{equation}
The common multiplier satisfies
\(\sup_\sigma\log^+|\sigma T_d|
            \ll_{q_0,d_{\rm K}} f(1+b_x/k)\log(2b_x)\).
Even summed over \(mk\) pairs its logarithmic cost divided by
\(sB_x^{\rm det}\) is \(O((\log b_x)^{-2})\).
\end{corollary}
\begin{proof}
Use \eqref{eq:allprimetable862} at \(\iota\) and
\(\iota\circ c\), with \(c\) complex conjugation. The same
\(T_d\) and \(S_p\) occur at both. Substitute in
\eqref{eq:tablevaluation855}; moving \(v_\iota(T_d\bar T_d)\)
inside the negative valuation gives \eqref{eq:allprimeperiod862}.
The cohomological term \(p_{p,\iota}\) is the one for these
very tests and twists there. The archimedean bound follows
from the estimate for \(A_d\) and the conjugate bound for
\(\gamma_x\) in Theorem~\ref{thm:integralleading857}.
The final assertion follows by substituting
\(B_x^{\rm det}=fxb_x(\log b_x)^3/(mk)\) and
\(m^2\ll x,\ k\le s\).
\end{proof}

Equation~\eqref{eq:allprimeperiod862} bounds the finite valuations
for the marked high-weight period. To use it in the terminal
determinant one also needs the archimedean size of the inverse
untwisted norm contained in \(H_{\nu,z}\), with the same \(b\).
\clearpage

\subsection{The relative different and the archimedean period scale}
\label{sec:jetscale863}
Fix a node in the shorter family of Section~\ref{sec:allprimetable862}.
Write \(f=[E^+:\QQ]\), \(d_{\rm K}=|\operatorname{disc}(K)|\),
\(\xi\) for the quadratic character of \(K\), and \(c\) for complex
conjugation at the distinguished complex embedding. Put
\[
 D_E=N_{E/\QQ}(\mathfrak D_{E/K})^{1/(2f)},\qquad
 \Omega_E=\frac{2\pi}{\operatorname{rd}(E)}
                  \exp\!\left(\gamma-\frac{L'(1,\xi)}{L(1,\xi)}\right).
\]
The tower discriminant formula gives
\(D_E=\operatorname{rd}(E)/\sqrt{d_{\rm K}}\).
Let \(J\subset\overline{\QQ}\) be a number field of definition for
the marked discriminant-one elliptic curve and its differential.
We may take \(J\) Galois over \(\QQ\). For an embedding
\(\rho:J\to\CC\), let \(A_\rho\) be the area of its period lattice
in the given differential and \(q_\rho=\pi/A_\rho\). Complex
conjugating both the curve and the differential conjugates the
period lattice, so \(q_{c\rho}=q_\rho\). By
Theorem~\ref{thm:mean56}, with the mean taken over all embeddings,
\begin{equation}
 \log D_E+\mathbb E_\rho\log q_\rho=-\log\Omega_E .
 \label{eq:areamean863}
\end{equation}

\begin{lemma}[A common unit for the area factors]
\leavevmode\label{lem:unitarea863}
For any \(H\ge0\) there is an algebraic unit \(u\) such that every
embedding \(\sigma:J(u)\to\CC\), restricting to \(\rho\), satisfies
\[
 \left|\log|\sigma u|-\tfrac12
        \big(\mathbb E_\eta\log q_\eta-\log q_\rho\big)\right|
            \le \frac{1}{2(H+1)} .
\]
\end{lemma}
\begin{proof}
The prescribed vector on the embeddings of \(J\) is real, is
constant on each conjugate pair and has sum zero. Dirichlet's
unit theorem says that the logarithms \((\log|\rho v|)_\rho\)
of units \(v\) of \(J\) span this real vector space. Here both
members of a complex pair are counted. Express the vector as a
real linear combination of finitely many such vectors. Approximate
their coefficients by rationals until each coordinate error is at
most \(1/(2(H+1))\). Write these rationals as \(m_i/N\) with
\(N>0\). Take any root \(u\) of \(u^N=\prod_i v_i^{m_i}\).
Both \(u\) and \(u^{-1}\) are integral. At every extension of
\(\rho\) its log modulus is exactly
\(\sum_i(m_i/N)\log|\rho v_i|\). This also covers the zero
dimensional logarithm space by taking \(u=1\).
\end{proof}

\begin{theorem}[A simultaneous finite and archimedean normalization]
\leavevmode\label{thm:jetscale863}
For all main pairs at this node choose a common unit as in
Lemma~\ref{lem:unitarea863} with \(H=\max_{\mathbf d}\sum_\tau d_\tau\).
There are algebraic numbers \(b_{\mathbf d}\ne0\), independent of
the critical twist, such that the table and projected period
inequalities \eqref{eq:allprimetable862} and
\eqref{eq:allprimeperiod862} hold with \(b=b_{\mathbf d}\) in the
definition of the table, the lift, and \(H_{\nu,z}\). The cohomology
bound and \(T_d\) are those of Corollary~\ref{cor:allprimeperiod862}.
For every \(\sigma\in\operatorname{Aut}(\CC)\) let
\(\sigma'=c\sigma c\), \(h=\sum_\tau d_\tau\), and use \(q_\sigma\)
for \(q_{\sigma|J}\). Then
\begin{equation}
 \left|\log\!\left(
    |\sigma b_{\mathbf d}|\,|\sigma' b_{\mathbf d}|\,
           (q_\sigma q_{\sigma'})^{h/2}\right)
               -h\log(M/\Omega_E)\right|\le 1 .
 \label{eq:jetpair863}
\end{equation}
\end{theorem}
\begin{proof}
Put \(\mathfrak j=\mathfrak m\overline{\mathfrak b}\) using the
integral ideals of Corollary~\ref{cor:splitperiod858}. One has
\(\mathfrak j\bar{\mathfrak j}=(M)\mathfrak D_{E/K}\).
For a rational prime \(p\) let \(r_p\) be the exponent of the
different at a prime above \(p\), and \(e_p\) its absolute
ramification index. These integers are independent of the choice
of prime because \(E/\QQ\) is Galois and \(K\) is stable. The prime
ideal norm formula gives \(D_E=\prod_p p^{r_p/e_p}\). Choose a
positive integer \(N\) with \(\mathfrak j^N=(w)\) and
\(e_p\mid N\) at every prime in this finite product. This is
possible by finiteness of the ideal class group. Then \(D_E^N\)
is a positive rational number and
\((D_E^N)\mathcal O_E=\mathfrak D_{E/K}^N\).

For each oriented embedding, viewed as an automorphism \(\tau\)
of \(E\) fixing \(K\), take any algebraic root
\[
             a_\tau^{2N}=(MD_E)^N\,\tau(w/\bar w).
\]
The fractional ideal of the right side is
\((\tau\mathfrak j)^{2N}\), since conjugation commutes with \(\tau\).
Thus \(v_\iota(a_\tau)=v_\iota(\tau\mathfrak j)\) for all normalized
finite valuations, including after extension. The conjugates of
\(w/\bar w\) have modulus one, also by commutativity in
\(\operatorname{Gal}(E/\QQ)\). Hence
\(|\sigma a_\tau|=(MD_E)^{1/2}\) at every complex embedding.
Take \(b_{\mathbf d}=u^h\prod_\tau a_\tau^{d_\tau}\).

At a split prime the full correction in
Corollary~\ref{cor:splitperiod858} is
\(-\sum_\tau d_\tau v_\iota(\tau\mathfrak j)\). This cancels
\(v_\iota(b_{\mathbf d})\). At a nonsplit prime the old marks
are units and the relative extension is unramified there.
Indeed the conductor primes of the auxiliary cyclotomic factor
split in \(K\), as in Section~\ref{sec:shortpairs859}, and base
change to \(K\) preserves unramifiedness elsewhere. Both integral
ideals in \(\mathfrak j\) are then units at the prime.
Consequently the proof of Theorem~\ref{thm:allprimetable862}
gives the identical bound there. Applying the table inequality
at \(\iota\) and \(\iota c\) as in
Corollary~\ref{cor:allprimeperiod862} proves the finite assertion.

For each \(\sigma\), the preceding modulus calculation and the
unit lemma give
\[
 \left|\log\big(|\sigma b_{\mathbf d}|q_\sigma^{h/2}\big)
   -\tfrac h2\big(\log(MD_E)+\mathbb E_\rho\log q_\rho\big)\right|
          \le h/(2(H+1)).
\]
Apply this at \(\sigma\) and \(\sigma'\), and add. Substitution of
\eqref{eq:areamean863} proves \eqref{eq:jetpair863}.
\end{proof}

\begin{corollary}[The untwisted norm at the distinguished embedding]
\leavevmode\label{cor:identitynorm863}
Take \(n=2z\) and \(z\ge2\) in the main-pair range. Let
\(\phi_{\nu,1}\) be the untwisted high-weight lift scaled by
\(b_{\mathbf d}a_n(\pi/A)^{h/2}\). At the same rank and with
the same old column take a degree-zero pair
\(\eta=\psi\upsilon_1^{-z}\) as in Theorem~\ref{thm:flattorus853},
scaled by \(a_n\). At the distinguished complex embedding,
\begin{align}
 \left|\log\frac{\|\phi_{\nu,1}\|_{\rm cl}^2}
                     {\|\phi_{\eta,1}\|_{\rm cl}^2}
       -\log\dim V_{\mathbf d}
       -\sum_\tau\log(d_\tau!)
       -h\log(M/\Omega_E)\right|
       &\le 1+4f\log\zeta(2).
 \label{eq:identitynorm863}
\end{align}
Moreover \(J_n(\psi,1)H_{\nu,z}T_d\overline{T_d}\) in place
of \(H_{\nu,z}T_d\overline{T_d}\) still satisfies
\eqref{eq:allprimeperiod862} for every prescribed twist.
\end{corollary}
\begin{proof}
In the quotient of the two identities in Theorem~\ref{thm:torusnorm852}
the scalar \(a_{E,n}\), the old-level mass, \(\zeta_F(n)\) and
\(|a_n|^2\) cancel. The remaining factors are
\(|b_{\mathbf d}|^2 q_1^{h}\), the dimension,
\(\prod_\tau d_\tau!\), and the ratio of the two partial Hecke
Euler products at \(z\). Here \(q_1=\pi/A\) denotes the elliptic
factor at the identity embedding in this formula only.
For either unitary character the absolute value of the logarithm
of the positive partial Euler product is bounded by
\[
 \sum_{\mathfrak p}\sum_{j\ge1}
            \frac{(N\mathfrak p)^{-2j}}{j}
       =\log\zeta_E(2)\le 2f\log\zeta(2).
\]
The last inequality holds prime by prime because there are at
most \(2f\) primes above a rational prime and each has norm at
least that prime. Apply \eqref{eq:jetpair863} at the identity,
where \(\sigma'=\sigma\), to obtain the stated error. Finally
\(J_n(\psi,1)\) is a nonzero algebraic integer by
Theorem~\ref{thm:flattorus853}. Multiplying the common multiplier
by it decreases each negative valuation.

The class norm in \eqref{eq:identitynorm863} is an algebraic
number. At a general embedding its image is a pairing of two
conjugated vectors. Extending the norm comparison to that image
requires identifying both conjugated theta lifts in the same
class metric.
\end{proof}
\clearpage

\subsection{Galois conjugation of the symmetric Weyl translates}
\label{sec:weyltransport864}
Take \(\sigma\in\operatorname{Aut}(\CC/K)\) and the elliptic
curve \(C=\CC/(\delta\mathcal O_K)\) of Section~\ref{sec:arch59}.
The coordinate on its universal cover is the integral of the
marked algebraic differential. On its conjugate \(C^\sigma\)
we use the conjugated differential. We first put both period
lattices in the same rational alternating space.

\begin{lemma}[The conjugated torsion in a norm-one ideal]
\leavevmode\label{lem:torsionidele864}
There is a fractional ideal \(\mathfrak a\) of \(K\) of norm
one and a complex scalar \(\delta_\sigma\ne0\) for which the
conjugated lattice is \(\delta_\sigma\mathfrak a\). Under the
uniformizations by these lattices, the map on the rational
finite Tate modules induced by \(\sigma\) is multiplication by
some \(b_\sigma\in\mathbb A_{K,f}^{\times}\). One has
\[
 b_\sigma\widehat{\mathcal O}_K=\widehat{\mathfrak a},\qquad
 b_\sigma\bar b_\sigma=\chi_{\rm cyc}(\sigma),\qquad
 \frac{|\delta_\sigma|^2}{A_\sigma}
       =\frac{|\delta|^2}{A}=\frac{2}{\sqrt{d_{\rm K}}}.
\]
Here \(\chi_{\rm cyc}(\sigma)\in\widehat{\mathbb Z}^{\times}\)
describes the action on all roots of unity.
\end{lemma}
\begin{proof}
The maximal-order CM lattice of \(C^\sigma\) is homothetic
to a fractional ideal. Every ideal class of \(K\) is a square.
Indeed a class of order dividing two has a representative
\(\mathfrak l\) with \(\bar{\mathfrak l}/\mathfrak l=(y)\).
Then \(y\bar y=1\), since this positive rational generates the
unit ideal. Write \(y=t/\bar t\) by Hilbert's theorem 90.
The ideal \(t\mathfrak l\) is stable under conjugation. In its
prime factorization the split pairs give rational ideals and
the inert primes are rational ideals. The only ramified prime
is \((\sqrt{-d_{\rm K}})\), since \(-d_{\rm K}\) is the
discriminant and \(d_{\rm K}\) is prime in this family.
The ideal is therefore principal. Squaring on the finite
class group is now bijective. If \([\mathfrak t]^2\) is the
class of the conjugated lattice, take
\(\mathfrak a=\mathfrak t/\bar{\mathfrak t}\).
Its norm is one. The area of \(\mathcal O_K\) is
\(\sqrt{d_{\rm K}}/2\), giving the area identity.

The finite Tate modules, in the period coordinates just chosen,
are \(\widehat{\mathcal O}_K\) and \(\widehat{\mathfrak a}\).
The action on torsion commutes with the \(\mathcal O_K\)-action,
since a CM endomorphism has its prescribed tangent multiplier
in \(K\), fixed by \(\sigma\). It therefore acts by an idele
on the rank-one rational Tate modules and identifies the two
integral modules. Both principal polarizations are given in
these coordinates by the same alternating form, namely the
imaginary part of \(x\bar y\) times \(2/\sqrt{d_{\rm K}}\)
with a common orientation. On \(p^r\)-division points the
Weil pairing is the exponential of this integral pairing
divided by \(p^r\) on representatives in the Tate lattice.
It commutes with field embeddings, as follows also from its
definition by evaluating the functions of the two torsion
divisors. Substituting \(b_\sigma x,b_\sigma y\) multiplies
the alternating form by \(b_\sigma\bar b_\sigma\).
Perfectness modulo \(p^r\) therefore gives
\(b_{\sigma,p}\bar b_{\sigma,p}\equiv\chi_{\rm cyc,p}(\sigma)
\bmod p^r\) for every \(r\). This proves the idele identity.
For the Tate-module and polarization conventions see
\href{https://math.stanford.edu/~conrad/vigregroup/vigre04/mainthm.pdf}
{B.~Conrad, \emph{Main theorem of complex multiplication}, Sections 1 and 6}.
\end{proof}

For the even rank \(n\) use \(D=fn\) product coordinates in
Lemma~\ref{lem:basis59}. Write \(W_f=\mathbb A_{K,f}^D\).
Let \(R(x)\) be the symmetric finite Weyl translation of
Theorem~\ref{thm:vac242}. Its multiplication law has scalar
\(\psi_f(B(x,y)/2)\), where \(B\) is the rational alternating
form in that theorem and \(\psi_f\) is the finite additive
character. In the marked model, \(R(x)\) at a rational vector
\(x\in K^D\) acts by \(U_{\delta x}\) in the product variables.
This follows by canceling the common factor
\(e^{\pi i h_K(x,x)}\) in the translation identity in the proof
of Theorem~\ref{thm:vac242}. Write \(\Theta\) for the product of
the marked elliptic sections, normalized to have value one at
zero. The space of finite vectors is spanned by \(R(x)\Theta\).
Use \(R'(x)\) for the corresponding translation in the model
with \(\delta_\sigma,\mathfrak a,\Theta^\sigma\). Its rational
action is \(U_{\delta_\sigma x}\). The rational alternating
form and the multiplication law are the same by the area identity.
To obtain the full finite model from the conjugate product,
use one translate per coset modulo
\(\widehat{\mathfrak a}^{\,D}\). These translates are independent
by their distinct commutator characters on this self-dual lattice.
They span the smooth model by finite Fourier expansion, exactly
as in the proof of Theorem~\ref{thm:vac242}. The Weyl action on
each finite vector is locally constant.

\begin{theorem}[Exact conjugation of symmetric Weyl translations]
\leavevmode\label{thm:weyltransport864}
Coefficientwise conjugation of the Taylor series in the product
differentials gives a \(\sigma\)-semilinear bijection \(L_\sigma\)
of the two smooth finite models. It sends \(\Theta\) to
\(\Theta^\sigma\), and for every \(x\in W_f\) and finite vector \(G\),
\begin{equation}
              L_\sigma(R(x)G)=R'(b_\sigma x)L_\sigma G.
 \label{eq:weyltransport864}
\end{equation}
In the matrix coordinates \((Z_{\tau i})\) this statement uses
the permutation \(\tau\mapsto\sigma\tau\) of the oriented
embeddings. The same map applies to every homogeneous degree.
\end{theorem}
\begin{proof}
For a positive integer \(N\) sufficiently divisible for a finite
list of vectors, the proof of Theorem~\ref{thm:transport59}
writes each in the form
\[
                  G(z)=\Theta(z)S([z/N]),\qquad S\in\CC(C^D).
\]
Indeed for a rational torsion translate \(U_w\Theta\) the ratio
to \(\Theta\), pulled back by \(z\mapsto Nz\), is rational by
Theorem~\ref{thm:trans57}. We allow arbitrary complex coefficients
in the finite linear span. Substitution gives the Taylor series
of \(\Theta^\sigma(z)S^\sigma([z/N])\) after coefficientwise
conjugation. The local differential parameters are functorial by
the formal integration calculation in Theorem~\ref{thm:trans57}.
When the quotient has a pole at zero this identity is an identity
in the fraction field of the formal local ring, with the product
regular. By that theorem each individual translated product
becomes a root of unity times a translate in the conjugate
model. Thus the conjugated series represents a finite vector.
It determines that vector uniquely by holomorphic continuation.
This constructs \(L_\sigma\), proves its injectivity, and
identifies its value on \(\Theta\) by Theorem~\ref{thm:germ57}.

We show first that the two sides of \eqref{eq:weyltransport864}
are proportional by a scalar \(\epsilon(x)\) independent of
\(G\). Include \(\Theta\) in the list containing \(G\).
By density in the finite adeles choose a rational \(a\in K^D\)
arbitrarily close to \(x\). Choose a rational \(a'\) arbitrarily
close to \(b_\sigma a\). These choices can be made so that \(R(a)\)
and \(R(x)\) agree on the list, and \(R'(a')\) and \(R'(b_\sigma x)\)
agree on its images. They can also be made to satisfy
\[
              [\delta_\sigma a'/N]=\sigma([\delta a/N])
                             \quad\hbox{on }(C^\sigma)^D .
\]
In fact this last condition is exactly
\(a'-b_\sigma a\in N\widehat{\mathfrak a}^{\,D}\) by
Lemma~\ref{lem:torsionidele864}, an open condition. In \(R(a)G\)
the rational function \(S\) is evaluated at
\([z/N]+[\delta a/N]\). Conjugating it gives precisely the
argument in \(R'(a')L_\sigma G\). The remaining factor is the
conjugate of \(U_{\delta a}\Theta\), proportional to
\(U_{\delta_\sigma a'}\Theta^\sigma\) by
Theorem~\ref{thm:trans57} in each coordinate. Its proportionality
factor is determined by evaluating this comparison as an
identity of sections for \(G=\Theta\). That section is nonzero.
Thus for the given \(x\) one scalar works for every \(G\).

By the idele identity of Lemma~\ref{lem:torsionidele864},
\[
 \sigma\!\left(\psi_f(B(x,y)/2)\right)
                   =\psi_f(B(b_\sigma x,b_\sigma y)/2).
\]
Applying the two Weyl multiplication laws to \(x,y\) therefore
gives \(\epsilon(x+y)=\epsilon(x)\epsilon(y)\).
The reflection \(I:G(z)\mapsto G(-z)\) commutes with
coefficientwise conjugation. It acts on both finite models:
the marked elliptic section is even since its symmetric line
has a unique normalized section with value one at zero, and
reflection interchanges its translates at opposite points.
The formulas for symmetric translation give
\(IR(x)I=R(-x)\) and the corresponding identity for \(R'\).
It follows that \(\epsilon(-x)=\epsilon(x)\). Each value has
square one. Since every \(x\) in the additive finite adelic
space is twice \(x/2\), \(\epsilon(x)=\epsilon(x/2)^2=1\).
This proves the exact identity, including above 2.
Its range contains all translates of \(\Theta^\sigma\), since
\(b_\sigma\) is invertible. These span the target.

Finally the matrix taking product to matrix coordinates has
columns \((\tau(v_i))_\tau\) for vectors of the \(K\)-orthonormal
basis in Lemma~\ref{lem:basis59}. Conjugating its entries
permutes the rows as stated. The inverse coordinate change
also transports because \(E/\QQ\) is abelian. Taking homogeneous
parts of the same series proves the last assertion.
\end{proof}
\clearpage

\subsection{The untwisted theta norm at all complex embeddings}
\label{sec:allnorm865}
Keep the two lifts and the field of Corollary~\ref{cor:identitynorm863}.
They have the same even rank \(n=2z\) and the same point column
at every old mark. Write \(F_\mu\) for the class sum with
\(S_{\mu,z}=1\) in Theorem~\ref{thm:torusnorm852}, with no
new places. Thus \(F_\mu\) uses the unit real vacuum and the
one-row coefficient representation for \(\mu\).

For an automorphism \(\sigma\) fixing \(K\) define \(\mu^\sigma\)
at finite ideles of the norm-one torus by applying \(\sigma\)
to their character values. This is again a unitary automorphic
character, with the weights permuted by \(\tau\mapsto\sigma\tau\).
Indeed, modulo a compact open kernel and rational torus points
the finite class group is finite. A positive power of any finite
idele is a rational norm-one element times a member of the kernel.
Its character value to this power is a product of integral
powers of CM embeddings of that rational element. All conjugates
have modulus one. Taking roots proves algebraicity and unit
modulus for every conjugate of every finite value. The finite
kernel is preserved. The inverse archimedean character on a
rational point is permuted in precisely the stated way, which
proves automorphy with these infinite weights. It also gives
\(\mu^{c\sigma c}=\mu^\sigma\): at finite places
\(\overline{\sigma(\overline{\mu(g)})}
   =1/\overline{\sigma(\mu(g))}=\sigma(\mu(g))\),
and the two permutations of the CM embeddings agree since
\(E/\QQ\) is abelian.

\begin{lemma}[Transport of the marked finite test]
\leavevmode\label{lem:markedconj865}
Identify the source and target models in
Theorem~\ref{thm:weyltransport864} with the same abstract finite
oscillator representation. For the common old point test \(\Phi\)
there are a scalar \(\ell_\sigma\ne0\) and a finite norm-one
idele \(u\) such that \(L_\sigma\Phi=\ell_\sigma D_u\Phi\).
Here \(D_u\) is normalized change of position by \(u\) at
split places and is the identity elsewhere. For split-supported
source and target representatives \(g,y\), away from the marks,
let \(A\) act on positions by \(X\mapsto yXg^{-1}\), and let
\(D_A\) denote normalized change of position without a splitting
character. Then
\begin{equation}
 L_\sigma(D_A\Phi)
     =\ell_\sigma
       \frac{\sigma(t_y)}{t_y}\,D_uD_A\Phi,\qquad
 t_y=\prod_{v\ {\rm moving}}|\det y_v|_v^{-1/2}.
 \label{eq:splittransport865}
\end{equation}
The constants \(u,\ell_\sigma\) are independent of the character,
the real degree and the representatives.
\end{lemma}
\begin{proof}
The identification with the finite representation can be made by
the theta realization in the proof of Theorem~\ref{thm:vac242}.
For the conjugate product use the product-period vector for
\(\widehat{\mathfrak a}^{\,D}\) with its symmetric period
character. This is a line in the same finite oscillator space:
in a lattice model a prescribed period character of a self-dual
lattice occurs once, by expanding in the translates on its
cosets, or by the normalized lattice intertwiner and its
translates. The product section and this vector are identified
up to a common nonzero scalar. The action on rational translations
is \(U_{\delta_\sigma x}\) by rational invariance of the theta
sum, and agrees on finite translations by density. This proves
the asserted identification in the same symmetric Weyl convention.

Regard \(b_\sigma\) of Lemma~\ref{lem:torsionidele864} as an
idele of \(E\). Write \(b_\sigma=u a\), with \(u\bar u=1\)
and \(a\in\widehat{\mathcal O}_E^\times\), choosing \(a=1\)
on the position member of every old split pair. Indeed, at
each nonsplit place \(b_\sigma\) is a unit because its norm
is a unit. Take \(u=1\) there. At an old split pair take for
\(u\) the position component of \(b_\sigma\) and its conjugate
inverse on the other member. Make the same choice at the
finitely many other split pairs at which \(b_\sigma\) has
nonzero valuation, and set \(u=1\) elsewhere.
Then \(a\bar a=\chi_{\rm cyc}(\sigma)\).
Put \(L^0=D_u^{-1}L_\sigma\). Conjugation of a symmetric
Weyl operator by \(D_u\) multiplies its argument by \(u\):
on positions this is change of variables and on frequencies
it is its inverse transpose. Thus \(L^0R(x)=R(ax)L^0\).

Let \(\phi_0\) be the standard unmarked finite vacuum. Its
period lattice \(\widehat M\) is stable under \(a\). In the
symmetric convention its period character, on this lattice,
is the sign given at 2 by the parity of \(h_K(m,m)\), as in
Lemmas~\ref{lem:basis59} and~\ref{lem:split61}. This formula
on the completion follows by continuity from its rational
periods. The local trace norm of \(am\) at 2 is
\(\chi_{\rm cyc,2}(\sigma) h_K(m,m)\), with an odd unit as
multiplier. Thus the sign is fixed. The period line occurs
once and \(L^0\) is injective. Consequently
\(L^0\phi_0=\ell_\sigma\phi_0\) with \(\ell_\sigma\ne0\).

Here is the calculation on a finite set of split positions.
Trivialize the local position model so that the original box
vector is the indicator of a lattice \(U\), retaining its
original scalar normalization in every subsequent vector.
For \(L=\varpi^jU\subseteq U\), the indicator of \(L\) is
the box multiplied by
\[
        [U:L]^{-1}\sum_{\eta\in L^\perp/U^\perp}
                                      \psi_f(\langle X,\eta\rangle).
\]
This is the finite-character orthogonality formula on \(U/L\).
Under \(L^0\) the frequency scalar \(a\) on the dual member
permutes this quotient. On the position member a translate
by \(t\) goes to a translate by \(a t\). Any lattice \(A_vU\)
is a union of cosets of some such \(L\), and multiplication
by the scalar unit \(a\) permutes these cosets. At an old mark
\(a=1\) on positions, so the specified individual coset
modulo \(\varpi U\) is fixed. Apply these expansions to all
the finitely many factors of the one global box. They show
that the corresponding unnormalized product of indicators
is fixed up to the common \(\ell_\sigma\).
The normalized change \(D_A\) adds
\(\prod_v|\det A_v|_v^{-1/2}
 =t_y\prod_v|g_v|_v^{n/2}\).
The second factor is rational since \(n\) is even. Conjugating
the first gives exactly the quotient in
\eqref{eq:splittransport865}. The identity before moving the
representatives gives the formula for \(\Phi\).
\end{proof}

\begin{theorem}[The conjugates of the high-weight to flat norm ratio]
\leavevmode\label{thm:allnorm865}
Let \(N_\nu=\|\phi_{\nu,1}\|_{\rm cl}^2\) and
\(N_\eta=\|\phi_{\eta,1}\|_{\rm cl}^2\) with the common
rank and the scalings of Corollary~\ref{cor:identitynorm863}.
For every embedding of these algebraic numbers in \(\CC\),
\begin{equation}
 \left|\log|\sigma(N_\nu/N_\eta)|
   -\log\dim V_{\mathbf d}
   -\sum_\tau\log(d_\tau!)
   -h\log(M/\Omega_E)\right|
                   \le 1+4f\log\zeta(2).
 \label{eq:allnorm865}
\end{equation}
\end{theorem}
\begin{proof}
Extend the embedding to an automorphism of \(\CC\).
Composing on the left with complex conjugation preserves
absolute values. We may thus suppose \(\sigma\) fixes \(K\).
Write \(\sigma'=c\sigma c\). Choose the same split-supported
representatives for source and target when computing the two
norms and their conjugates, by Theorem~\ref{thm:class61}.
Avoid the old marks, all ramified places, and the support of
the ideles \(u\) for both automorphisms. The original class
weights are rational reciprocals of stabilizer orders.

In these representatives the finite part of the untwisted
kernel is the test \(D_A\Phi\) multiplied by
\(\mu_f(g)^{-1}\delta_n(y)\), where \(\delta_n(y)\) is the
intrinsic unitary target splitting factor. To verify the source
factor, at a split pair the restrictions of the splitting
character to the two members are inverse since its restriction
to the diagonal is trivial. Its \(n=2z\) position factors on a
source scalar give \(\upsilon_1^z\). The external factor is
\((\mu\upsilon_1^z)^{-1}\), giving the asserted cancellation
as in Theorem~\ref{thm:torusnorm852}. The same identity on
split units follows from the same character restriction.
There is no external target factor. Put
\(r_\sigma(y)=\sigma(t_y\delta_n(y))/(t_y\delta_n(y))\).
The product \(t_y^2\delta_n(y)\overline{\delta_n(y)}=t_y^2\)
is rational, so
\begin{equation}
                 r_\sigma(y)\overline{r_{\sigma'}(y)}=1 .
 \label{eq:targetcancel865}
\end{equation}

After the transport in Lemma~\ref{lem:markedconj865}, the extra
operator \(D_u\) is absorbed by sending \(g\) to \(g u^{-1}\):
its position matrix multiplies \(yXg^{-1}\) by \(u\). The
change of variables is normalized in both expressions and
the source character contributes one fixed unit-modulus
factor, independent of \(g,y\). This identity is valid also
on the split support of \(u\), where the other representatives
are the identity. Alternatively it is the normalized central
translation of Lemma~\ref{lem:centraltest736}. Right translation
permutes the torus classes and preserves their counting measure
since the source is abelian. We have therefore obtained the
actual class sum for \(\mu^\sigma\) at \(y\).

We spell out the real and scalar factors. The original
\(a_n q_{\rm id}^{h/2}\) converts each unscaled coefficient
to the homogeneous derivative of the marked section by
Theorem~\ref{thm:vac242}, where \(q_{\rm id}=\pi/A\).
Theorem~\ref{thm:weyltransport864} takes this derivative to
the corresponding derivative in the conjugate differential,
with the oriented labels permuted. In that differential the
same theta-functional calculation converts the derivative
back to the unit oscillator coefficient times
\(C_{\sigma,n}^{-1}q_\sigma^{h/2}\). Indeed rational Weyl
translations intertwine the two realizations and the product
period line is one dimensional, fixing \(C_{\sigma,n}\ne0\)
already in degree zero. Differentiation of
\(\exp(q_\sigma\bar w\cdot Z)\) gives exactly the indicated
power for \(P(\sqrt{q_\sigma}\bar w)\). In the rational matrix
coordinates \(w=\delta_\sigma\Phi(X)\), the Gaussian scale
\(q_\sigma|\delta_\sigma|^2\) is constant by
Lemma~\ref{lem:torsionidele864}. Thus the unit polynomial
changes from the original realization only by a degree-dependent
scalar of modulus one and the permutation of the oriented
indices. All components of the vector are covered by this
calculation: a symmetric tensor is obtained by polarization
of a pure power, and the symmetrizing projector averages
permutations with rational coefficients.

It follows that the vector of conjugated high-weight class
values at \(y\) equals the permuted vector of \(F_{\nu^\sigma}(y)\)
times
\[
  (\sigma b_{\mathbf d})q_\sigma^{h/2}
     C_{\sigma,n}^{-1}\ell_\sigma\,
                \epsilon_{\sigma,\nu}r_\sigma(y),
       \qquad |\epsilon_{\sigma,\nu}|=1.
\]
For the flat vector the identical formula holds with \(b=1\)
and \(h=0\), with the same \(C_{\sigma,n}^{-1}\ell_\sigma\).
Apply this description also to \(\sigma'\). In a monomial
basis the symmetric tensor metric has rational Gram matrix:
it is induced on symmetrized tensors by the orthonormal tensor
basis, with the pure highest tensor of norm one. Tensoring
over \(\tau\) gives the same metric after permuting the places.
In the class pairing \(\sigma(\bar v)=\overline{\sigma'(v)}\).
Both vectors then have the same character \(\mu^\sigma\) and
the same labels. Equation~\eqref{eq:targetcancel865} cancels
the target-dependent factors at each \(y\). The two products
of the constants \(C^{-1}\ell\) cancel upon division by the
flat norm, and the two remaining phases have modulus one.
Consequently
\[
 |\sigma(N_\nu/N_\eta)|=
    |\sigma b_{\mathbf d}|\,|\sigma' b_{\mathbf d}|\,
       (q_\sigma q_{\sigma'})^{h/2}\,
       \frac{\|F_{\nu^\sigma}\|_{\rm cl}^2}
            {\|F_{\eta^\sigma}\|_{\rm cl}^2}.
\]
The norms on the right are positive by
Theorem~\ref{thm:torusnorm852}. In their quotient the rank and
level factors cancel. The dimension and the factorials are
unchanged by permutation, and the logarithm of the ratio of
the partial Euler products is bounded in absolute value by
\(4f\log\zeta(2)\), as in Corollary~\ref{cor:identitynorm863}.
Insert \eqref{eq:jetpair863} to obtain \eqref{eq:allnorm865}.
\end{proof}
\clearpage

\subsection{A signed bound for the terminal determinant ratio}
\label{sec:terminaldirection866}
Fix a node \(x\) and the simultaneous family of twists of
Corollaries~\ref{cor:actualdenom803} and~\ref{cor:detheight810}.
For nonzero algebraic \(U\) independent of the coefficient
embedding, and a family \(Y=(Y_{p,\iota})\) of nonzero algebraic
numbers, write
\[
 \mathcal H_x(Y;U)=\sup_\sigma\log|\sigma U|
       +\sum_p\log p\,\sup_{\iota\mid p}
                       \max(0,-v_\iota(UY_{p,\iota})).
\]
A sum of nonnegative terms denotes the supremum over finite
subsums. The first term uses the ordinary logarithm.
For two such data,
\(\mathcal H_x(YY';UU')\le\mathcal H_x(Y;U)+\mathcal H_x(Y';U')\),
by the triangle inequality for positive parts and the product
law for absolute values. The same observation applies to any
nonnegative integral power. A family of powers of the new
conductor has \(\mathcal H_x(\,\cdot\,;1)=0\), since the
conductor of the twist attached to \(p,\iota\) is prime to \(p\).

For the two positive integral weights \(\kappa_\tau\) chosen in
Section~\ref{sec:shortpairs859} put
\begin{align}
 J_x&=\sum_\tau\big(\log(\kappa_\tau!)
                         +\kappa_\tau\log(M/\Omega_E)\big),\notag\\
 \mathcal A_x&=J_x-2\sum_p S_p\log p .
 \label{eq:direction866}
\end{align}
This last sum is finite. More explicitly its subtracted term is
\begin{equation}
 2\sum_p S_p\log p
 =\sum_\tau\sum_{\substack{1\le r\le xb_x\\(r,q_1)=1}}
       \Lambda(r)(1-\xi(r))
                \big(\kappa_\tau/r-\kappa_{[-r]\tau}\big)_+ .
 \label{eq:vonmangoldt866}
\end{equation}
Here \(\Lambda\) is the von Mangoldt function, \(\xi\) is the
quadratic character of \(K\), and \([-r]\) multiplies the
embedding label modulo \(q_1\). In fact, for \(r=p^a\) the
coefficient \(1-\xi(r)\) is two for an odd inert power, one
for every ramified power and zero otherwise. These are exactly
\(2/e_0\) on the allowed exponents of Section~\ref{sec:shortpairs859}.
All primes dividing \(q_1\) split in \(K\). No term with
\(p^a>b_x\) contributes because each \(\kappa_\tau\) is
at most \(b_x\) and at least one.

\begin{proposition}[A common multiplier for a main critical ratio]
\leavevmode\label{prop:mainheight866}
For each main critical pair, with its prescribed twist power
\(\beta=\beta_{p,\iota}\), there is a nonzero algebraic
\(U_{\nu,z}\), independent of \(p,\iota,\beta\), such that
\[
 \mathcal H_x\big((R_{\nu,z}(\beta)^{-1});U_{\nu,z}\big)
                    \le -\mathcal A_x+C_{q_0,d_{\rm K}}B_x^{\rm det}.
\]
This holds at every node of every sufficiently distant shorter
interval and for all main pairs there.
\end{proposition}
\begin{proof}
Use the high-weight normalization of Theorem~\ref{thm:jetscale863}
and any one flat comparison pair of rank \(n=2z\) in
Theorem~\ref{thm:allnorm865}. Denote its tame character by
\(\psi\). With the constants of Section~\ref{sec:twelftransfer855}
take
\[
           U_{\nu,z}= J_n(\psi,1)\,H_{\nu,z}\,T_d\overline{T_d}.
\]
All factors are nonzero and independent of the twist.
Since \(J_n(\psi,1)=B_n b_n^2 N_\eta\) and
\(H_{\nu,z}=c_*^2/(|C|^4 N_\nu)\), Theorem~\ref{thm:allnorm865}
gives
\[
 \sup_\sigma\log|\sigma U_{\nu,z}|
 \le -J_x+
 C_{q_0,d_{\rm K}}\big[
       x^{70}(1+\log b_x)+ f(1+b_x/k)\log(2b_x)\big].
\]
Here are the bounds on the factors. The representation dimension
is at least one, \(d_\tau\ge\kappa_\tau\) by
Theorem~\ref{thm:allprimetable862}, and \(M/\Omega_E\ge1\)
eventually since \(K\) is fixed, \(\operatorname{rd}(E)\ge1\)
and \(M\ge b_x^{20}\). These facts allow the factorials and
linear terms to be replaced in the required direction.
The finite integer \(|C|\ge1\) only improves the upper bound.
The bound for \(T_d\) and its conjugate is in
Corollary~\ref{cor:allprimeperiod862}. The factor \(c_*\)
satisfies the bound with \(x^{70}\) by
Corollary~\ref{cor:flatmean853} and Theorem~\ref{thm:flattorus853}.
Finally \(\log B_n\le(2fn)^2\log3\) and
\(\log b_n=fn\log(Mt_E)\), by definition. Here \(fn\ll x^{17}\),
\(\log t_E\ll\log(2x)\) and \(\log M\le22\log b_x\), as in the
proof of the same theorem. All these costs, including the
Euler product error, fit in the displayed bound. The bracket
is \(o(B_x^{\rm det})\) uniformly: divide by
\(fxb_x(\log b_x)^3/(mk)\) and use \(b_x\ge x^{400}\),
\(m\ll\sqrt x,\ k\ll x^5\).

By Corollaries~\ref{cor:allprimeperiod862} and~\ref{cor:identitynorm863},
\[
 -v_\iota\!\big(U_{\nu,z}\rho_\beta^2 R_{\nu,z}(\beta)^{-1}\big)
                         \le p_{p,\iota}+2S_p .
\]
The sum of the logarithmic cohomology bounds
\(\sum_p(\log p)\sup_\iota p_{p,\iota}\) is at most
\(B_x^{\rm det}\) by Corollary~\ref{cor:actualdenom803},
using its common bound on all integral torus tests.
For the four-row factor defining \(\rho_\beta\), take
\(\psi_2\) with the original prescribed tame character
of order \(Q\) as allowed in Section~\ref{sec:twelftransfer855}.
Its critical integer is 2 and its absolute angular weight
is 2. Its character is ramified at both old marks, and the
new twist powers are ramified at the new marks. Thus the
incomplete Euler ratio in Theorem~\ref{thm:flattorus853}
is the \(R_J\) of Theorem~\ref{thm:criticalheights734} for
this very twist, with inverse character conventions giving
the same positive absolute value. That theorem gives
\(\mathcal H_x((\rho_\beta^{-1});1)\le4\log2+1\)
using total tolerance at most one in the simultaneous
choice. Multiplying by \(\rho_\beta^{-2}\), summing the
bounds, and using \eqref{eq:direction866} proves the claim.
\end{proof}

\begin{theorem}[The signed estimate for the terminal class determinant]
\leavevmode\label{thm:terminaldirection866}
Let \(\mathfrak Z_\vartheta\) be the actual terminal class Gram
determinant of Corollary~\ref{cor:terminalcomparison851}, and
let \(D_x\) be its source dimension. On the same shorter
interval, with the same fields, weights and simultaneous
twists, there is a common nonzero algebraic \(U_x\) such that
\begin{equation}
 \mathcal H_x\!\left(
    (\mathfrak Z_1/\mathfrak Z_{\vartheta_{p,\iota}});U_x\right)
       \le -D_x m k\,\mathcal A_x
                         +C_{q_0,d_{\rm K}}D_x m k B_x^{\rm det}.
 \label{eq:terminaldirection866}
\end{equation}
In dimension zero both determinants and \(U_x\) equal one.
\end{theorem}
\begin{proof}
Assume \(D_x\ge1\) and fix a lift index \(j\). In
\(\mathcal L_j(\beta)/\mathcal L_j(1)\) from
Section~\ref{sec:compressionbudget851} there are three products
of absolute finite Hecke values, each squared, of lengths
\(m/2,m,(k+7m)/2\). The remaining factor is a power of the new
conductor, a unit at the coefficient prime. On the third product
take \(7m/2<z\le(k+7m)/2\). The \(k/2\) resulting ratios are
the main ratios of the preceding proposition. Indeed the
third character is \(\alpha^{-1}\upsilon_1^{3m+(k+7m)/2}\);
the factor \(\alpha\) changes by \(\beta\) and its tame part at
the identity twist is \(\chi_M^j\). Here one fixes the original
character extension first and multiplies by the prescribed
twist, as in Section~\ref{sec:refinedclass824}. Passing to the
inverse character gives the negative-weight convention of
Section~\ref{sec:torusnorm852}. Its angular weights are the
ones computed in Section~\ref{sec:shortpairs859}.

There remain \(m/2+m+7m/2=5m\) ratios. In the first two
products the characters are \(c\upsilon_1^{m/2}\) and
\(c^{-1}\upsilon_1^m\). Their absolute angular weights
are \(m/2,m\), since \(\upsilon_1\) has circle weight one.
In the third the angular weight at every place is at least
\((k+7m)/2\) by the main-pair window. Hence every remaining
integer is positive and at most its absolute angular weight.
The old tame parts in these three products are
\(\chi_M^{\pm j}\), nontrivial at each old pair since the
prescribed order \(Q>m\), and the twist powers are among
those in Theorem~\ref{thm:criticalheights734}. Old ramified
places contribute Euler factor one both before and after
the twist. At the new pairs the numerator is ramified
and the denominator has its original local factors.
Thus each is exactly one of that theorem's critical ratios.
Applying the theorem with the simultaneous list of both signs
costs at most \(4\log2+1\) for each inverse ratio with
multiplier one. Squaring the products now gives
\[
 \mathcal H_x\big((\mathcal L_j(1)/\mathcal L_j(\beta));U_j\big)
       \le -k\mathcal A_x+C_{q_0,d_{\rm K}}k B_x^{\rm det}
                                     +10m(4\log2+1),
 \qquad U_j=\prod_{7m/2<z\le(k+7m)/2} U_{\nu,z}^{\,2}.
\]

In Corollary~\ref{cor:terminalcomparison851} use tolerance
at most one. Its family \(y\) is the terminal ratio times
\(\prod_j(\mathcal L_j(\vartheta^j)/\mathcal L_j(1))^{D_x}\).
Its multiplier is \(H_x=\prod_j\Delta_j(1)\). The two
estimates of that corollary imply
\(\mathcal H_x(y;H_x)\ll_{q_0,d_{\rm K}}D_x m k B_x^{\rm det}\):
for the archimedean term use
\(f m^2b_x\log b_x/(mk B_x^{\rm det})
=m^2/(x(\log b_x)^2)\), and for the finite term use
\(s=k+m-1\ll k\). Multiply \(y\) by the \(m\) inverse
families just bounded, each to power \(D_x\). With
\(U_x=H_x\prod_j U_j^{D_x}\), subadditivity gives
\eqref{eq:terminaldirection866}, since \(m\ll k\) and
\(B_x^{\rm det}\to\infty\). For \(D_x=0\) the empty
determinants give the assertion directly.
\end{proof}

The estimate here is one side of a comparison. A bound on the
other side for the same terminal lattice is needed to turn
\eqref{eq:terminaldirection866} into the prime-sum inequality
of Theorem~\ref{thm:criterion16}.
\clearpage

\subsection{A product formula for full and partial terminal frames}
\label{sec:longframe867}
Use the terminal map from the common good source of dimension
\(d=D_x\), with the metric of Proposition~\ref{prop:terminalmetric807}.
Write \(A_\vartheta(Y)\) for the matrix with columns the terminal
values of the lifts of a tuple \(Y\), with rows indexed by a finite
class and a tensor coordinate of its compact representation.
The terminal rank is \(s+mt\) and the class set depends on the twist.

\begin{lemma}[Integral pairings on the actual terminal classes]
\leavevmode\label{lem:terminalintegral867}
Every stabilizer in the terminal class set is trivial. In tensor
coordinates the class Gram matrix of a tuple \(Y\) is
\(A_\vartheta(Y)^*A_\vartheta(Y)\). The local pairing of a Weyl
vector integral at \(\iota\) with a Weyl vector integral at
\(\iota c\), in the corresponding slots, is integral.
\end{lemma}
\begin{proof}
At an original mark \(v\mid p_0\) the terminal level imposes
identity reduction, as specified immediately before
Proposition~\ref{prop:terminalmetric807}. Here \(p_0\ge3\)
splits completely in \(E\). A conjugate of the local level
lies in a conjugate of \(1+p_0 M_{s+mt}(\mathbb Z_{p_0})\).
This group is torsion free. In fact a nonidentity element
has the form \(1+p_0^aB\), \(a\ge1\), with some unit entry
of \(B\). Taking a positive power prime to \(p_0\) preserves
the least entry valuation \(a\) of its difference from the
identity. Taking the \(p_0\)-th power increases it to \(a+1\)
by the binomial formula and oddness of \(p_0\). No positive
power can be the identity. The finite rational class stabilizer
injects into this local matrix group, so it is trivial.
In the class metric each class therefore has weight one.

For each column height \(i\), the coordinate wedges of the
standard orthonormal basis are an orthonormal basis of
\(\bigwedge^i\). Include the coefficient representation in
the tensor of these exterior powers, with the required
determinant twist if necessary. In the dual flag-minor
construction of Lemma~\ref{lem:dualcartan811}, taking products
of flag minors gives a surjection over \(\mathbb Z\) from the
tensor of the column-coordinate modules onto their product
span. Its dual embeds the Weyl lattice integrally in the
dual tensor. The extremal vector has a single unit coordinate,
since only the choice of the top minor in each column gives
that extremal monomial. Thus the restricted tensor metric
has the prescribed highest vector of norm one and equals the
class coefficient metric by irreducibility. Determinant lines
and their duals have the same unit normalization. Tensor the
construction over all real places. At a finite place the two
local frames in a Hermitian contraction are dual: a change
by a unitary matrix on the rational space acts by the inverse
transpose on the conjugate slot. Alternatively choose all
class representatives with identity component above the
coefficient prime by Theorem~\ref{thm:class61}. At that prime
the ordinary integral frames suffice. The coefficient pairing
in the ambient tensor has identity matrix, so summing over the
classes proves the two assertions.
\end{proof}

We specify the hypotheses on columns for the product formula.
At any coefficient prime use the same twist
\(\vartheta=\vartheta_{p,\iota}\) in the two slots labelled by
\(\iota,\iota c\). The local evaluations may be made on a
number field containing the coordinates when they are algebraic.
More generally one can extend an embedding to a field
isomorphism \(\CC\to\overline{\QQ}_p\) and evaluate there:
both algebraically closed fields have characteristic zero and
transcendence degree the continuum over \(\QQ\), so an extension
is obtained by matching transcendence bases. Indeed
\(\overline{\QQ}_p\) has the cardinality of the continuum and
the algebraic closure of a field generated over a countable
field by an infinite set has cardinality the maximum of that
cardinality and the countable cardinal. All hypotheses on
local evaluations below are required at every such embedding.
For the ordinary archimedean norm of a matrix or vector of
values, \(\sigma A\) means applying \(\sigma\) to the entries
in the original tensor coordinates, with the same fixed
representatives. For a rectangular complex matrix \(A\) with
\(n\) columns denote its squared exterior norm by
\(V(A)=\det(A^*A)\), including \(V(\varnothing)=1\).

\begin{theorem}[Covolume and index of a completed terminal frame]
\leavevmode\label{thm:longframe867}
Let \((Y,Z)=(y_1,\ldots,y_n,z_1,\ldots,z_r)\) be any basis
of the common good source, so \(n+r=d>0\). Suppose that
\[
 \sup_\sigma\log V(\sigma A_1(Y))\le c_0,\qquad
 \sup_{\sigma,i}\log\|\sigma A_1(z_i)\|^2\le b_0 .
\]
For each \(p,\iota\) suppose integral square matrices on
the \(n\) columns, separately at \(\iota\) and \(\iota c\),
take the twisted lifts of \(Y\) to integral vectors. Suppose
their determinants have valuation at most \(e_p\ge0\), with
\(\sum_p e_p\log p\le\delta\). On each of the \(r\) additional
columns let a common scalar \(\gamma_\vartheta\ne0\) make
its twisted lift integral at both evaluations. Suppose
\(\gamma_\vartheta\overline{\gamma_\vartheta}\) is algebraic
and that for some fixed nonzero algebraic \(V_0'\)
\[
 \mathcal H_x\big(((\gamma_\vartheta
       \overline{\gamma_\vartheta})^{-1});V_0'\big)\le b_1.
\]
For \(r=0\) omit all assumptions on these additional columns.
On the shorter intervals of Theorem~\ref{thm:terminaldirection866},
\begin{equation}
 c_0+2\delta+r(b_0+b_1)
      \ \ge\ d m k\,\mathcal A_x-C_{q_0,d_{\rm K}}d m k B_x^{\rm det}.
 \label{eq:longframe867}
\end{equation}
\end{theorem}
\begin{proof}
Set \(G_\vartheta=V(A_\vartheta(Y,Z))>0\). Positivity follows
from injectivity of the row map on the good source, proved in
Proposition~\ref{prop:chainpairing810}. A change of basis on
the full source multiplies both determinants by the same
squared modulus, so \(G_1/G_\vartheta=\mathfrak Z_1/\mathfrak
Z_\vartheta\). We first check algebraicity of these individual
determinants under the hypotheses. With \(\sigma'=c\sigma c\),
\(\sigma G_1\) is the mixed exterior pairing of the columns of
\(\sigma' A_1\) with those of \(\sigma A_1\). The
Cauchy--Schwarz inequality on the exterior power bounds it
by \(\sqrt{V(\sigma' A_1)V(\sigma A_1)}\). Each factor
\(V(\sigma A_1(Y,Z))\) is at most
\(V(\sigma A_1(Y))\prod_i\|\sigma A_1(z_i)\|^2\): project the
successive additional columns orthogonally off the preceding
ones in the Gram determinant formula. Hence
\[
                      \sup_\sigma\log|\sigma G_1|\le c_0+r b_0.
\]
A transcendental number has unbounded orbit under field
automorphisms of \(\CC\), since on a transcendence basis
containing it one may multiply it by any positive integer
and extend to the algebraic closure. Thus \(G_1\) is algebraic,
and so is \(G_\vartheta\) by the algebraic ratio of
Theorem~\ref{thm:terminaldirection866}.

At a fixed \(p,\iota\), use the two clearing matrices on
\(Y\). Multiply each remaining column by \(\gamma_\vartheta\)
in both evaluations. Lemma~\ref{lem:terminalintegral867}
makes every entry of the mixed Gram matrix integral.
Taking its determinant costs at most \(2e_p\) from the
two clearing matrices. More precisely,
\[
   -v_\iota\!\left((\gamma_\vartheta
       \overline{\gamma_\vartheta})^{r}G_\vartheta\right)
                           \le 2e_p .
\]
In the conjugate slot one uses the values of \(\iota c\),
so scalar multiplication contributes exactly
\(\iota(\overline{\gamma_\vartheta})\). This calculation
uses the same actual twist and class representatives in
the two slots, also when their clearing matrices differ.

The constant family \(G_1\) is the product of this last
family, the \(r\)-th power of
\((\gamma_\vartheta\overline{\gamma_\vartheta})^{-1}\),
and the terminal ratio. Subadditivity from
Section~\ref{sec:terminaldirection866} bounds its
\(\mathcal H_x\) with multiplier \(U=U_x(V_0')^r\) by
\[
 2\delta+r b_1-dmk\,\mathcal A_x
                               +C_{q_0,d_{\rm K}}dmk B_x^{\rm det}.
\]
For any nonzero algebraic \(\alpha\), the product formula
applied to \(U\alpha\) gives
\[
  0\le\sup_\sigma\log|\sigma\alpha|
                    +\mathcal H_x((\alpha);U).
\]
Indeed the average archimedean logarithm of \(U\alpha\)
is bounded by the two separate suprema and the negative
finite valuations in the formula by the displayed positive
parts, with the normalized local degree weights summing to
one above each rational prime. Apply this inequality to
\(G_1\) and use its archimedean bound to obtain
\eqref{eq:longframe867}.
\end{proof}

In particular any full frame whose twisted terminal columns
are integral in the two slots at all the prescribed places
has \(\sup_\sigma\log V(\sigma A_1(Y))\) at least the
right side of \eqref{eq:longframe867}. For a partial frame
the error is explicitly proportional to the number of
columns used to complete it. Applying this inequality in the
zero criterion still requires independent columns with
simultaneously small upper bounds for their untwisted norms
and their local indices, in positive dimension on the same nodes.
\clearpage

\clearpage
\section{References}\label{nav:bibliography-XVI}
\begingroup\def\section*#1{}\begin{thebibliography}{99}\small
\bibitem{XVI-XIII-IV-oh199} H. Oh, \emph{Uniform pointwise bounds for matrix coefficients of unitary representations and applications to Kazhdan constants}, Duke Math. J. \textbf{113} (2002), 133--192. \href{https://bpb-us-w2.wpmucdn.com/campuspress.yale.edu/dist/5/5177/files/2025/11/mat.pdf}{Theorem 1.1 and Corollary 1.5}.
\end{thebibliography}\endgroup
\end{document}
